% Lesson 57 — Vocabulary: Words and expressions related to using ICTs as recreational resources at home or in the community — Anglais Tle A — Cours signé OnBuch+
% Programme officiel MINESEC. Compiler avec : tectonic EN57-vocabulary-words-and-expressions-related-to-using-icts.tex
\newcommand{\DOCMATIERE}{Anglais}
\newcommand{\DOCNIVEAU}{Tle A}
\newcommand{\DOCMODULE}{Chapter 5 : Media and Communication}
\newcommand{\DOCLECON}{EN57}
\newcommand{\DOCTITRE}{Lesson 57 — Vocabulary: Words and expressions related to using ICTs as recreational resources at home or in the community}
% OnBuch+ — préambule « cours signés » ENRICHI (compilé avec tectonic / XeLaTeX)
% Système visuel « Pop » d'OnBuch+ : fond crème (jamais blanc), bordures encre
% épaisses, ombres « dures » décalées, titres Archivo Black en capitales.
% Version MATHÉMATIQUES : graphes (pgfplots/tikz), boîtes pédagogiques
% (définition, propriété, méthode, attention, le savais-tu, exercice résolu,
% exercice, TP, bilan), page de garde et signature OnBuch+ ancrée en bas.
%
% Macros attendues dans main.tex AVANT \input{preamble} :
%   \DOCMATIERE  \DOCNIVEAU  \DOCMODULE  \DOCLECON (numéro)  \DOCTITRE
\documentclass[11pt]{article}

\usepackage{fontspec}
\usepackage[english,french]{babel} % français = langue principale ; l'anglais via \eng{..} / textebox
\usepackage[a4paper,margin=2cm,top=3.4cm,bottom=2.8cm,headheight=1.4cm,footskip=1.3cm]{geometry}
\usepackage{amsmath,amssymb,amsfonts}
\usepackage{mathtools} % \xmapsto, \coloneqq... souvent utilisés par les modèles
\usepackage{cancel} % simplifications barrées (bilans de forces)
% \abs{z} (module, valeur absolue) et \norm{u} (norme), utilisés par les modèles
\providecommand{\abs}{}\let\abs\relax\DeclarePairedDelimiter{\abs}{\lvert}{\rvert}
\providecommand{\norm}{}\let\norm\relax\DeclarePairedDelimiter{\norm}{\lVert}{\rVert}
\usepackage[table]{xcolor} % [table] : fournit aussi \rowcolors (alternance de lignes)
\usepackage{pagecolor}
\usepackage{tikz}
\usetikzlibrary{arrows.meta,positioning,calc,shapes.geometric,shapes.misc,shapes.arrows,decorations.pathmorphing,decorations.pathreplacing,decorations.markings,patterns,shadows,mindmap,trees,fit,backgrounds,matrix,intersections,angles,quotes}
\usepackage{pgfplots}
\usepackage[version=4]{mhchem} % équations nucléaires \ce{^{235}_{92}U}
\usepackage[european,siunitx]{circuitikz} % schémas électriques
\pgfplotsset{compat=1.18}
\usepgfplotslibrary{fillbetween,groupplots} % aires entre courbes (calcul intégral)
\usepackage{enumitem}
\usepackage{fancyhdr}
% [most] charge toutes les bibliothèques tcolorbox (listings, minted, raster,
% documentation...) et gonfle inutilement la mémoire TeX sur un document long
% (~300 boîtes) : on ne charge que ce qui est réellement utilisé.
\usepackage[skins,breakable]{tcolorbox}
\usepackage{graphicx}
\usepackage{colortbl}
\usepackage{booktabs}
\usepackage{tabularx}
\usepackage{diagbox} % case d'en-tête diagonale des tableaux à double entrée
% Colonnes extensibles centrée / gauche / droite (C, L, R), souvent utilisées par les modèles
\newcolumntype{C}{>{\centering\arraybackslash}X}
\newcolumntype{L}{>{\raggedright\arraybackslash}X}
\newcolumntype{R}{>{\raggedleft\arraybackslash}X}
\usepackage{array}
\usepackage{multicol}
\usepackage{multirow}
\usepackage{siunitx}
% parse-numbers=false : accepte aussi \SI{2,5 \times 10^{-3}}{...}
\sisetup{locale=FR,output-decimal-marker={,},group-minimum-digits=4,parse-numbers=false}
\DeclareSIUnit\ohm{\ensuremath{\symup{\Omega}}} % U+2126 absent de la police
\DeclareSIUnit\division{div} % réglages d'oscilloscope (ms/div, V/div)
\DeclareSIUnit\tr{tr} % tours (vitesse de rotation, ex. \SI{1500}{\tr\per\minute})
\DeclareSIUnit\molar{M} % concentration molaire (ex. \SI{3}{\molar}, \SI{1}{\micro\molar})
\DeclareSIUnit\an{an} % année (le modèle confond parfois avec \year, un compteur TeX) — ex. \SI{5}{\kilo\meter\per\an}
\DeclareSIUnit\FCFA{FCFA} % franc CFA (ex. \SI{500}{\FCFA\per\kilo\gram})
\DeclareSIUnit\degreeBrix{\textdegree Brix} % teneur en sucre d'un sirop/jus (réfractométrie)
\DeclareSIUnit\month{mois} % le modèle utilise parfois \month, un compteur TeX, comme unité de durée
\usepackage{qrcode}
% \degree : commande fréquemment utilisée par les modèles pour un angle ou une
% température en dehors de \SI{}{} (non fournie par les paquets chargés).
\providecommand{\degree}{^{\circ}}
% \qed (fin de démonstration), utilisé par les modèles sans amsthm
\providecommand{\qed}{\hfill$\square$}
% \barre : double barre des valeurs interdites dans un tableau de variations (tkz-tab)
\providecommand{\barre}{\ensuremath{\|}}
% environnement proof (amsthm non chargé) utilisé par les modèles
\ifdefined\proof\else\newenvironment{proof}[1][Démonstration]{\par\noindent\textit{#1.}\ }{\qed\par}\fi
\usepackage{lastpage}
\usepackage{hyperref}
\hypersetup{hidelinks}

% ============================================================
% Palette « Pop » OnBuch+ — mêmes hex que la classe P (plus_ui.dart)
% ============================================================
\definecolor{poporange}{HTML}{F59321}
\definecolor{popdark}{HTML}{E07A0C}
\definecolor{popcream}{HTML}{FFF6E8}
\definecolor{popink}{HTML}{1C1714}
\definecolor{poppurple}{HTML}{7A5AE0}
\definecolor{popgreen}{HTML}{2E9E6A}
\definecolor{poppink}{HTML}{E0457B}
\definecolor{popblue}{HTML}{2F7DE1}
\definecolor{popgold}{HTML}{C98A12}
\definecolor{popmuted}{HTML}{8A7F76}
\definecolor{popline}{HTML}{EADFD2}
\definecolor{popgray}{HTML}{8A7F76}
\colorlet{popgrey}{popgray}
% Alias tolérants (le modèle nomme parfois les couleurs différemment)
\colorlet{popwhite}{white}
\colorlet{popblack}{popink}
\colorlet{popred}{poppink}
\colorlet{popyellow}{popgold}
% Teintes claires (fonds de tableaux/encadrés) — le modèle nomme parfois
% les couleurs avec un suffixe L (ex. popblueL) sans les avoir définies.
\colorlet{poporangeL}{poporange!20}
\colorlet{popdarkL}{popdark!20}
\colorlet{poppurpleL}{poppurple!20}
\colorlet{popgreenL}{popgreen!20}
\colorlet{poppinkL}{poppink!20}
\colorlet{popblueL}{popblue!20}
\colorlet{popgoldL}{popgold!20}
\colorlet{popredL}{popred!20}
\colorlet{popgrayL}{popgray!20}
% Teintes claires pour fonds de tableaux / zones
\colorlet{popblueL}{popblue!10!white}
\colorlet{poporangeL}{poporange!14!white}
\colorlet{popgreenL}{popgreen!12!white}
\colorlet{poppurpleL}{poppurple!12!white}
\colorlet{poppinkL}{poppink!12!white}
\colorlet{popgoldL}{popgold!14!white}

\pagecolor{popcream}

% ============================================================
% Typographie
% ============================================================
\defaultfontfeatures{Ligatures=TeX}
\setmainfont{PlusJakartaSans-Medium.ttf}[Path=../../../fonts/,BoldFont=PlusJakartaSans-Bold.ttf]
% Maths en sans-serif (Fira Math) avec chiffres et lettres droites de Plus
% Jakarta, pour que les formules se fondent dans le texte.
\usepackage[math-style=ISO,bold-style=ISO]{unicode-math}
\setmathfont{FiraMath-Regular.otf}[Path=../../../fonts/]
\setmathfont{PlusJakartaSans-Medium.ttf}[Path=../../../fonts/,range={up/num,up/{Latin,latin},\mathup}]
\usepackage{icomma} % virgule décimale sans espace en mode math
% Caractères mathématiques Unicode absents des polices de texte (Plus Jakarta,
% Archivo Black) : rendus par leur équivalent mathématique, en texte comme en
% formule — sinon ils s'affichent en carré vide (ex. « Arithmétique dans ℤ »).
\usepackage{newunicodechar}
\newunicodechar{ℝ}{\ensuremath{\mathbb{R}}}
\newunicodechar{ℤ}{\ensuremath{\mathbb{Z}}}
\newunicodechar{ℂ}{\ensuremath{\mathbb{C}}}
\newunicodechar{ℕ}{\ensuremath{\mathbb{N}}}
\newunicodechar{ℚ}{\ensuremath{\mathbb{Q}}}
\newunicodechar{𝔻}{\ensuremath{\mathbb{D}}}
\newunicodechar{α}{\ensuremath{\alpha}}
\newunicodechar{β}{\ensuremath{\beta}}
\newunicodechar{γ}{\ensuremath{\gamma}}
\newunicodechar{δ}{\ensuremath{\delta}}
\newunicodechar{ε}{\ensuremath{\varepsilon}}
\newunicodechar{θ}{\ensuremath{\theta}}
\newunicodechar{λ}{\ensuremath{\lambda}}
\newunicodechar{μ}{\ensuremath{\mu}}
\newunicodechar{σ}{\ensuremath{\sigma}}
\newunicodechar{τ}{\ensuremath{\tau}}
\newunicodechar{φ}{\ensuremath{\varphi}}
\newunicodechar{ψ}{\ensuremath{\psi}}
\newunicodechar{ω}{\ensuremath{\omega}}
\newunicodechar{Σ}{\ensuremath{\Sigma}}
% majuscules produites par \MakeUppercase dans les titres (α-aminés -> Α)
\newunicodechar{Α}{\ensuremath{\alpha}}
\newunicodechar{Β}{\ensuremath{\beta}}
\newunicodechar{Γ}{\ensuremath{\Gamma}}
\newunicodechar{Δ}{\ensuremath{\Delta}}
\newunicodechar{Ε}{\ensuremath{\varepsilon}}
\newunicodechar{Θ}{\ensuremath{\Theta}}
\newunicodechar{Λ}{\ensuremath{\Lambda}}
\newunicodechar{Μ}{\ensuremath{\mu}}
\newunicodechar{Τ}{\ensuremath{\tau}}
\newunicodechar{Φ}{\ensuremath{\Phi}}
\newunicodechar{Ψ}{\ensuremath{\Psi}}
\newunicodechar{Ω}{\ensuremath{\Omega}}
\newunicodechar{↦}{\ensuremath{\mapsto}}
\newunicodechar{∈}{\ensuremath{\in}}
\newunicodechar{∉}{\ensuremath{\notin}}
\newunicodechar{∧}{\ensuremath{\wedge}}
\newunicodechar{⇔}{\ensuremath{\Leftrightarrow}}
\newunicodechar{⟺}{\ensuremath{\Longleftrightarrow}}
\newunicodechar{⇒}{\ensuremath{\Rightarrow}}
\newunicodechar{⟹}{\ensuremath{\Longrightarrow}}
\newunicodechar{∘}{\ensuremath{\circ}}
\newunicodechar{∪}{\ensuremath{\cup}}
\newunicodechar{∩}{\ensuremath{\cap}}
\newunicodechar{≡}{\ensuremath{\equiv}}
\newunicodechar{⊂}{\ensuremath{\subset}}
\newunicodechar{⊥}{\ensuremath{\perp}}
\newunicodechar{∃}{\ensuremath{\exists}}
\newunicodechar{∀}{\ensuremath{\forall}}
\newunicodechar{∛}{\ensuremath{\sqrt[3]{\,}}}
\newunicodechar{∜}{\ensuremath{\sqrt[4]{\,}}}
\newunicodechar{⃗}{\ensuremath{{}^{\to}}}
\newunicodechar{ⁿ}{\ifmmode^{n}\else\textsuperscript{\ensuremath{n}}\fi}
\newunicodechar{ˣ}{\ifmmode^{x}\else\textsuperscript{\ensuremath{x}}\fi}
\newunicodechar{ᵇ}{\ifmmode^{b}\else\textsuperscript{\ensuremath{b}}\fi}
\newunicodechar{ʳ}{\ifmmode^{r}\else\textsuperscript{\ensuremath{r}}\fi}
\newunicodechar{ᵘ}{\ifmmode^{u}\else\textsuperscript{\ensuremath{u}}\fi}
\newunicodechar{ʸ}{\ifmmode^{y}\else\textsuperscript{\ensuremath{y}}\fi}
\newunicodechar{ᵃ}{\ifmmode^{a}\else\textsuperscript{\ensuremath{a}}\fi}
\newunicodechar{ᵉ}{\ifmmode^{e}\else\textsuperscript{\ensuremath{e}}\fi}
\newunicodechar{ᵏ}{\ifmmode^{k}\else\textsuperscript{\ensuremath{k}}\fi}
\newunicodechar{ᵖ}{\ifmmode^{p}\else\textsuperscript{\ensuremath{p}}\fi}
\newunicodechar{ᴺ}{\ifmmode^{N}\else\textsuperscript{\ensuremath{N}}\fi}
\newunicodechar{ᶜ}{\ifmmode^{c}\else\textsuperscript{\ensuremath{c}}\fi}
\newunicodechar{ʰ}{\ifmmode^{h}\else\textsuperscript{\ensuremath{h}}\fi}
\newunicodechar{ᵠ}{\ifmmode^{q}\else\textsuperscript{\ensuremath{q}}\fi}
\newunicodechar{ₙ}{\ifmmode_{n}\else\textsubscript{\ensuremath{n}}\fi}
\newunicodechar{ₐ}{\ifmmode_{a}\else\textsubscript{\ensuremath{a}}\fi}
\newunicodechar{ᵢ}{\ifmmode_{i}\else\textsubscript{\ensuremath{i}}\fi}
\newunicodechar{ᵣ}{\ifmmode_{r}\else\textsubscript{\ensuremath{r}}\fi}
\newunicodechar{ₖ}{\ifmmode_{k}\else\textsubscript{\ensuremath{k}}\fi}
\newunicodechar{ᵦ}{\ifmmode_{\beta}\else\textsubscript{\ensuremath{\beta}}\fi}
\newunicodechar{ₓ}{\ifmmode_{x}\else\textsubscript{\ensuremath{x}}\fi}
\newfontfamily\popdisplay{ArchivoBlack-Regular.ttf}[Path=../../../fonts/]
\newfontfamily\poplogo{SpaceGrotesk-Bold.ttf}[Path=../../../fonts/]
\newfontfamily\popsemi{PlusJakartaSans-SemiBold.ttf}[Path=../../../fonts/]
\newfontfamily\popxbold{PlusJakartaSans-ExtraBold.ttf}[Path=../../../fonts/]
\newfontfamily\popxboldls{PlusJakartaSans-ExtraBold.ttf}[Path=../../../fonts/,LetterSpace=3.0]

\setlist[enumerate]{leftmargin=1.7em,itemsep=5pt,topsep=4pt}
\setlist[itemize]{leftmargin=1.7em,itemsep=4pt,topsep=4pt,label={\color{poporange}\textbullet}}
\setlength{\parindent}{0pt}
\setlength{\parskip}{5pt}
\renewcommand{\arraystretch}{1.25}

% pgfplots : style Pop par défaut
\pgfplotsset{
  every axis/.append style={axis line style={popink,line width=1pt},tick style={popink},
    grid style={popline,line width=0.5pt},label style={font=\small},tick label style={font=\footnotesize}},
  popcurve/.style={poporange,line width=1.8pt,no markers,smooth},
}
% Même style disponible dans un \draw TikZ simple (hors pgfplots)
\tikzset{popcurve/.style={poporange,line width=1.8pt}}
\tikzset{popgrid/.style={popline,very thin}}
\colorlet{popcurve}{poporange} % le nom du style est parfois utilisé comme couleur

% ============================================================
% Pill / squiggle
% ============================================================
\newcommand{\poppill}[3]{%
  \tikz[baseline=-0.55ex]\node[draw=popink,line width=1.2pt,rounded corners=2.6pt,%
    fill=#2,inner xsep=6.5pt,inner ysep=3pt]{%
    \popxboldls\fontsize{7}{7}\selectfont\color{#3}\MakeUppercase{#1}};%
}
\newcommand{\popsquiggle}[1]{%
  \tikz[baseline=-0.4ex]{
    \draw[#1,line width=2.6pt,line cap=round]
      plot [smooth] coordinates {(0,0.09) (0.34,0.01) (0.68,0.21) (1.02,0.01) (1.36,0.10)};
  }%
}
% Mot-clé surligné (marqueur orange sous le texte)
\newcommand{\cle}[1]{{\popxbold\color{popdark}#1}}

% ============================================================
% En-tête / pied
% ============================================================
\pagestyle{fancy}
\fancyhf{}
\renewcommand{\headrulewidth}{0pt}
\renewcommand{\footrulewidth}{0pt}
\lhead{\raisebox{-0.15ex}{\poplogo\fontsize{13}{13}\selectfont\color{popink}OnBuch\color{poporange}+}}
\rhead{\poppill{\DOCMATIERE\ \textbullet\ \DOCNIVEAU\ \textbullet\ Leçon \DOCLECON}{white}{popink}}
\lfoot{\popxbold\fontsize{7.6}{7.6}\selectfont\color{popmuted}\MakeUppercase{Cours signé OnBuch+}\ \textbullet\ {\popsemi\DOCTITRE}}
\rfoot{\tikz[baseline=-0.6ex]\node[rounded corners=8.5pt,fill=popink,minimum height=17pt,minimum width=17pt,inner xsep=5pt,inner sep=0]{\popxbold\fontsize{8}{8}\selectfont\color{popcream}\thepage\,/\,\pageref*{LastPage}};}

% ============================================================
% Titres
% ============================================================
\newcommand{\chaptitre}[2]{%
  \begin{tcolorbox}[enhanced,colback=poporange,colframe=popink,boxrule=2.6pt,arc=11pt,
      drop shadow=popink,left=15pt,right=15pt,top=12pt,bottom=11pt,before skip=3pt,after skip=18pt]
    \poppill{#2}{popink}{popcream}\par\vspace{9pt}
    {\raggedright\hyphenpenalty=10000\exhyphenpenalty=10000\popdisplay\fontsize{22}{24}\selectfont\color{popcream}\MakeUppercase{#1}\par}
  \end{tcolorbox}
}
% Section numérotée : pastille orange avec le numéro + titre Archivo
\newcounter{popsec}
\newcommand{\coursec}[1]{%
  \stepcounter{popsec}\setcounter{popsub}{0}%
  \par\vspace{16pt}\noindent
  \begin{minipage}{\linewidth}
  \tikz[baseline=-0.7ex]\node[draw=popink,line width=1.6pt,fill=poporange,rounded corners=4pt,minimum size=22pt,inner sep=0]{\popdisplay\fontsize{12}{12}\selectfont\color{popcream}\thepopsec};%
  \hspace{8pt}{\popdisplay\fontsize{15}{17}\selectfont\color{popink}\MakeUppercase{#1}}\par
  \vspace{4pt}\noindent{\color{poporange}\rule{36pt}{3.6pt}}
  \end{minipage}\par\nopagebreak\vspace{8pt}
}
\newcounter{popsub}
\newcommand{\courssub}[1]{%
  \stepcounter{popsub}%
  \par\vspace{9pt}\noindent{\popxbold\fontsize{11.5}{13}\selectfont\color{popink}{\color{poporange}\thepopsec.\thepopsub}\hspace{6pt}#1}\par\nopagebreak\vspace{4pt}
}

% ============================================================
% Boîtes pédagogiques — même recette : fond blanc, bordure épaisse colorée,
% ombre dure encre, badge plein. Titre optionnel [#1].
% ============================================================
\tcbset{popbase/.style={breakable,enhanced,colback=white,boxrule=2.3pt,arc=9pt,
  shadow={3pt}{-3pt}{0pt}{popink},left=14pt,right=14pt,top=11pt,bottom=11pt,before skip=11pt,after skip=15pt,
  fontupper=\popsemi\fontsize{10.6}{14.5}\selectfont\color{popink},
  fontlower=\popsemi\fontsize{10.6}{14.5}\selectfont\color{popink}}}
% Coche dessinée (le glyphe ✓ n'existe pas dans les polices du document)
\newcommand{\popcheck}[1]{\tikz[baseline=-0.5ex]\draw[#1,line width=1.6pt,line cap=round,line join=round](-0.13,0.0)--(-0.04,-0.09)--(0.13,0.1);}
\providecommand{\coche}{\popcheck{popgreen}}
\providecommand{\resultat}[1]{\par\smallskip\noindent\fcolorbox{popgreen}{popgreenL}{\parbox{\dimexpr\linewidth-2\fboxsep-2\fboxrule\relax}{\textbf{Résultat.} #1}}\par\smallskip}
\newcommand{\popboxhead}[3]{\poppill{#1}{#2}{white}\ifx\relax#3\relax\else\hspace{6pt}{\popxbold\fontsize{10.6}{12}\selectfont\color{popink}#3}\fi\par\vspace{7pt}}

\newtcolorbox{aretenir}[1][]{popbase,colframe=popgreen,before upper={\popboxhead{À retenir}{popgreen}{#1}}}
% Texte en anglais (document de lecture, dialogue, extrait) : cadre
% d'étude. Un titre optionnel (source, auteur) s'affiche dans l'en-tête.
\newtcolorbox{textebox}[1][]{popbase,colframe=popblue,before upper={\popboxhead{Text}{popblue}{#1}\begin{otherlanguage}{english}},after upper={\end{otherlanguage}}}
% Marques ✓ / ✗ (forme correcte / fautive) souvent utilisées par les modèles
\providecommand{\cross}{\textcolor{poppink}{\ensuremath{\times}}}
\providecommand{\xmark}{\cross}\providecommand{\crossmark}{\cross}\providecommand{\cmark}{\ensuremath{\checkmark}}
% Ligne de réponse / justification à compléter (inventée par les modèles)
\providecommand{\justification}[1]{\par\noindent\textit{Justification~:} \dotfill\par}
% \textipa (package tipa non chargé) : la police ne couvre pas l'API -> simple passage
\providecommand{\textipa}[1]{#1}
% Soulignement (ulem non chargé) : \uline, \ul
\providecommand{\uline}[1]{\underline{#1}}\providecommand{\ul}[1]{\underline{#1}}
% \glossaire{\item ...} inventée par les modèles : liste de glossaire
\providecommand{\glossaire}[1]{\begin{itemize}#1\end{itemize}}
% \alert (beamer) inventée par les modèles : texte mis en évidence
\providecommand{\alert}[1]{\textbf{\textcolor{poppink}{#1}}}
% variantes ASCII d'ordinaux écrites en macro
\providecommand{\iere}{\textsuperscript{re}}\providecommand{\ieme}{\textsuperscript{e}}\providecommand{\ere}{\textsuperscript{re}}\providecommand{\eme}{\textsuperscript{e}}\providecommand{\er}{\textsuperscript{er}}\providecommand{\ie}{\textsuperscript{e}}
% Ordinaux français écrits en macro par les modèles : 1\ière, 3\ième
\newcommand{\ière}{\textsuperscript{re}}\newcommand{\ième}{\textsuperscript{e}}
% \exemple (inventée par les modèles) : étiquette « Exemple : »
\providecommand{\exemple}{\textbf{Exemple~:}\ }
% \square utilisable aussi hors mode math (cases à cocher)
\AtBeginDocument{\let\squareMath\square\renewcommand{\square}{\ensuremath{\squareMath}}}
% Mot ou phrase en anglais à l'intérieur d'un paragraphe français
\newcommand{\eng}[1]{\ifmmode\text{\textit{#1}}\else\begingroup\selectlanguage{english}\itshape #1\endgroup\fi}
\let\esp\eng % alias toléré
% flèches d'intonation inventées par les modèles
\providecommand{\textsearrow}{}\renewcommand{\textsearrow}{\ensuremath{\searrow}}\providecommand{\textnearrow}{}\renewcommand{\textnearrow}{\ensuremath{\nearrow}}
% \fr{..} : mot français (inventé par les modèles)
\providecommand{\fr}[1]{\textit{#1}}
\newtcolorbox{exemplebox}[1][]{popbase,colframe=poporange,before upper={\popboxhead{Exemple}{poporange}{#1}}}
\newtcolorbox{definition}[1][]{popbase,colframe=popblue,before upper={\popboxhead{Définition}{popblue}{#1}}}
\newtcolorbox{propriete}[1][]{popbase,colframe=poppurple,before upper={\popboxhead{Propriété}{poppurple}{#1}}}
\newtcolorbox{methode}[1][]{popbase,colframe=popgold,before upper={\popboxhead{Méthode}{popgold}{#1}}}
\newtcolorbox{attention}[1][]{popbase,colframe=poppink,before upper={\popboxhead{Attention}{poppink}{#1}}}
\newtcolorbox{experience}[1][]{popbase,colframe=popblue,colback=popblueL,before upper={\popboxhead{Expérience}{popblue}{#1}}}
% « Le savais-tu ? » : carte violette pleine (contexte, culture, Cameroun)
\newtcolorbox{savaistu}[1][]{popbase,colback=poppurple,colframe=popink,
  fontupper=\popsemi\fontsize{10.4}{14.2}\selectfont\color{white},
  before upper={\poppill{Le savais-tu\,?}{white}{poppurple}\ifx\relax#1\relax\else\hspace{6pt}{\popxbold\color{white}#1}\fi\par\vspace{7pt}}}
% Exercice résolu : énoncé en haut, \tcblower puis la solution sur fond crème
\newcounter{popexo}\newcounter{popexr}
\newtcolorbox{exoresolu}[1][]{popbase,colframe=poporange,colbacklower=popcream,
  segmentation style={popink,line width=1.2pt,dashed},
  before upper={\stepcounter{popexr}\popboxhead{Exercice résolu \thepopexr}{poporange}{#1}},
  before lower={\poppill{Solution}{popink}{popcream}\par\vspace{6pt}}}
% Exercice (non résolu en place) — la correction est dans la partie Corrigés
\newtcolorbox{exercice}[1][]{popbase,colframe=popink,
  before upper={\stepcounter{popexo}\popboxhead{Exercice \thepopexo}{popink}{#1}}}
% Corrigé
\newtcolorbox{corrige}[1][]{popbase,colframe=popgreen,colback=popgreenL,
  before upper={\popboxhead{Corrigé}{popgreen}{#1}}}
% Objectifs / prérequis / bilan
\newtcolorbox{objectifs}{popbase,colframe=popink,colback=poporangeL,
  before upper={\popboxhead{Objectifs de la leçon}{popink}{}}}
\newtcolorbox{prerequis}{popbase,colframe=popink,colback=popblueL,
  before upper={\popboxhead{Prérequis}{popblue}{}}}
\newtcolorbox{bilan}{popbase,colframe=popink,colback=white,boxrule=2.8pt,
  before upper={\popboxhead{Fiche bilan}{poporange}{L'essentiel à connaître pour l'examen}}}
% Figure encadrée avec légende
\newtcolorbox{popfigure}[1][]{enhanced,breakable=false,colback=white,colframe=popline,boxrule=1.4pt,arc=8pt,
  left=10pt,right=10pt,top=10pt,bottom=8pt,before skip=10pt,after skip=12pt,halign=center,
  lower separated=false,halign lower=center,
  fontlower=\popsemi\fontsize{9}{11.5}\selectfont\color{popmuted},
  #1}
\newcommand{\legende}[1]{\par\vspace{6pt}{\popsemi\fontsize{9}{11.5}\selectfont\color{popmuted}#1}}

% ============================================================
% Signature OnBuch+
% ============================================================
\newcommand{\poplogoinline}[1]{{\poplogo\fontsize{#1}{#1}\selectfont\color{popink}OnBuch\color{poporange}+}}
\newcommand{\signatureonbuch}{%
  \begin{tcolorbox}[enhanced,colback=popink,colframe=popink,boxrule=2.2pt,arc=10pt,
      drop shadow=poporange,left=14pt,right=14pt,top=10pt,bottom=10pt,before skip=0pt,after skip=0pt]
    \tikz[baseline=-0.7ex]\node[circle,fill=popgreen,draw=popcream,line width=1.3pt,minimum size=22pt,inner sep=0]{\popcheck{white}};%
    \hspace{9pt}\begin{minipage}[c]{0.62\linewidth}
      {\popxbold\fontsize{12}{13}\selectfont\color{popcream}Signé\ \textbullet\ {\poplogo OnBuch\color{poporange}+}}\par\vspace{2pt}
      {\raggedright\popsemi\fontsize{8}{10.5}\selectfont\color{popline}\MakeUppercase{Équipe pédagogique OnBuch+}\\\MakeUppercase{Conforme au programme officiel MINESEC}\par}
    \end{minipage}\hfill
    {\poplogo\fontsize{22}{22}\selectfont\color{popcream}OnBuch\color{poporange}+}
  \end{tcolorbox}%
}

% ============================================================
% Page de garde
% #1 titre · #2 matière · #3 niveau · #4 module · #5 n° leçon · #6 durée ·
% #7 description / objectifs (texte ou itemize)
% ============================================================
\newcommand{\pagegarde}[7]{%
  \thispagestyle{empty}
  \newgeometry{a4paper,margin=1.7cm,top=1.6cm,bottom=1.4cm}%
  \begin{tikzpicture}[remember picture,overlay]
    \fill[poporange!18] ([xshift=-3.2cm,yshift=-3.6cm]current page.north east) circle (4.6cm);
    \fill[poppurple!14] ([xshift=2.4cm,yshift=7.6cm]current page.south west) circle (3.8cm);
  \end{tikzpicture}%
  {\poplogo\fontsize{30}{30}\selectfont\color{popink}OnBuch\color{poporange}+}\hfill
  \raisebox{0.6ex}{\poppill{Cours signé \textbullet\ Premium}{popink}{popcream}}\par
  \vspace{1pt}\popsquiggle{poppurple}\par
  \vspace{8pt}
  \begin{tcolorbox}[enhanced,colback=poporange,colframe=popink,boxrule=2.8pt,arc=13pt,
      drop shadow=popink,left=18pt,right=18pt,top=12pt,bottom=13pt,after skip=10pt]
    \poppill{#2 \textbullet\ #3}{popink}{popcream}\hspace{5pt}\poppill{#4}{white}{popink}\par\vspace{8pt}
    {\popdisplay\fontsize{14}{14}\selectfont\color{popink}LEÇON #5}\par\vspace{4pt}
    {\raggedright\hyphenpenalty=10000\exhyphenpenalty=10000\popdisplay\fontsize{25}{27}\selectfont\color{popcream}\MakeUppercase{#1}\par}
    \vspace{8pt}
    {\popxbold\fontsize{9.5}{9.5}\selectfont\color{popink}\MakeUppercase{Durée conseillée : #6}}
  \end{tcolorbox}
  \begin{tcolorbox}[enhanced,colback=white,colframe=popink,boxrule=2.2pt,arc=10pt,
      drop shadow=popink,left=16pt,right=16pt,top=10pt,bottom=10pt,after skip=8pt]
    \poppill{Dans cette leçon}{poppurple}{white}\par\vspace{5pt}
    \popsemi\fontsize{9.3}{12.6}\selectfont\color{popink}#7
  \end{tcolorbox}
  \noindent\begin{minipage}[t]{0.487\linewidth}
    \begin{tcolorbox}[enhanced,colback=popgreen,colframe=popink,boxrule=2.2pt,arc=10pt,
        drop shadow=popink,left=11pt,right=11pt,top=10pt,bottom=10pt,height=3.1cm,valign=center]
      \tcbox[colback=white,colframe=popink,boxrule=1.3pt,arc=5pt,boxsep=0pt,nobeforeafter,
        left=3pt,right=3pt,top=3pt,bottom=3pt,valign=center]{\qrcode[height=1.9cm]{https://play.google.com/store/apps/details?id=com.luvvix.onbuch&hl=fr}}%
      \hspace{8pt}\begin{minipage}[c]{0.5\linewidth}
        {\popdisplay\fontsize{11}{12.5}\selectfont\color{white}\MakeUppercase{Télécharge OnBuch}}\par\vspace{4pt}
        {\raggedright\popsemi\fontsize{8.4}{11}\selectfont\color{white}Gratuit \textbullet\ Google Play\\Tuteur IA, annales, concours, résultats\par}
      \end{minipage}
    \end{tcolorbox}
  \end{minipage}\hfill
  \begin{minipage}[t]{0.487\linewidth}
    \begin{tcolorbox}[enhanced,colback=poppurple,colframe=popink,boxrule=2.2pt,arc=10pt,
        drop shadow=popink,left=11pt,right=11pt,top=10pt,bottom=10pt,height=3.1cm,valign=center]
      \tikz[baseline=-0.7ex]\node[circle,fill=white,draw=popink,line width=1.3pt,minimum size=1.6cm,inner sep=0]{\popdisplay\fontsize{24}{24}\selectfont\color{poppurple}+};%
      \hspace{8pt}\begin{minipage}[c]{0.6\linewidth}
        {\popdisplay\fontsize{11}{12.5}\selectfont\color{white}\MakeUppercase{Passe à OnBuch+}}\par\vspace{4pt}
        {\raggedright\popsemi\fontsize{8.4}{11}\selectfont\color{white}Tous les cours signés, Tuteur IA illimité, fiches PDF — onglet OnBuch+ de l'app\par}
      \end{minipage}
    \end{tcolorbox}
  \end{minipage}
  \vspace{8pt}
  \popcontenu
  \vfill
  \signatureonbuch
  \restoregeometry
  \newpage
}

% Rangée « Ce cours contient » (page de garde)
\newcommand{\poptile}[3]{% #1 couleur #2 gros texte #3 légende
  \begin{tcolorbox}[enhanced,colback=#1,colframe=popink,boxrule=2pt,arc=9pt,shadow={2.5pt}{-2.5pt}{0pt}{popink},
      left=6pt,right=6pt,top=9pt,bottom=9pt,halign=center,valign=center,height=1.75cm,width=\linewidth,nobeforeafter]
    {\popdisplay\fontsize{12.5}{13}\selectfont\color{white}\MakeUppercase{#2}}\par\vspace{3pt}
    {\centering\popsemi\fontsize{7.8}{9.5}\selectfont\color{white}#3\par}
  \end{tcolorbox}}
\newcommand{\popcontenu}{%
  \par\noindent{\popxboldls\fontsize{7.5}{7.5}\selectfont\color{popmuted}CE COURS SIGNÉ CONTIENT}\par\vspace{7pt}
  \noindent\begin{minipage}[t]{0.235\linewidth}\poptile{popblue}{Cours}{complet, illustré, pas à pas}\end{minipage}\hfill
  \begin{minipage}[t]{0.235\linewidth}\poptile{poporange}{Méthodes}{et exercices résolus}\end{minipage}\hfill
  \begin{minipage}[t]{0.235\linewidth}\poptile{poppink}{Exercices}{type examen + corrigés}\end{minipage}\hfill
  \begin{minipage}[t]{0.235\linewidth}\poptile{popgreen}{Fiche bilan}{l'essentiel à retenir}\end{minipage}\par}

% Sommaire
\newtcolorbox{sommaire}{enhanced,colback=white,colframe=popink,boxrule=2.2pt,arc=10pt,
  drop shadow=popink,left=18pt,right=18pt,top=13pt,bottom=13pt,before skip=6pt,after skip=20pt,
  before upper={\poppill{Sommaire}{poppurple}{white}\par\vspace{9pt}\popsemi\fontsize{10.6}{14.5}\selectfont\color{popink}}}

% Signature de fin de cours — poussée tout en bas de la dernière page
\newcommand{\findecours}{%
  \par\vfill\nopagebreak
  \begin{center}\popsquiggle{poporange}\end{center}\vspace{4pt}
  \signatureonbuch
}
\AtBeginDocument{\def\llbracket{\mathopen{[\mkern-3.5mu[}}\def\rrbracket{\mathclose{]\mkern-3.5mu]}}} % intervalles d'entiers (glyphe ⟦ absent de la police)
% Glyphes absents de Fira Math : redéfinis à partir de glyphes disponibles
\AtBeginDocument{%
  \def\setminus{\mathbin{\backslash}}\let\smallsetminus\setminus
  \def\vdots{\mathord{\vbox{\baselineskip3.2pt\lineskiplimit0pt\kern4pt\hbox{.}\hbox{.}\hbox{.}}}}%
  \def\ddots{\mathinner{\mkern1mu\raise6.5pt\hbox{.}\mkern2mu\raise3.6pt\hbox{.}\mkern2mu\raise0.7pt\hbox{.}\mkern1mu}}%
  \def\triangle{\mathord{\scalebox{0.9}{$\symup{\Delta}$}}}%
  \def\nabla{\mathord{\raisebox{\depth}{\scalebox{1}[-1]{$\symup{\Delta}$}}}}%
  \def\checkmark{\popcheck{popdark}}%
  \def\star{\mathbin{\ast}}\def\bigstar{\mathord{\ast}}%
  \def\diamond{\mathbin{\scalebox{0.6}{$\Diamond$}}}%
  \def\bigcup{\mathop{\mathchoice{\raisebox{-0.3em}{\scalebox{1.7}{$\cup$}}}{\scalebox{1.25}{$\cup$}}{\cup}{\cup}}\displaylimits}%
  \def\bigcap{\mathop{\mathchoice{\raisebox{-0.3em}{\scalebox{1.7}{$\cap$}}}{\scalebox{1.25}{$\cap$}}{\cap}{\cap}}\displaylimits}%
  \def\lesseqqgtr{\mathrel{\lessgtr}}\def\bigoplus{\mathop{\oplus}\displaylimits}%
}
\providecommand{\ssi}{si et seulement si} % abréviation fréquente
\AtBeginDocument{\providecommand{\wideparen}{\overparen}} % arc de cercle

%% FIGFIX-BEGIN v3 — correctifs globaux des figures (docs/onbuch-generation/tools/figfix)
\usepackage{adjustbox}\usepackage{etoolbox}
% toute figure TikZ/pgfplots plus large que la ligne est réduite à la largeur disponible
\BeforeBeginEnvironment{tikzpicture}{\begin{adjustbox}{max width=\linewidth}}
\AfterEndEnvironment{tikzpicture}{\end{adjustbox}}
% légendes pgfplots lisibles par-dessus les courbes
\pgfplotsset{every axis/.append style={legend style={fill=white,fill opacity=0.92,text opacity=1,draw=black!15,inner sep=3pt}}}
% étiquettes d'arêtes (syntaxe edge["..."]) avec fond blanc
\tikzset{every edge quotes/.append style={fill=white,inner sep=1.2pt}}
% bulles de cartes mentales : texte à la ligne dans la bulle, bulle un peu plus grande
\tikzset{every concept/.append style={text width=2.0cm,align=center,minimum size=2.3cm,inner sep=2pt}}
%% FIGFIX-END
\begin{document}

\pagegarde{Lesson 57 — Vocabulary: Words and expressions related to using ICTs as recreational resources at home or in the community}{Anglais}{Tle A}{Chapter 5 : Media and Communication}{EN57}{2 h}{Cette leçon de vocabulaire dote les élèves de Terminale A du lexique précis des TIC comme ressources de loisir, à la maison ou dans la communauté (streaming, gaming, réseaux sociaux, cybercafés). Elle développe champs lexicaux, collocations, expressions idiomatiques, familles de mots et faux amis, avec réemploi systématique en contexte oral et écrit.
\vspace{4pt}
\begin{multicols}{2}\small
\begin{itemize}
\item À la fin de cette leçon, je sais nommer en anglais les principaux équipements et services TIC utilisés pour les loisirs (smartphone, game console, streaming platform...).
\item À la fin de cette leçon, je sais employer le vocabulaire des activités récréatives numériques (to stream, to download, to chat, to browse...).
\item À la fin de cette leçon, je sais utiliser les collocations courantes du thème (to play online games, to binge-watch a series, to scroll through social media...).
\item À la fin de cette leçon, je sais reconnaître et employer des expressions idiomatiques liées aux loisirs numériques (to be glued to the screen, to kill time online...).
\item À la fin de cette leçon, je sais éviter les faux amis et calques du français (console vs console, to assist vs assister, actually vs actuellement...).
\item À la fin de cette leçon, je sais construire des mots de la même famille (entertain, entertainment, entertaining, entertained...).
\end{itemize}\end{multicols}}

\begin{sommaire}
\begin{multicols}{2}
\begin{enumerate}
\item ICT equipment and places for recreation
\item Recreational activities online and offline
\item Collocations and idiomatic expressions
\item False friends and French traps
\item Word families, pronunciation and spelling
\item Using the vocabulary in context
\item Activité d'intégration
\item Méthodes et exercices résolus
\item Exercices
\item Corrigés des exercices
\item Fiche bilan
\end{enumerate}
\end{multicols}
\end{sommaire}

\begin{objectifs}
\begin{itemize}
\item À la fin de cette leçon, je sais nommer en anglais les principaux équipements et services TIC utilisés pour les loisirs (smartphone, game console, streaming platform...).
\item À la fin de cette leçon, je sais employer le vocabulaire des activités récréatives numériques (to stream, to download, to chat, to browse...).
\item À la fin de cette leçon, je sais utiliser les collocations courantes du thème (to play online games, to binge-watch a series, to scroll through social media...).
\item À la fin de cette leçon, je sais reconnaître et employer des expressions idiomatiques liées aux loisirs numériques (to be glued to the screen, to kill time online...).
\item À la fin de cette leçon, je sais éviter les faux amis et calques du français (console vs console, to assist vs assister, actually vs actuellement...).
\item À la fin de cette leçon, je sais construire des mots de la même famille (entertain, entertainment, entertaining, entertained...).
\item À la fin de cette leçon, je sais prononcer et orthographier correctement les mots clés du thème (leisure, download, Wi-Fi, website...).
\item À la fin de cette leçon, je sais réemployer ce lexique dans des phrases, dialogues et courts textes sur mes propres loisirs numériques.
\end{itemize}
\end{objectifs}

\begin{prerequis}
\begin{itemize}
\item Vocabulaire général des médias et de la communication vu en début de chapitre (media, the press, the Internet).
\item Présent simple et présent continu (habitudes et actions en cours) vus en Première.
\item Verbes à particule de base (to log in, to switch on/off) rencontrés en Première.
\item Notions de base sur les faux amis déjà signalées en classe de Première (actually, eventually).
\end{itemize}
\end{prerequis}

\newpage

\coursec{ICT equipment and places for recreation}

\begin{textebox}[Situation d'accroche — A Saturday evening in Bamenda]
It is 7 p.m. in Bamenda. Junior has just finished his homework and grabs his smartphone: thirty minutes of FIFA Mobile, then a Nollywood series on his laptop, and finally a video call with his cousin in Buea to plan the weekend. His grandmother shakes her head: “In my days, leisure meant storytelling under the moonlight!” Like Junior, millions of young Cameroonians now use ICTs — phones, computers, game consoles, the Internet — not only to study, but above all to relax and have fun, at home or in community cybercafés and gaming centres.
\end{textebox}

\textbf{Problématique.} Can you name, in correct English, all these devices, places and services, and talk about your digital leisure without falling into French traps? That is exactly what this lesson will teach you: the essential vocabulary of \cle{ICTs as recreational resources}, at home and in the community.

\begin{definition}[What are ICTs?]
\cle{ICTs} is the abbreviation of \emph{Information and Communication Technologies}: all the technologies used to handle information and to communicate — computers, mobile phones, tablets, the Internet, television and radio networks.\\
\eng{ICTs have changed the way young people spend their free time.} « Les TIC ont changé la façon dont les jeunes occupent leur temps libre. »\\
Prononciation : on épele les lettres — « aï-si-tiz ».
\end{definition}

\courssub{Observing the vocabulary in a text}

Avant d'apprendre les listes de mots, lisons un court texte authentique et relevons le lexique des TIC utilisées comme ressources de loisir.

\begin{textebox}[Reading text — Digital leisure in Yaoundé]
Every evening, after dinner, my family gathers around the smart TV to stream comedies. My father prefers the news, but my little sister always wins the remote control! On Saturdays, my friends and I meet at the gaming centre near the market to play football video games on the game consoles. The place is always crowded, and we sometimes wait an hour for a free screen.

When there is no Wi-Fi at home, I buy a data bundle and chat on WhatsApp with my classmates. My cousin, who lives in a quarter without network coverage, goes to the cybercafé twice a week: there, for a small fee, he watches music videos on YouTube and downloads films onto his laptop. Last month, the community multimedia centre near the council office organised a free film show, and more than a hundred young people came with their headphones and tablets. Digital leisure is everywhere in Yaoundé — you only need a device, a connection, and a little money for data!
\end{textebox}

\textbf{Glossaire (mots difficiles) :}
\begin{itemize}
\item \eng{to gather} (v.) : se rassembler
\item \eng{remote control} (n.) : télécommande
\item \eng{crowded} (adj.) : bondé, plein de monde
\item \eng{fee} (n.) : droit, tarif
\item \eng{network coverage} (n.) : couverture réseau
\item \eng{council office} (n.) : mairie
\end{itemize}

\textbf{Questions de compréhension :}
\begin{enumerate}
\item What does the family watch on the smart TV?
\item Where do the narrator and his friends play video games?
\item What does the narrator do when there is no Wi-Fi at home?
\item Why does the cousin go to the cybercafé?
\item What did the community multimedia centre organise last month?
\end{enumerate}

\begin{corrige}[Questions de compréhension]
\begin{enumerate}
\item The family watches comedies (and the father would prefer the news).
\item They play video games at the gaming centre near the market.
\item He buys a data bundle and chats on WhatsApp with his classmates.
\item Because there is no network coverage in his quarter; he goes there to watch music videos and download films.
\item It organised a free film show.
\end{enumerate}
\end{corrige}

\courssub{Devices: naming the equipment}

\begin{definition}[Device]
A \cle{device} (prononcé « di-\textbf{vaïce} ») is a machine or tool designed for a particular purpose, especially an electronic one: a smartphone, a tablet, a computer. Pluriel : \eng{devices} (« di-vaï-siz »). Ne pas confondre avec \eng{to devise} (v.) : concevoir, imaginer.
\end{definition}

Voici les appareils essentiels du loisir numérique. Observez d'abord les exemples, puis mémorisez le tableau.

\eng{She watches films on her tablet in bed.} « Elle regarde des films sur sa tablette au lit. »\\
\eng{My brother plays FIFA on his game console.} « Mon frère joue à FIFA sur sa console de jeux. »\\
\eng{Put on your headphones; the baby is sleeping.} « Mets tes écouteurs ; le bébé dort. »

\begin{center}
\begin{tabularx}{\linewidth}{|l|X|X|}
\hline
\rowcolor{popblueL} \textbf{English} & \textbf{Français} & \textbf{Exemple en contexte} \\
\hline
smartphone & smartphone & \eng{I chat on my smartphone.} \\
\hline
laptop & ordinateur portable & \eng{She does her homework on her laptop.} \\
\hline
tablet & tablette & \eng{He reads e-books on his tablet.} \\
\hline
desktop computer & ordinateur de bureau & \eng{The cybercafé has ten desktop computers.} \\
\hline
game console & console de jeux vidéo & \eng{The game console is connected to the TV.} \\
\hline
smart TV & télévision connectée & \eng{We stream films on the smart TV.} \\
\hline
headphones / earphones & casque / écouteurs & \eng{He listens to music with his earphones.} \\
\hline
speakers & haut-parleurs, enceintes & \eng{The speakers are too loud!} \\
\hline
\end{tabularx}
\end{center}

\begin{attention}[Pluriels et articles piégeux]
\begin{itemize}
\item \eng{headphones}, \eng{earphones} et \eng{speakers} sont presque toujours au \textbf{pluriel} (comme « des écouteurs »), car il y a deux écouteurs ou deux enceintes : \eng{My headphones are new.} et non \emph{my headphone is new}.
\item \eng{smart TV} se prononce « smarte ti-vi » ; on dit aussi \eng{a connected TV}.
\item Un \eng{laptop} se pose sur les genoux (\eng{lap} = genoux) ; un \eng{desktop} reste sur un bureau (\eng{desk}).
\end{itemize}
\end{attention}

\courssub{Places in the community: shared access to ICTs}

Quand on n'a pas d'appareil ou de connexion à la maison, on se rend dans des lieux communautaires. Voici les quatre lieux-clés.

\begin{definition}[Community places for digital leisure]
\begin{itemize}
\item \cle{cybercafé} / \cle{Internet café} : a place where you pay to use computers connected to the Internet.
\item \cle{gaming centre} : a place with game consoles and computers where people pay to play video games.
\item \cle{video club} : a shop where you can rent or watch films (very common in Cameroonian quarters).
\item \cle{community multimedia centre} : a public centre offering computers, Internet, radio and training to the local population.
\end{itemize}
\end{definition}

\eng{The gaming centre is full on Sundays.} « Le centre de jeux est plein le dimanche. »\\
\eng{She goes to the cybercafé to print her CV and check her e-mails.} « Elle va au cybercafé pour imprimer son CV et consulter ses e-mails. »\\
\eng{The video club near the stadium shows Premier League matches.} « Le vidéo-club près du stade diffuse les matchs de Premier League. »

\begin{attention}[Faux amis et calques du français]
\begin{itemize}
\item Ne traduisez JAMAIS « cybercafé » par \emph{cyber coffee} ! On dit \eng{cybercafé} ou \eng{Internet café}. \eng{Coffee} signifie « café » (la boisson).
\item \eng{data bundle} signifie « forfait internet » : \eng{I bought a 2 GB data bundle.} « J'ai acheté un forfait de 2 Go. » Ne dites pas \emph{data packet} dans ce sens (un \eng{data packet} est un « paquet de données » technique, invisible pour l'utilisateur).
\item « Un forfait » ne se dit pas \emph{a package} ni \emph{a formula} : dites \eng{a data bundle} ou \eng{a data plan}.
\item \eng{library} = bibliothèque ; une « librairie » se dit \eng{bookshop}.
\end{itemize}
\end{attention}

\courssub{Services and platforms}

\begin{definition}[Online services for recreation]
\begin{itemize}
\item \cle{streaming platform} : a service that lets you watch films and series directly online, without downloading them (Netflix, YouTube). \eng{to stream} = regarder/écouter en streaming.
\item \cle{social media} : websites and apps where people share messages, photos and videos (Facebook, WhatsApp, TikTok). Toujours au pluriel : \eng{Social media are very popular.}
\item \cle{gaming website} : a website where you play games online.
\item \cle{music app} : an application for listening to music (Spotify, Boomplay, Audiomack).
\end{itemize}
\end{definition}

\eng{We stream Nollywood films on a streaming platform.} « Nous regardons des films nollywoodiens en streaming sur une plateforme de streaming. »\\
\eng{She spends two hours a day on social media.} « Elle passe deux heures par jour sur les réseaux sociaux. »

\begin{propriete}[Noms de marques : pas d'article, pas de préposition calquée]
\begin{itemize}
\item Les noms de plateformes s'emploient \textbf{sans article} : \eng{on Facebook}, \eng{on WhatsApp}, \eng{on YouTube}, \eng{on TikTok}.
\item On dit \eng{on social media}, \eng{on the Internet} (avec \eng{the}), mais \eng{online} sans article : \eng{I watched it online.}
\item \eng{to chat on WhatsApp}, \eng{to post on Facebook}, \eng{to download from YouTube}.
\end{itemize}
\end{propriete}

\courssub{At home vs in the community}

\begin{propriete}[Deux contextes, deux vocabulaires]
\begin{itemize}
\item \textbf{At home} : you use your \cle{personal devices} (your own smartphone, laptop, smart TV) and usually a private Wi-Fi connection. \eng{At home, I watch series on my laptop.}
\item \textbf{In the community} : you use \cle{shared access} — computers and consoles that many people use — and you often pay for the connection or the time: \eng{At the cybercafé, I pay 200 francs for one hour of Internet.}
\end{itemize}
Expressions utiles : \eng{free Wi-Fi} (Wi-Fi gratuit), \eng{to pay a fee} (payer un droit), \eng{public access} (accès public), \eng{to share a computer} (partager un ordinateur).
\end{propriete}

\begin{popfigure}
\begin{center}
\begin{tikzpicture}[scale=0.9]
% Maison
\draw[thick, popblue] (0,0) rectangle (2.2,1.6);
\draw[thick, popblue] (-0.2,1.6) -- (1.1,2.6) -- (2.4,1.6);
\draw[thick, popblue] (0.9,0) rectangle (1.4,0.8);
\node[popblue, font=\bfseries] at (1.1,-0.5) {HOME};
\node[align=center, font=\small] at (1.1,3.1) {personal devices\\ private Wi-Fi};
% Cybercafé
\draw[thick, poporange] (8,0) rectangle (10.6,1.6);
\draw[thick, poporange] (7.8,1.6) rectangle (10.8,2.1);
\node[poporange, font=\bfseries] at (9.3,-0.5) {CYBERCAFÉ / GAMING CENTRE};
\node[align=center, font=\small] at (9.3,2.7) {shared access\\ paid connection};
% Flèches
\draw[-{Stealth[length=3mm]}, very thick, popgreen] (2.4,1.1) -- (7.8,1.1) node[midway, above, font=\small] {no Wi-Fi at home};
\draw[-{Stealth[length=3mm]}, very thick, poppurple] (7.8,0.5) -- (2.4,0.5) node[midway, below, font=\small] {back home with downloaded files};
\end{tikzpicture}
\end{center}
\legende{Du domicile au lieu communautaire : appareils personnels contre accès partagé.}
\end{popfigure}

\courssub{The language of connection}

\begin{definition}[Connection vocabulary]
\begin{itemize}
\item \cle{Wi-Fi} (prononcé « \textbf{ouaï-faï} ») : wireless Internet connection. \eng{Is there free Wi-Fi here?}
\item \cle{data bundle} : a paid amount of mobile Internet (un forfait). \eng{My data bundle is finished.}
\item \cle{broadband} : a fast, permanent Internet connection (le haut débit). \eng{Few homes have broadband in my quarter.}
\item \cle{to connect / to disconnect} : se connecter / se déconnecter. \eng{Connect to the Wi-Fi.} / \eng{Disconnect your phone at night.}
\item \cle{network coverage} : the area where a mobile network works (la couverture réseau). \eng{There is no network coverage in the village.}
\end{itemize}
\end{definition}

\begin{propriete}[Verbes et prépositions de la connexion]
\begin{center}
\begin{tabularx}{\linewidth}{|X|X|}
\hline
\rowcolor{popblueL} \textbf{Expression anglaise} & \textbf{Français} \\
\hline
to connect to the Wi-Fi & se connecter au Wi-Fi \\
\hline
to disconnect from the Internet & se déconnecter d'Internet \\
\hline
to log in / to log out & se connecter / se déconnecter (d'un compte) \\
\hline
to go online / to stay offline & aller en ligne / rester hors ligne \\
\hline
to run out of data & épuiser son forfait \\
\hline
to top up (one's phone) & recharger (son crédit) \\
\hline
\end{tabularx}
\end{center}
Attention à la préposition : on se connecte \eng{to} quelque chose, jamais \emph{connect on the Wi-Fi}.
\end{propriete}

\begin{attention}[Prononciations piégeuses]
\begin{itemize}
\item \eng{leisure} se prononce « \textbf{lé-djeur} » (une seule syllabe et demie) — jamais « lé-i-seur ».
\item \eng{device} : « di-\textbf{vaïce} » avec un « s » sonore écrit -ce ; ne pas confondre avec \eng{advice} (conseil, invariable) et \eng{to advise} (conseiller).
\item \eng{Wi-Fi} : « ouaï-faï », deux syllabes, accent sur la seconde.
\item \eng{data} : « déi-ta » en anglais britannique courant.
\end{itemize}
\end{attention}

\coursec{Recreational activities online and offline}

\courssub{Online activities: the key verbs}

Le loisir numérique, ce sont d'abord des \textbf{actions}. Chaque activité a son verbe et sa préposition : c'est ce qu'on appelle les \cle{collocations} (associations habituelles de mots).

\begin{center}
\begin{tabularx}{\linewidth}{|X|X|}
\hline
\rowcolor{popblueL} \textbf{Collocation anglaise} & \textbf{Français} \\
\hline
to play video games (on a console) & jouer à des jeux vidéo (sur une console) \\
\hline
to watch films / series online & regarder des films / des séries en ligne \\
\hline
to stream music & écouter de la musique en streaming \\
\hline
to download songs / films & télécharger des chansons / des films \\
\hline
to upload photos / videos & mettre en ligne des photos / des vidéos \\
\hline
to chat with friends (on WhatsApp) & discuter avec des amis (sur WhatsApp) \\
\hline
to browse the Internet / to surf the Net & naviguer sur Internet \\
\hline
to post comments / pictures & publier des commentaires / des photos \\
\hline
to share videos & partager des vidéos \\
\hline
to listen to music / podcasts & écouter de la musique / des podcasts \\
\hline
to scroll through one's news feed & faire défiler son fil d'actualité \\
\hline
to make video calls & passer des appels vidéo \\
\hline
\end{tabularx}
\end{center}

\eng{She uploads photos of the school feast every year.} « Elle met en ligne des photos de la fête de l'école chaque année. »\\
\eng{Stop scrolling through your news feed and come to dinner!} « Arrête de faire défiler ton fil d'actualité et viens dîner ! »

\begin{attention}[Prépositions à mémoriser]
\begin{itemize}
\item \eng{to listen \textbf{to} music} — jamais \emph{to listen music} (le « to » est obligatoire).
\item \eng{to chat \textbf{with} friends}, \eng{to play \textbf{with} friends}, mais \eng{to play a game} (sans préposition devant le jeu).
\item \eng{to spend time \textbf{on} social media} : \eng{He spends three hours on TikTok.} « Il passe trois heures sur TikTok. »
\item \eng{to be addicted \textbf{to} video games} : être accro aux jeux vidéo.
\end{itemize}
\end{attention}

\courssub{Offline and mixed activities}

Toutes les activités numériques ne demandent pas une connexion permanente.

\begin{definition}[Offline digital leisure]
\begin{itemize}
\item \cle{to play offline games} : jouer à des jeux hors ligne (installés sur l'appareil).
\item \cle{to watch downloaded films} : regarder des films téléchargés.
\item \cle{to listen to saved music} : écouter de la musique enregistrée.
\item \cle{to transfer files via Bluetooth} : transférer des fichiers par Bluetooth.
\item \cle{to charge one's phone} : charger son téléphone. \eng{My battery is low; I must charge my phone.}
\end{itemize}
\end{definition}

\eng{During power cuts, the children play offline games on the tablet.} « Pendant les coupures de courant, les enfants jouent à des jeux hors ligne sur la tablette. »\\
\eng{He transferred the song to my phone via Bluetooth.} « Il a transféré la chanson sur mon téléphone par Bluetooth. »

\coursec{Collocations and idiomatic expressions}

\courssub{Verb + noun collocations}

\begin{propriete}[Les collocations à connaître pour le Bac]
\begin{itemize}
\item \eng{to \textbf{spend} time online} : passer du temps en ligne (jamais \emph{to pass time}).
\item \eng{to \textbf{waste} time on games} : perdre son temps sur des jeux.
\item \eng{to \textbf{kill} time} : tuer le temps. \eng{I play games to kill time at the bus stop.}
\item \eng{to \textbf{have} fun} : s'amuser. \eng{We had a lot of fun at the gaming centre.}
\item \eng{to \textbf{take} a break (from screens)} : faire une pause (loin des écrans).
\item \eng{to \textbf{make} friends online} : se faire des amis en ligne.
\item \eng{to \textbf{keep in touch} with someone} : rester en contact avec quelqu'un.
\end{itemize}
\end{propriete}

\courssub{Idiomatic expressions about digital leisure}

\begin{definition}[Idioms in context]
\begin{itemize}
\item \cle{to be glued to one's screen} : to watch one's screen all the time (être rivé à son écran). \eng{The boy is glued to his screen all day.}
\item \cle{to binge-watch a series} : to watch many episodes of a series one after another (regarder une série en rafale). \eng{We binge-watched the whole series during the holidays.}
\item \cle{screen time} : the amount of time spent in front of a screen (le temps d'écran). \eng{Parents should limit their children's screen time.}
\item \cle{a game addict / a gaming addict} : a person who cannot stop playing (un accro aux jeux).
\item \cle{to go viral} : to spread very quickly on the Internet (devenir viral). \eng{Her dance video went viral on TikTok.}
\end{itemize}
\end{definition}

\begin{exemplebox}[Mini-dialogue — At the gaming centre]
\eng{— Hi, Paul! Are you coming to the gaming centre this afternoon?\\
— I can't. My father says I spend too much time on video games. He wants me to take a break from screens.\\
— He is right, in a way. You were glued to your screen all weekend!\\
— I know. But we had so much fun binge-playing the new football game!\\
— OK. See you online this evening, then.\\
— Sure. I'll log in after dinner.}\\
« — Salut Paul ! Tu viens au centre de jeux cet après-midi ? — Je ne peux pas. Mon père dit que je passe trop de temps sur les jeux vidéo. Il veut que je fasse une pause loin des écrans. — Il a raison, d'une certaine façon. Tu étais rivé à ton écran tout le week-end ! — Je sais. Mais on s'est tellement amusés à jouer en rafale au nouveau jeu de foot ! — Bon. À ce soir en ligne, alors. — Bien sûr. Je me connecterai après le dîner. »
\end{exemplebox}

\coursec{False friends and French traps}

\begin{attention}[Les faux amis du loisir numérique]
\begin{center}
\begin{tabularx}{\linewidth}{|X|X|X|}
\hline
\rowcolor{popblueL} \textbf{Français} & \textbf{English (correct)} & \textbf{Piège à éviter} \\
\hline
un forfait internet & a data bundle / a data plan & \emph{a package}, \emph{a formula} \\
\hline
une librairie & a bookshop & \emph{a library} (= bibliothèque) \\
\hline
une console & a game console & \emph{a consol} (n'existe pas) \\
\hline
recharger son téléphone (crédit) & to top up one's phone & \emph{to recharge} = recharger la batterie \\
\hline
recharger la batterie & to charge one's phone & \emph{to top up the battery} \\
\hline
un abonnement & a subscription & \emph{an abonnement} (n'existe pas) \\
\hline
le haut débit & broadband & \emph{high debit} \\
\hline
un écran & a screen & \emph{an ecran} (orthographe française) \\
\hline
actuellement & currently, at present & \emph{actually} (= en fait) \\
\hline
\end{tabularx}
\end{center}
\end{attention}

\begin{propriete}[Deux verbes à ne pas confondre : to charge / to top up]
\begin{itemize}
\item \eng{to charge one's phone} = brancher le téléphone pour recharger la \textbf{batterie} : \eng{I charge my phone every night.}
\item \eng{to top up one's phone} = acheter du \textbf{crédit} téléphonique : \eng{I topped up my phone with 500 francs.}
\end{itemize}
En français, « recharger » sert pour les deux ; en anglais, il faut choisir !
\end{propriete}

\coursec{Word families, pronunciation and spelling}

\courssub{Word families}

\begin{center}
\begin{tabular}{|l|l|l|}
\hline
\rowcolor{popblueL} \textbf{Nom} & \textbf{Verbe} & \textbf{Adjectif} \\
\hline
connection & to connect & connected \\
\hline
communication & to communicate & communicative \\
\hline
recreation & to recreate & recreational \\
\hline
entertainment & to entertain & entertaining \\
\hline
relaxation & to relax & relaxing / relaxed \\
\hline
addiction & --- & addicted / addictive \\
\hline
download & to download & downloadable \\
\hline
player & to play & playful \\
\hline
\end{tabular}
\end{center}

\eng{Video games can be very addictive.} « Les jeux vidéo peuvent rendre très accro. »\\
\eng{The show was really entertaining.} « Le spectacle était vraiment divertissant. »

\begin{attention}[Relaxing ou relaxed ?]
Comme tous les adjectifs en \eng{-ing} / \eng{-ed} :
\begin{itemize}
\item \eng{relaxing} décrit la \textbf{chose} qui détend : \eng{Music is relaxing.} « La musique est relaxante. »
\item \eng{relaxed} décrit la \textbf{personne} détendue : \eng{I feel relaxed after the game.} « Je me sens détendu après le match. »
\end{itemize}
Même logique : \eng{boring} (ennuyeux) / \eng{bored} (ennuyé) ; \eng{interesting} / \eng{interested}.
\end{attention}

\courssub{Pronunciation and spelling points}

\begin{propriete}[Orthographe et prononciation à retenir]
\begin{itemize}
\item \eng{game} (« guéïme ») s'écrit avec un seul \textbf{m} ; mais \eng{gaming} garde le \eng{-ing} simple : \eng{gaming centre}.
\item \eng{to download} et \eng{to upload} s'écrivent en \textbf{un seul mot} ; l'accent tombe sur la seconde syllabe pour le verbe (« daoun-\textbf{lode} ») et sur la première pour le nom (« \textbf{daoun}-lode »).
\item \eng{online} et \eng{offline} s'écrivent en un seul mot, sans trait d'union.
\item \eng{website} s'écrit en un seul mot ; \eng{web page} en deux mots.
\item \eng{e-mail} (ou \eng{email}) : les deux graphies existent ; gardez la même dans toute une copie.
\item \eng{headphones} : le \eng{h} initial se prononce (« héde-fonez »), contrairement au français.
\end{itemize}
\end{propriete}

\coursec{Using the vocabulary in context}

\courssub{Lexical mind map}

\begin{popfigure}
\begin{center}
\begin{tikzpicture}[mindmap, grow cyclic, every node/.style={concept}, concept color=popblue!60, text=white,
level 1/.append style={level distance=3.2cm, sibling angle=120},
level 2/.append style={level distance=2.2cm, sibling angle=45},
scale=0.85, transform shape]
\node[concept, font=\bfseries] {ICTs for recreation}
child[concept color=popgreen!80] { node {Devices}
  child { node {smartphone} }
  child { node {laptop} }
  child { node {game console} }
  child { node {smart TV} }
  child { node {headphones} }
}
child[concept color=poporange!85] { node {Places}
  child { node {cybercafé} }
  child { node {gaming centre} }
  child { node {video club} }
  child { node {multimedia centre} }
}
child[concept color=poppurple!75] { node {Services}
  child { node {streaming platform} }
  child { node {social media} }
  child { node {gaming website} }
  child { node {music app} }
};
\end{tikzpicture}
\end{center}
\legende{Carte mentale du lexique : les TIC comme ressources de loisir.}
\end{popfigure}

\courssub{Guided practice}

\begin{exoresolu}[Vocabulary in context]
\textbf{Énoncé.} Complete the sentences with the correct word from the lesson: \eng{data bundle}, \eng{gaming centre}, \eng{headphones}, \eng{network coverage}, \eng{stream}.
\begin{enumerate}
\item There is no \ldots\ldots\ldots in the village, so I cannot call my cousin.
\item I bought a 2 GB \ldots\ldots\ldots to chat on WhatsApp.
\item The boys meet at the \ldots\ldots\ldots to play football video games.
\item Put on your \ldots\ldots\ldots; your music is too loud.
\item We \ldots\ldots\ldots comedies on the smart TV every evening.
\end{enumerate}
\tcblower
\textbf{Solution détaillée.}
\begin{enumerate}
\item \eng{network coverage} — il est question de l'absence de réseau pour téléphoner.
\item \eng{data bundle} — « 2 GB » indique un forfait internet (jamais \emph{data packet}).
\item \eng{gaming centre} — lieu communautaire où l'on joue aux jeux vidéo.
\item \eng{headphones} — au pluriel ; le casque permet d'écouter sans déranger.
\item \eng{stream} — regarder en streaming sur une télévision connectée.
\end{enumerate}
\end{exoresolu}

\begin{exoresolu}[False friends — Correct the mistakes]
\textbf{Énoncé.} Each sentence contains one French trap. Correct it.
\begin{enumerate}
\item I bought a new package of Internet yesterday.
\item She went to the library to buy a novel.
\item I must recharge my phone: I have no more credit.
\item Actually, most students have a smartphone.
\end{enumerate}
\tcblower
\textbf{Solution détaillée.}
\begin{enumerate}
\item \eng{I bought a new \textbf{data bundle} yesterday.} — « forfait » ne se dit pas \emph{package}.
\item \eng{She went to the \textbf{bookshop} to buy a novel.} — \eng{library} = bibliothèque (on n'y achète rien).
\item \eng{I must \textbf{top up} my phone: I have no more credit.} — \eng{to recharge} concerne la batterie.
\item \eng{\textbf{Currently}, most students have a smartphone.} — \eng{actually} signifie « en fait », pas « actuellement ».
\end{enumerate}
\end{exoresolu}

\begin{exercice}[Your turn — At home or in the community?]
Classify these activities under \textbf{At home} or \textbf{In the community}, then write one full sentence for each:
\begin{enumerate}
\item playing on your own game console
\item paying for one hour at the Internet café
\item watching a match at the video club
\item chatting on WhatsApp with your personal smartphone
\item using a shared computer at the multimedia centre
\end{enumerate}
\end{exercice}

\begin{corrige}[Exercice « Your turn »]
\textbf{At home} : 1 et 4 (personal devices). \textbf{In the community} : 2, 3 et 5 (shared access, paid connection).\\
Exemples de phrases : \eng{At home, I play on my own game console.} / \eng{In the community, I pay for one hour at the Internet café.} / \eng{We watch football matches at the video club.} / \eng{At home, I chat on WhatsApp with my personal smartphone.} / \eng{Students use shared computers at the multimedia centre.}
\end{corrige}

\begin{exercice}[Word families]
Complete the table with the missing word, then use each word in a sentence of your own.
\begin{enumerate}
\item verb : to entertain $\rightarrow$ noun : \ldots\ldots\ldots
\item noun : addiction $\rightarrow$ adjective : \ldots\ldots\ldots
\item verb : to connect $\rightarrow$ noun : \ldots\ldots\ldots
\item noun : relaxation $\rightarrow$ verb : \ldots\ldots\ldots
\end{enumerate}
\end{exercice}

\begin{corrige}[Exercice « Word families »]
\begin{enumerate}
\item \eng{entertainment} — \eng{Television is a popular form of entertainment.}
\item \eng{addictive} (ou \eng{addicted}) — \eng{Some online games are very addictive.}
\item \eng{connection} — \eng{The Internet connection is slow today.}
\item \eng{to relax} — \eng{I listen to music to relax after school.}
\end{enumerate}
\end{corrige}

\begin{experience}[Speaking — A class survey on digital leisure]
Par groupes de quatre :
\begin{enumerate}
\item Préparez cinq questions avec le lexique de la leçon : \eng{How much time do you spend on social media every day?} / \eng{Do you prefer watching films at home or at the video club?} / \eng{How often do you go to the cybercafé?}
\item Interrogez vos camarades et notez les réponses.
\item Rapportez les résultats à la classe : \eng{In our group, three students out of four use a data bundle every week. Two students are glued to their screens at weekends!}
\end{enumerate}
\end{experience}

\begin{savaistu}[Digital leisure in Cameroon]
Au Cameroun, les \eng{call boxes} (les cabines téléphoniques tenues par des vendeurs de crédit) et les cybercafés restent des lieux de loisir numérique essentiels dans les quartiers où le Wi-Fi domestique est rare. Beaucoup de jeunes y téléchargent musiques et films, y rechargent leur \eng{data bundle} ou y regardent des matchs. Le mot anglais \eng{leisure} (« lé-djeur »), qui signifie « loisir », vient du vieux français \emph{leisir} (« être permis ») — un mot que le français a donné à l'anglais au Moyen Âge !
\end{savaistu}

\begin{aretenir}[Key points of this lesson]
\begin{itemize}
\item \textbf{Devices} : smartphone, laptop, tablet, desktop computer, game console, smart TV, headphones, speakers.
\item \textbf{Places} : cybercafé / Internet café (jamais \emph{cyber coffee}), gaming centre, video club, community multimedia centre.
\item \textbf{Services} : streaming platform, social media, gaming website, music app.
\item \textbf{Connection} : Wi-Fi (« ouaï-faï »), data bundle (forfait), broadband, to connect to, network coverage.
\item \textbf{Collocations} : to spend time online, to have fun, to kill time, to be glued to one's screen, to binge-watch, to go viral.
\item \textbf{Faux amis} : data bundle (\emph{pas} package), bookshop (\emph{pas} library), to top up (crédit) / to charge (batterie), currently (\emph{pas} actually).
\item \textbf{Distinction clé} : \eng{personal devices at home} vs \eng{shared access in the community}.
\end{itemize}
\end{aretenir}

\coursec{Recreational activities online and offline}

Dans la section précédente, nous avons découvert le matériel et les lieux des TIC (\eng{smartphone, laptop, cybercafé, gaming centre}…). Mais à quoi servent tous ces outils ? C'est précisément la question de cette section : quelles \emph{activités de loisir} les jeunes — à Yaoundé, à Buea ou à Londres — pratiquent-ils en ligne et hors ligne ? Pour en parler correctement en anglais, il faut maîtriser un ensemble de \cle{verbes d'action} très fréquents, chacun associé à un complément typique. Observons d'abord un texte authentique.

\courssub{Observation : a teenager's digital weekend}

\begin{textebox}[My digital weekend — a blog post by Miriam, 17, Bamenda]
Last weekend, I binge-watched a whole Nigerian series with my sister. We streamed eight episodes in two days — we just could not stop! On Saturday evening, I went live on TikTok for ten minutes, and my followers sent funny comments. One of them even shared my video with his friends.

On Sunday morning, I listened to a podcast about African music while my brother was playing an online game. He was competing against a team from Ghana, and he levelled up twice. He was so happy! In the afternoon, I downloaded three songs and uploaded the photos I had taken at my cousin's birthday party. I edited them, chose the best one and posted a selfie. My friends liked it and left nice comments.

Before sleeping, I made a video call to my aunt in Douala. We chatted for almost an hour. Then I sent a voice note to my best friend and finally logged out at half past ten. My mother says I spend too much time online — maybe she is right!
\end{textebox}

\textbf{Glossaire.} \eng{binge-watch} (v.) : regarder (une série) en boulimie, épisode après épisode ; \eng{to go live} : diffuser en direct ; \eng{followers} (n.) : abonnés ; \eng{to level up} : monter de niveau (jeu vidéo) ; \eng{a voice note} : un message vocal ; \eng{to log out} : se déconnecter.

\textbf{Comprehension questions.}
\begin{enumerate}
\item How many episodes did Miriam stream?
\item What did her followers do when she went live?
\item Who was Miriam's brother competing against?
\item What did she do with the photos of the birthday party?
\item At what time did she log out?
\end{enumerate}

\begin{corrige}[Comprehension questions]
\begin{enumerate}
\item She streamed eight episodes (in two days).
\item They sent funny comments, and one of them shared her video with his friends.
\item He was competing against a team from Ghana.
\item She edited them, chose the best one and posted a selfie.
\item She logged out at half past ten (10.30 p.m.).
\end{enumerate}
\end{corrige}

\courssub{The action verbs of digital recreation}

Le texte de Miriam contient presque tout le vocabulaire de cette leçon. Classons ces verbes selon leur \emph{intention} : regarder/écouter, jouer, communiquer, créer. C'est la meilleure façon de les mémoriser.

\textbf{Watching and listening.}

\begin{definition}[Verbes pour regarder et écouter]
\begin{itemize}
\item \cle{to stream} (films, music) : regarder ou écouter en flux continu, sans télécharger. \eng{We stream films on Friday nights.} « Nous regardons des films en streaming le vendredi soir. »
\item \cle{to watch series/films online} : regarder des séries en ligne. \eng{She watches American series online.} « Elle regarde des séries américaines en ligne. »
\item \cle{to binge-watch} : regarder une série entière d'un coup. \eng{They binge-watched the whole season.} « Ils ont enchaîné toute la saison. »
\item \cle{to listen to music/podcasts} : écouter de la musique, des podcasts. Attention à la préposition \eng{to} ! \eng{I listen to podcasts on the bus.} « J'écoute des podcasts dans le bus. »
\item \cle{to subscribe to a channel} : s'abonner à une chaîne. \eng{He subscribed to a cooking channel.} « Il s'est abonné à une chaîne de cuisine. »
\end{itemize}
\end{definition}

\textbf{Playing.}

\begin{definition}[Verbes et expressions du jeu]
\begin{itemize}
\item \cle{to play video games / online games} : jouer à des jeux vidéo / en ligne. \eng{My brother plays online games every evening.} « Mon frère joue à des jeux en ligne chaque soir. »
\item \cle{to compete against other players} : affronter d'autres joueurs. \eng{They are competing against a team from Ghana.} « Ils affrontent une équipe du Ghana. »
\item \cle{to win / to lose a game} : gagner / perdre une partie. \eng{Our team won the game.} « Notre équipe a gagné la partie. »
\item \cle{to level up} : monter de niveau. \eng{I levelled up after three hours of playing.} « Je suis monté de niveau après trois heures de jeu. »
\item \cle{multiplayer mode} (n.) : le mode multijoueur. \eng{They prefer the multiplayer mode.} « Ils préfèrent le mode multijoueur. »
\end{itemize}
\end{definition}

\textbf{Communicating.}

\begin{definition}[Verbes de communication en ligne]
\begin{itemize}
\item \cle{to chat (with friends)} : bavarder, discuter (par écrit). \eng{She chats with her cousins every evening.} « Elle discute avec ses cousins chaque soir. »
\item \cle{to make video calls} : passer des appels vidéo. \eng{We made a video call to our uncle in London.} « Nous avons passé un appel vidéo à notre oncle à Londres. »
\item \cle{to send voice notes} : envoyer des messages vocaux. \eng{He sent me a voice note this morning.} « Il m'a envoyé un message vocal ce matin. »
\item \cle{to join a group / a forum} : rejoindre un groupe / un forum. \eng{He joined a football forum.} « Il a rejoint un forum de football. »
\item \cle{to comment / to like / to share / to follow} : commenter, aimer, partager, suivre (quelqu'un). \eng{My friends liked and shared my post.} « Mes amis ont aimé et partagé ma publication. »
\end{itemize}
\end{definition}

\textbf{Creating and managing files.}

\begin{definition}[Créer, télécharger, publier]
\begin{itemize}
\item \cle{to download} : télécharger (vers son appareil) ; \cle{to upload} : mettre en ligne (depuis son appareil). \eng{I downloaded three songs and uploaded my photos.} « J'ai téléchargé trois chansons et mis mes photos en ligne. »
\item \cle{to post} : publier (une photo, un message). \eng{She posted a selfie.} « Elle a publié un selfie. »
\item \cle{to take selfies} : prendre des selfies.
\item \cle{to edit photos/videos} : retoucher, monter des photos ou des vidéos.
\item \cle{to record skits} : enregistrer des sketchs.
\item \cle{to go live} : diffuser en direct. \eng{He went live on TikTok.} « Il a fait un direct sur TikTok. »
\item \cle{to browse the Net / to surf the Internet} : naviguer sur Internet. \eng{I surf the Internet after school.} « Je navigue sur Internet après l'école. »
\end{itemize}
\end{definition}

\begin{attention}[Download ou upload ? Le sens de la flèche]
\eng{To download} = faire venir un fichier \emph{vers} ton appareil (une chanson, un film). \eng{To upload} = envoyer un fichier \emph{depuis} ton appareil vers Internet (une photo, une vidéo). Astuce : \eng{down} = vers le bas, vers toi ; \eng{up} = vers le haut, vers le réseau. Ne les inverse pas : on \eng{downloads} une chanson et on \eng{uploads} une photo.
\end{attention}

\begin{popfigure}
\begin{tikzpicture}[node distance=0.15cm]
\node[draw=popblue, fill=popblueL, rounded corners, font=\bfseries, minimum width=6.2cm] (t1) {Watching \& Listening};
\node[draw=poppurple, fill=poppurpleL, rounded corners, font=\bfseries, minimum width=6.2cm, right=0.6cm of t1] (t2) {Playing \& Creating};
\node[draw=popline, fill=popcream, rounded corners, align=left, minimum width=6.2cm, below=0.25cm of t1, anchor=north west, text width=5.8cm] (c1) {to stream a film\\to binge-watch a series\\to listen to a podcast\\to subscribe to a channel\\to download a song};
\node[draw=popline, fill=popcream, rounded corners, align=left, minimum width=6.2cm, below=0.25cm of t2, anchor=north east, text width=5.8cm] (c2) {to play online games\\to compete against a team\\to level up\\to post a selfie\\to edit a video\\to go live};
\end{tikzpicture}
\legende{Classer les verbes par intention pour mieux les mémoriser.}
\end{popfigure}

\courssub{Grammar point : regular verbs and spelling traps}

\begin{propriete}[Conjugaison de ces verbes au prétérit]
La plupart des verbes du numérique sont \textbf{réguliers} : on ajoute \eng{-ed}.
\begin{center}
\begin{tabular}{|l|l|l|}
\hline
\rowcolor{popblueL} \textbf{Base verbale} & \textbf{Prétérit / participe} & \textbf{Remarque} \\
\hline
stream & streamed & régulier \\
\hline
download / upload & downloaded / uploaded & régulier \\
\hline
post / like / share & posted / liked / shared & \eng{like} et \eng{share} : e final + d \\
\hline
chat & chatted & \textbf{doublement du t} (une syllabe, CVC) \\
\hline
log (in/out) & logged & \textbf{doublement du g} \\
\hline
subscribe & subscribed & e final + d \\
\hline
win / lose & won / lost & \textbf{irréguliers !} \\
\hline
\end{tabular}
\end{center}
Règle du doublement : un verbe d'une syllabe terminée par \emph{consonne-voyelle-consonne} double la consonne finale : \eng{chat $\rightarrow$ chatted}, \eng{log $\rightarrow$ logged}, \eng{stop $\rightarrow$ stopped}. En revanche, \eng{stream} se termine par deux consonnes (\emph{-am} précédé de deux voyelles) : pas de doublement, \eng{streamed}.
\end{propriete}

\begin{attention}[Faux amis et calques du français]
\begin{itemize}
\item \eng{to surf the Internet} est correct, mais on ne dit \textbf{jamais} \eng{to navigate on Internet} (calque du français « naviguer sur Internet »). On dit \eng{to browse the Internet} ou \eng{to surf the Net}.
\item \eng{to listen} exige la préposition \eng{to} : on dit \eng{I listen to music}, jamais \eng{I listen music}.
\item \eng{to compete} se construit avec \eng{against} : \eng{to compete against a team}, pas \eng{with} dans ce sens.
\item \eng{to spend time online} : « passer du temps » se traduit par \eng{to spend}, pas par \eng{to pass}. \eng{She spends two hours online every day.} « Elle passe deux heures en ligne chaque jour. »
\item « Dans le bus » se dit \eng{on the bus} (et non \eng{in the bus}) : \eng{I listen to podcasts on the bus.}
\end{itemize}
\end{attention}

\begin{methode}[Mémoriser les verbes du numérique]
\begin{enumerate}
\item \textbf{Classe} chaque verbe par intention : \emph{watch / play / communicate / create}. Un verbe rangé est un verbe retenu.
\item \textbf{Apprends chaque verbe avec son complément typique} (collocation) : \eng{to stream + a film}, \eng{to download + a song}, \eng{to post + a photo}, \eng{to make + a video call}, \eng{to send + a voice note}.
\item \textbf{Fabrique une phrase personnelle} pour chaque verbe : \eng{I download gospel songs every Sunday.}
\item \textbf{Réutilise} les verbes à l'oral : raconte ta soirée d'hier en cinq phrases.
\end{enumerate}
\end{methode}

\courssub{Frequency and duration : talking about screen time}

Pour décrire ses habitudes, on utilise des expressions de fréquence et de durée :

\begin{center}
\begin{tabularx}{\linewidth}{|X|X|X|}
\hline
\rowcolor{popblueL} \textbf{English} & \textbf{Français} & \textbf{Exemple} \\
\hline
to spend hours online & passer des heures en ligne & \eng{He spends hours online.} \\
\hline
screen time (n.) & le temps d'écran & \eng{My screen time is too high.} \\
\hline
to be online / offline & être en ligne / hors ligne & \eng{Are you online now?} \\
\hline
to log in / to log out & se connecter / se déconnecter & \eng{Log in with your password.} \\
\hline
every day / at night & chaque jour / le soir & \eng{She chats at night.} \\
\hline
\end{tabularx}
\end{center}

\begin{exemplebox}[Exemples traduits]
\eng{She spends two hours online every day.} « Elle passe deux heures en ligne chaque jour. »

\eng{They are competing against a team from Ghana.} « Ils affrontent une équipe du Ghana. »

\eng{I logged out at half past ten.} « Je me suis déconnecté à dix heures et demie. »

\eng{We binge-watched a Cameroonian comedy series.} « Nous avons enchaîné tous les épisodes d'une série comique camerounaise. »
\end{exemplebox}

\begin{popfigure}
\begin{tikzpicture}[font=\small]
\draw[-{Stealth}, thick, popdark] (0,0) -- (12.5,0);
\foreach \x/\h/\act in {0.5/{6 p.m.}/{log in}, 3.5/{7 p.m.}/{stream a film}, 7/{9 p.m.}/{video call}, 11/{10 p.m.}/{log out}}{
  \fill[poporange] (\x,0) circle (0.09);
  \node[above, popblue, font=\bfseries] at (\x,0.15) {\h};
  \node[below, align=center, text width=2.2cm] at (\x,-0.15) {\act};
}
\end{tikzpicture}
\legende{A typical digital evening : log in, stream, call, log out.}
\end{popfigure}

\begin{savaistu}[Le mot « binge-watch »]
Le verbe \eng{to binge-watch} est né avec les plateformes de streaming : \eng{binge} signifie « excès, frénésie ». Regarder huit épisodes d'affilée, comme Miriam, c'est exactement cela ! Le mot est entré dans les grands dictionnaires anglais dans les années 2010 — preuve que le vocabulaire du numérique évolue très vite. Au Cameroun, les séries nigérianes (Nollywood) sont parmi les plus « binge-watchées » par les jeunes.
\end{savaistu}

\courssub{Practice}

\begin{exoresolu}[Complete with the correct verb]
Énoncé. Complète avec : \eng{stream, chatted, download, levelled up, went live}.
\begin{enumerate}
\item Yesterday, my brother \dots\ in his favourite game.
\item We \dots\ on WhatsApp for an hour last night.
\item Can I \dots\ this song on your phone?
\item The comedian \dots\ on Facebook and 500 fans watched him.
\item They \dots\ films every weekend.
\end{enumerate}
\tcblower
Solution détaillée.
\begin{enumerate}
\item \eng{levelled up} — contexte de jeu vidéo (\eng{game}) ; verbe régulier.
\item \eng{chatted} — \eng{last night} appelle le prétérit ; attention au doublement du \emph{t}.
\item \eng{download} — après \eng{can}, base verbale ; on télécharge une chanson \emph{vers} le téléphone.
\item \eng{went live} — prétérit irrégulier de \eng{go live} ; les fans regardent un direct.
\item \eng{stream} — habitude au présent (\eng{every weekend}) ; on regarde des films en streaming.
\end{enumerate}
\end{exoresolu}

\begin{exercice}[Your turn]
\begin{enumerate}
\item Traduis : « Hier soir, j'ai passé deux heures en ligne, puis je me suis déconnecté à onze heures. »
\item Trouve l'erreur : \eng{She navigated on Internet and listened music.}
\item Écris quatre phrases sur ta soirée d'hier avec : \eng{log in, stream, chat, log out}.
\end{enumerate}
\end{exercice}

\begin{corrige}[Exercice « Your turn »]
\begin{enumerate}
\item \eng{Last night, I spent two hours online, then I logged out at eleven o'clock.}
\item Deux erreurs : on dit \eng{She browsed the Internet} (ou \eng{surfed the Net}) et \eng{listened to music} (préposition \eng{to} obligatoire).
\item Modèle : \eng{Yesterday I logged in at 7 p.m. I streamed a Nigerian film with my sister. Then I chatted with my best friend for an hour. I logged out at 10 p.m.}
\end{enumerate}
\end{corrige}

\begin{aretenir}[L'essentiel de la section]
\begin{itemize}
\item Chaque verbe du numérique s'apprend avec son complément : \eng{to stream a film, to download a song, to post a photo, to make a video call}.
\item Ces verbes sont réguliers, mais attention : \eng{chat $\rightarrow$ chatted}, \eng{log $\rightarrow$ logged} ; \eng{win $\rightarrow$ won} et \eng{lose $\rightarrow$ lost} sont irréguliers.
\item Jamais \eng{navigate on Internet} : dis \eng{browse the Internet} ou \eng{surf the Net}.
\item \eng{to listen to}, \eng{to compete against}, \eng{to spend time online} : retiens les prépositions et les constructions.
\end{itemize}
\end{aretenir}

\coursec{Collocations and idiomatic expressions}

Dans la section précédente, nous avons découvert les activités récréatives en ligne et hors ligne : \eng{playing video games}, \eng{watching films}, \eng{chatting with friends}, \eng{listening to music}. Mais savoir nommer une activité ne suffit pas : pour parler comme un anglophone, il faut savoir \emph{combiner} les mots correctement. On ne dit pas \eng{do a game} ni \eng{make a game} : on dit \eng{play a game}. Ces combinaisons naturelles s'appellent des \cle{collocations}, et elles sont au cœur de cette section, avec les \cle{expressions idiomatiques} qui donnent du naturel à votre expression orale et écrite — un atout décisif pour l'essai du Bac.

\courssub{What is a collocation?}

\begin{definition}[Collocation]
A \cle{collocation} is a natural combination of words that native speakers use together. The words « choose » each other: we say \eng{play a game} (correct), never \eng{do a game} or \eng{make a game}. We say \eng{waste time}, never \eng{lose time} (which means « être en retard »). Learning collocations means learning words \emph{in chunks}, not in isolation.
\end{definition}

\begin{attention}[Pourquoi les francophones se trompent]
Le français et l'anglais ne combinent pas les mêmes verbes avec les mêmes noms. Le français dit « faire un jeu », « passer du temps », « perdre son temps » ; l'anglais dit \eng{play a game}, \eng{spend time}, \eng{waste time}. Traduire mot à mot produit des fautes typiques : \eng{*do a game}, \eng{*pass time}, \eng{*loose time}. La seule solution : mémoriser chaque nom \emph{avec} son verbe, comme on mémorise un nom français avec son genre.
\end{attention}

\begin{center}
\begin{tabularx}{\linewidth}{|X|X|X|}\hline
\rowcolor{popblueL} Français (calque à éviter) & English (correct) & Remarque \\ \hline
faire un jeu & to \cle{play} a game & jamais \eng{do/make a game} \\ \hline
passer du temps en ligne & to \cle{spend} time online & \eng{pass time} est un calque \\ \hline
perdre son temps & to \cle{waste} time & \eng{lose time} = être en retard \\ \hline
regarder un match en direct & to \cle{watch} a live match & \eng{see} est possible mais moins naturel \\ \hline
faire le buzz & to \cle{go viral} & expression idiomatique \\ \hline
\end{tabularx}
\end{center}

\courssub{Observing collocations in context}

Lisons d'abord un court texte authentique. Repérez les combinaisons verbe + nom et les expressions idiomatiques : elles sont signalées par \cle{...}.

\begin{textebox}[A Saturday at the cybercafé — reading text]
Every Saturday afternoon, the small cybercafé near the post office in Buea fills up with teenagers. Some come to \cle{play online} games; others just want to \cle{kill time} on their phones before evening classes. Sixteen-year-old Collins is the champion of the place: he loves to \cle{play against} older students, and he usually wins. “I only \cle{play for fun},” he says, “but my little brother is \cle{glued to the screen} all day. He is completely \cle{hooked on} that racing game. My mother says he \cle{wastes his data} and should kill time reading instead!”

In a corner, two girls \cle{binge-watch} a Nigerian series on a shared laptop, while their friend Sandra prefers to \cle{stream a film} on her phone. “I \cle{spend hours online} every weekend,” she admits. “Yesterday I watched a live match, then I \cle{scrolled through my feed} for an hour, \cle{updated my status} and \cle{tagged my friends} in our photos. One of our comedy skits \cle{went viral} overnight — two thousand shares!”

The owner, Mr Ekane, is not worried. “ICTs are fine as recreational resources,” he explains, “as long as young people do not become \cle{screen addicts}. Play, watch, chat — but also read a book sometimes!”

\emph{Glossary}: skit (n.) = petit sketch comique ; share (n.) = partage ; owner (n.) = propriétaire ; as long as = tant que ; feed (n.) = fil d'actualité.

\emph{Comprehension questions}: 1. What does Collins do at the cybercafé? 2. Why is Collins's mother unhappy with his little brother? 3. List three things Sandra did online yesterday. 4. What happened to the comedy skit? 5. What advice does Mr Ekane give?
\end{textebox}

\begin{corrige}[Comprehension questions — réponses attendues]
1. \eng{He plays online games and plays against older students.} 2. \eng{Because he is glued to the screen all day and wastes his data.} 3. \eng{She watched a live match, scrolled through her feed, updated her status and tagged her friends in photos.} 4. \eng{It went viral overnight — it got two thousand shares.} 5. \eng{Young people should enjoy ICTs but not become screen addicts; they should also read books sometimes.}
\end{corrige}

\courssub{Collocations with PLAY}

\begin{propriete}[Les combinaisons de PLAY]
\begin{itemize}
\item \eng{play a game} — jouer à un jeu. \eng{The children are playing a game of Ludo.} « Les enfants jouent au Ludo. »
\item \eng{play online} — jouer en ligne. \eng{Thousands of Cameroonians play online every evening.} « Des milliers de Camerounais jouent en ligne chaque soir. »
\item \eng{play against someone} — jouer contre quelqu'un. \eng{Our school team played against GBHS Bamenda.} « Notre équipe scolaire a joué contre le GBHS de Bamenda. »
\item \eng{play for fun} — jouer pour le plaisir (pas pour l'argent). \eng{I don't bet; I just play for fun.} « Je ne parie pas ; je joue juste pour le plaisir. »
\end{itemize}
Remarque : on dit \eng{play} + sport/jeu sans article (\eng{play football}, \eng{play chess}), mais \eng{play} + \textbf{the} + instrument de musique (\eng{play the guitar}).
\end{propriete}

\courssub{Collocations with WATCH and STREAM}

\begin{propriete}[Regarder à l'ère du streaming]
\begin{itemize}
\item \eng{watch a live match} — regarder un match en direct. \eng{We watched a live match at the video club in Douala.} « Nous avons regardé un match en direct au club vidéo de Douala. »
\item \eng{stream a film} — regarder un film en streaming. \eng{She streamed a film on her tablet.} « Elle a regardé un film en streaming sur sa tablette. »
\item \eng{binge-watch a series} — regarder une série en rafale (plusieurs épisodes d'affilée). \eng{We binge-watched the whole series during the holidays.} « Nous avons enchaîné toute la série pendant les vacances. »
\end{itemize}
Distinction utile : \eng{watch} implique une attention soutenue (un match, un film), \eng{see} est plus général, et \eng{look at} s'emploie pour une image fixe (\eng{look at this photo}).
\end{propriete}

\begin{savaistu}[D'où vient « binge-watching » ?]
\eng{Binge-watching} (regarder plusieurs épisodes d'affilée) vient de \eng{binge}, qui signifie « excès, beuverie » (à l'origine \eng{a drinking binge} = une beuverie). Le mot est entré dans les grands dictionnaires anglais vers 2013, avec l'essor des plateformes de streaming qui publient des saisons entières d'un coup. Il s'emploie aussi comme nom : \eng{We had a binge-watch last night.} Prononciation : « bin-tch ouo-tching ».
\end{savaistu}

\courssub{Collocations with SPEND and WASTE}

\begin{propriete}[SPEND et WASTE : la structure]
\eng{spend} + temps/argent/données + \textbf{on} + nom, ou + \textbf{-ing} : \eng{I spend two hours on my phone.} / \eng{I spend two hours chatting.} « Je passe deux heures à discuter. »

\eng{waste} + temps/argent/données + \textbf{on} + nom, ou + \textbf{-ing} : \eng{Don't waste your data on useless videos.} « Ne gaspille pas tes données en vidéos inutiles. »

\eng{kill time} — tuer le temps : \eng{He killed time on his phone while waiting for the bendskin.} « Il a tué le temps sur son téléphone en attendant la moto-taxi. »

Attention, verbes irréguliers : \eng{spend — spent — spent} ; \eng{waste} est régulier (\eng{waste — wasted — wasted}).
\end{propriete}

\begin{popfigure}
\begin{tikzpicture}[
  central/.style={rectangle, rounded corners, draw=popblue, fill=popblueL, thick, align=center, font=\bfseries, minimum width=2.2cm, minimum height=1cm},
  comp/.style={rectangle, rounded corners, draw=poporange, fill=poporangeL, align=center, minimum width=2.4cm},
  prep/.style={rectangle, rounded corners, draw=popgreen, fill=popgreenL, align=center, minimum width=2cm},
  fleche/.style={-{Stealth[length=3mm]}, thick, popdark}]
\node[central] (s) {SPEND};
\node[comp, above left=1.1cm and 2.6cm of s] (c1) {time};
\node[comp, left=1.1cm of s] (c2) {hours};
\node[comp, below left=1.1cm and 2.6cm of s] (c3) {data / the evening};
\node[prep, right=2.2cm of s, align=center] (p1){on + noun\\ \eng{on games}};
\node[prep, above right=0.4cm and 2.2cm of s, align=center] (p2){-ing\\ \eng{playing}};
\node[prep, below right=0.4cm and 2.2cm of s, align=center] (p3){online\\ \eng{spend hours online}};
\draw[fleche] (s) -- (c1);
\draw[fleche] (s) -- (c2);
\draw[fleche] (s) -- (c3);
\draw[fleche] (s) -- (p1);
\draw[fleche] (s) -- (p2);
\draw[fleche] (s) -- (p3);
\end{tikzpicture}
\legende{Organigramme : le verbe \eng{spend}, ses compléments possibles et ses constructions.}
\end{popfigure}

\courssub{Collocations with social media}

\begin{propriete}[Le vocabulaire des réseaux sociaux]
\begin{itemize}
\item \eng{scroll through one's feed} — faire défiler son fil d'actualité. \eng{She scrolled through her feed during the break.} « Elle a fait défiler son fil pendant la récréation. »
\item \eng{update one's status} — mettre à jour son statut. \eng{He updated his status after the match.} « Il a mis à jour son statut après le match. »
\item \eng{go viral} — devenir viral, faire le buzz. \eng{Their comedy skit went viral overnight.} « Leur sketch comique a fait le buzz en une nuit. »
\item \eng{tag a friend} — identifier un ami (sur une photo, une publication). \eng{She tagged three friends in the photo.} « Elle a identifié trois amis sur la photo. »
\end{itemize}
\end{propriete}

\begin{exemplebox}[Exemples traduits]
\eng{Their comedy skit went viral overnight.} « Leur sketch comique a fait le buzz en une nuit. »

\eng{Stop being glued to your phone!} « Arrête d'être rivé à ton téléphone ! »

\eng{I updated my status and tagged my best friend.} « J'ai mis à jour mon statut et identifié mon meilleur ami. »
\end{exemplebox}

\courssub{Idiomatic expressions}

\begin{definition}[Idiomatic expression]
An \cle{idiomatic expression} (or \cle{idiom}) is a fixed phrase whose meaning cannot be guessed from the individual words: \eng{to be glued to the screen} does not involve real glue (« colle ») — it means « être rivé à l'écran, incapable de s'en détacher ».
\end{definition}

\begin{center}
\begin{tabularx}{\linewidth}{|X|X|X|}\hline
\rowcolor{popblueL} Idiomatic expression & Literal sense & Équivalent français \\ \hline
to be glued to the screen & être collé à l'écran & être rivé à l'écran \\ \hline
to be hooked on a game & être accroché à un jeu & être accro à un jeu \\ \hline
to kill time online & tuer le temps en ligne & passer/tuer le temps en ligne \\ \hline
to be a screen addict & être un drogué de l'écran & être accro aux écrans \\ \hline
to be a gaming freak & être un fou des jeux & être un mordu de jeux vidéo \\ \hline
to go viral & devenir « viral » & faire le buzz, devenir viral \\ \hline
\end{tabularx}
\end{center}

\begin{attention}[Faux ami : « to go viral »]
\eng{To go viral} signifie « faire le buzz, se répandre très vite sur Internet », et non « devenir dangereux » ou « attraper un virus ». Attention : \eng{viral} ne se traduit pas toujours par « viral » au sens médical. \eng{A viral video} = « une vidéo qui a fait le buzz » ; \eng{a viral infection} = « une infection virale ». Le contexte décide.
\end{attention}

\courssub{Register: hooked on vs addicted to}

\begin{propriete}[Nuances de registre — essentiel pour l'essai du Bac]
\eng{to be hooked on} (un jeu, une série) est \textbf{familier} : parfait dans un dialogue ou une lettre à un ami. \eng{My brother is hooked on that racing game.}

\eng{to be addicted to} est \textbf{neutre à soutenu} : c'est le terme à employer dans un essai ou un débat formel. \eng{Many teenagers are addicted to social media.} « Beaucoup d'adolescents sont dépendants des réseaux sociaux. »

De même, \eng{a gaming freak} est familier et plaisant, tandis que \eng{a screen addict} est neutre et peut servir dans un texte argumentatif sur les dangers des TIC.
\end{propriete}

\begin{exoresolu}[Choose the correct collocation]
Énoncé — Complete each sentence with the correct collocation from the box: \emph{play against}, \emph{waste}, \emph{went viral}, \emph{glued to}, \emph{binge-watched}, \emph{tagged}.

1. Our team will \dots\ GBHS Yaoundé in the final.
2. Don't \dots\ your data on silly videos!
3. Their dance video \dots\ in just two days.
4. My sister is always \dots\ the screen of her phone.
5. We \dots\ the whole series last weekend.
6. She \dots\ me in her birthday photos.
\tcblower
Solution détaillée :

1. \eng{play against} — on joue \emph{contre} une équipe : \eng{Our team will play against GBHS Yaoundé in the final.}
2. \eng{waste} — collocation \eng{waste data} (gaspiller ses données) ; \eng{lose} serait un calque du français.
3. \eng{went viral} — passé irrégulier de \eng{go viral} (go — went — gone) : la vidéo « a fait le buzz ».
4. \eng{glued to} — expression idiomatique \eng{be glued to the screen} (être rivé à l'écran).
5. \eng{binge-watched} — verbe régulier au prétérit : regarder toute une série d'affilée.
6. \eng{tagged} — \eng{tag a friend} : identifier quelqu'un sur une photo ; noter le doublement du \emph{g} (\eng{tag} $\rightarrow$ \eng{tagged}).
\end{exoresolu}

\begin{attention}[Pièges à éviter]
\begin{itemize}
\item Jamais \eng{*do a game} ni \eng{*make a game} : toujours \eng{play a game}.
\item Jamais \eng{*pass time} : dites \eng{spend time} (passer du temps) ou \eng{kill time} (tuer le temps).
\item \eng{addicted} se construit avec \textbf{to}, pas avec \eng{of} : \eng{addicted to video games}.
\item \eng{hooked} se construit avec \textbf{on} : \eng{hooked on a series}.
\item Orthographe : \eng{tag} $\rightarrow$ \eng{tagged} ; \eng{chat} $\rightarrow$ \eng{chatted} (doublement de la consonne finale après une voyelle courte accentuée).
\end{itemize}
\end{attention}

\begin{exercice}[Your turn — speaking and writing]
1. Traduisez en anglais : « Mon petit frère est accro aux jeux vidéo ; il passe des heures en ligne et gaspille ses données. » (deux versions : une familière avec \eng{hooked on}, une soutenue avec \eng{addicted to}).
2. Rédigez cinq phrases sur vos loisirs numériques en utilisant au moins quatre collocations de cette section (\eng{stream}, \eng{scroll through}, \eng{play for fun}, \eng{kill time}, \eng{update my status}\dots).
3. À l'oral, par binôme : l'un décrit un « screen addict » imaginaire, l'autre lui donne des conseils.
\end{exercice}

\begin{corrige}[Exercice — éléments de réponse]
1. Version familière : \eng{My little brother is hooked on video games; he spends hours online and wastes his data.} Version soutenue : \eng{My little brother is addicted to video games; he spends hours online and wastes his data.}
2. Modèle : \eng{At weekends, I stream films on my phone and scroll through my feed. I play online games for fun, and I kill time chatting with my friends. Sometimes I update my status after a football match.}
3. Exemple de conseil : \eng{You should not be glued to the screen all day; kill time reading instead!}
\end{corrige}

\begin{aretenir}[L'essentiel de la section]
Une \cle{collocation} est une combinaison naturelle de mots : \eng{play a game}, \eng{watch a live match}, \eng{stream a film}, \eng{binge-watch a series}, \eng{spend time online}, \eng{waste data}, \eng{scroll through one's feed}, \eng{go viral}, \eng{tag a friend}. Les idiomes \eng{be glued to the screen}, \eng{be hooked on}, \eng{kill time} et \eng{be a screen addict} enrichissent l'expression. Pour l'essai du Bac, préférez le registre neutre : \eng{addicted to} plutôt que \eng{hooked on}.
\end{aretenir}

\coursec{False friends and French traps}

Dans la section précédente, nous avons appris à combiner correctement les mots entre eux grâce aux \cle{collocations} (\eng{to play online games}, \eng{to stream a film}, \eng{to scroll through a feed}) et aux expressions idiomatiques (\eng{to be glued to the screen}). Mais il existe un danger encore plus sournois pour le francophone : les mots qui \emph{ressemblent} à des mots français mais qui ne \emph{signifient} pas la même chose. Ce sont les fameux \cle{faux amis} (\eng{false friends}). Dans le domaine des TIC et des loisirs numériques, ils sont particulièrement nombreux. Cette section vous apprend à les repérer, à les corriger et à ne plus jamais vous faire piéger.

\courssub{Observing the trap: a short text}

\begin{textebox}[A weekend of digital recreation — with mistakes!]
Last weekend, I assisted a gaming tournament at the community centre in Buea. Actually, I go there every Saturday because it is my favourite place for recreation. My cousin came with me; he is very sensible about video games, so he played for five hours without stopping! After the tournament, we went to the library near the market to buy a new game for my console. The shop assistant was very kind and he consoled us with a free poster. At home, I made a photo of the poster and posted it in Internet. Then I looked at TV until midnight. What a wonderful weekend!
\end{textebox}

\textbf{Glossaire :} \emph{tournament} (nom, tournoi) ; \emph{community centre} (nom, centre communautaire, maison des jeunes) ; \emph{shop assistant} (nom, vendeur, vendeuse) ; \emph{poster} (nom, affiche).

\textbf{Questions de compréhension :}
\begin{enumerate}
\item Where did the narrator go last Saturday?
\item What did the cousin do for five hours?
\item What did the narrator do with the poster at home?
\end{enumerate}

\textbf{Réponses attendues :} 1. He went to the community centre in Buea. 2. He played video games for five hours without stopping. 3. He took a photo of it and posted it online.

Ce texte contient \textbf{huit erreurs} typiques du francophone : \eng{assisted}, \eng{Actually}, \eng{sensible}, \eng{library}, \eng{consoled}, \eng{made a photo}, \eng{in Internet}, \eng{looked at TV}. Nous allons les corriger une par une.

\courssub{The great false friends of the ICT topic}

\begin{definition}[False friend (faux ami)]
Un \cle{faux ami} est un mot anglais qui ressemble par sa forme à un mot français mais dont le sens est différent. Exemple : \eng{actually} ne veut pas dire « actuellement » mais « en fait, réellement ».
\end{definition}

\begin{propriete}[Les cinq faux amis à connaître par cœur]
\begin{itemize}
\item \eng{a (game) console} = une \textbf{console de jeux}. Mais \eng{to console someone} = \textbf{consoler} quelqu'un (réconforter). \eng{My parents bought me a game console.} « Mes parents m'ont acheté une console de jeux. » / \eng{She consoled her friend who had lost the game.} « Elle a consolé son amie qui avait perdu la partie. »
\item \eng{actually} = \textbf{en fait, réellement}. Pour « actuellement », dites \eng{currently}, \eng{at present} ou \eng{at the moment}. \eng{I actually prefer offline games.} « En fait, je préfère les jeux hors ligne. »
\item \eng{to assist} = \textbf{aider} (sens soutenu). « Assister à » (être présent) se dit \eng{to attend}. \eng{He attended a video game competition in Yaoundé.} « Il a assisté à une compétition de jeux vidéo à Yaoundé. »
\item \eng{a library} = une \textbf{bibliothèque} (on y emprunte des livres). Une « librairie » (on y achète) se dit \eng{a bookshop}. \eng{a digital library} = une bibliothèque numérique.
\item \eng{sensible} = \textbf{raisonnable, sensé}. « Sensible » (émotif, fragile) se dit \eng{sensitive}. \eng{a sensible use of screen time} = « un usage raisonnable du temps d'écran ».
\end{itemize}
\end{propriete}

\begin{exemplebox}[Exemples traduits]
\eng{I am currently watching a documentary about hackers.} « Je regarde actuellement un documentaire sur les pirates informatiques. »

\eng{Thousands of fans attended the e-sports final in Douala.} « Des milliers de fans ont assisté à la finale d'e-sport à Douala. »

\eng{The technician assisted me when my laptop crashed.} « Le technicien m'a aidé quand mon ordinateur portable a planté. »

\eng{You can borrow e-books from the digital library for free.} « Tu peux emprunter gratuitement des livres numériques à la bibliothèque numérique. »

\eng{Be sensible: do not spend all night on social media.} « Sois raisonnable : ne passe pas toute la nuit sur les réseaux sociaux. »

\eng{Teenagers are often sensitive to negative comments online.} « Les adolescents sont souvent sensibles aux commentaires négatifs en ligne. »
\end{exemplebox}

\begin{attention}[Le piège de « actually »]
\eng{I am actually watching a film.} signifie « Je suis en train de regarder un film, \textbf{en fait} » — et non « actuellement ». Pour dire « actuellement », utilisez \eng{currently} ou \eng{at present} : \eng{I am currently watching a film.} Retenez : \eng{actually} répond à « en réalité », jamais à « en ce moment ». Prononciation : \eng{actually} se dit « ak-tchou-a-li » et \eng{currently} « keu-reunt-li ».
\end{attention}

\courssub{Summary table and mind map}

\begin{center}
\begin{tabularx}{\linewidth}{|X|X|X|}
\hline
\rowcolor{popblueL} Faux ami & Sens anglais réel & Mot anglais pour le sens français \\
\hline
actually & en fait, réellement & currently, at present \\
\hline
to assist & aider & to attend (assister à) \\
\hline
a library & une bibliothèque & a bookshop (librairie) \\
\hline
sensible & raisonnable & sensitive (sensible) \\
\hline
a console & une console de jeux & to console (consoler) \\
\hline
to demand & exiger & to ask for (demander) \\
\hline
a deception & une tromperie & a disappointment (déception) \\
\hline
\end{tabularx}
\end{center}

\begin{popfigure}
\begin{tikzpicture}[scale=0.95, every node/.style={font=\small}]
\draw[fill=popblueL, draw=popblue] (0,0) node[align=center] {actually};
\draw[fill=popgreenL, draw=popgreen] (4,0) node[align=center] {en fait};
\draw[fill=poporangeL, draw=poporange] (8,0) node[align=center] {currently};
\draw[fill=popblueL, draw=popblue] (0,-1.6) node[align=center] {to assist};
\draw[fill=popgreenL, draw=popgreen] (4,-1.6) node[align=center] {aider};
\draw[fill=poporangeL, draw=poporange] (8,-1.6) node[align=center] {to attend};
\draw[fill=popblueL, draw=popblue] (0,-3.2) node[align=center] {library};
\draw[fill=popgreenL, draw=popgreen] (4,-3.2) node[align=center] {bibliothèque};
\draw[fill=poporangeL, draw=poporange] (8,-3.2) node[align=center] {bookshop};
\draw[fill=popblueL, draw=popblue] (0,-4.8) node[align=center] {sensible};
\draw[fill=popgreenL, draw=popgreen] (4,-4.8) node[align=center] {raisonnable};
\draw[fill=poporangeL, draw=poporange] (8,-4.8) node[align=center] {sensitive};
\draw[fill=popblueL, draw=popblue] (0,-6.4) node[align=center] {console (n.)};
\draw[fill=popgreenL, draw=popgreen] (4,-6.4) node[align=center] {console de jeux};
\draw[fill=poporangeL, draw=poporange] (8,-6.4) node[align=center] {to console (v.)};
\node[font=\bfseries\small, text=popdark] at (0,1) {Faux ami};
\node[font=\bfseries\small, text=popdark] at (4,1) {Sens anglais réel};
\node[font=\bfseries\small, text=popdark] at (8,1) {Pour le sens français};
\end{tikzpicture}
\legende{Tableau des faux amis : mot anglais piégeux, son vrai sens, et le mot correct pour exprimer le sens français.}
\end{popfigure}

\courssub{French calques to banish}

Au-delà des mots isolés, le francophone calque aussi des \textbf{constructions entières}. Ces calques sont immédiatement repérés par un correcteur.

\begin{propriete}[Les prépositions et les verbes ne se calquent pas du français]
Dites \eng{on the Internet} (jamais \emph{in Internet}), \eng{on social media}, \eng{at home}, \eng{addicted to video games}, \eng{hooked on a series}. Dites \eng{to take a photo} (jamais \emph{make a photo}), \eng{to watch TV} (jamais \emph{look at TV}), \eng{to listen to music} (la préposition \eng{to} est obligatoire), \eng{to play a game} (jamais \emph{play with a game} pour « jouer à un jeu »).
\end{propriete}

\begin{center}
\begin{tabularx}{\linewidth}{|X|X|X|}
\hline
\rowcolor{popblueL} Calque français (à bannir) & Anglais correct & Traduction \\
\hline
make a photo & take a photo & prendre une photo \\
\hline
in Internet & on the Internet & sur Internet \\
\hline
look at TV & watch TV & regarder la télévision \\
\hline
listen music & listen to music & écouter de la musique \\
\hline
I am addict to games & I am addicted to games & je suis accro aux jeux \\
\hline
assist a match & attend a match & assister à un match \\
\hline
\end{tabularx}
\end{center}

\begin{attention}[Watch, look at, see : ne confondez plus]
\eng{to watch} = regarder attentivement quelque chose qui bouge ou dure : \eng{watch TV}, \eng{watch a film}, \eng{watch a match}. \eng{to look at} = porter les yeux sur quelque chose de fixe : \eng{look at the screen}, \eng{look at this photo}. \eng{to see} = voir (percevoir) ou « voir » un film au cinéma : \eng{We saw a great film at the cinema.} On dit donc \eng{watch TV} mais \eng{look at your screen}.
\end{attention}

\begin{methode}[Comment piéger les faux amis avant qu'ils ne vous piègent]
\begin{enumerate}
\item Repérez tout mot anglais qui ressemble à un mot français (\eng{console}, \eng{assist}, \eng{sensible}...) : méfiance automatique.
\item Vérifiez son sens réel dans un dictionnaire (papier ou en ligne) avant de l'employer.
\item Apprenez toujours le mot \textbf{dans un exemple complet}, jamais isolément : \eng{I attended a gaming tournament} et non simplement \eng{attend = assister à}.
\item Construisez-vous une liste personnelle de faux amis, en deux colonnes : sens anglais réel / mot anglais pour le sens français.
\item Relisez vos productions écrites en traquant spécifiquement ces mots : c'est le réflexe du bon candidat.
\end{enumerate}
\end{methode}

\begin{exoresolu}[Correct the mistakes]
Énoncé — Corrigez les phrases suivantes, toutes fautives : 1. \emph{I assisted a concert in Bamenda last night.} 2. \emph{Actually, more and more teenagers use tablets for recreation.} 3. \emph{She made a photo of her new console.} 4. \emph{We found the game in a library in town.} 5. \emph{He is very sensible: he cries when his team loses online.}
\tcblower
Solution détaillée :
\begin{enumerate}
\item \eng{I \textbf{attended} a concert in Bamenda last night.} — « Assister à » (être présent) = \eng{to attend} ; \eng{to assist} = aider.
\item \eng{\textbf{Currently}, more and more teenagers use tablets for recreation.} — « Actuellement » = \eng{currently} ; \eng{actually} = en fait.
\item \eng{She \textbf{took} a photo of her new console.} — « Prendre une photo » = \eng{to take a photo} ; prétérit irrégulier : \eng{take, took, taken}.
\item \eng{We found the game in a \textbf{bookshop} in town.} — On achète dans une \eng{bookshop} (librairie, boutique) ; une \eng{library} est une bibliothèque où l'on emprunte.
\item \eng{He is very \textbf{sensitive}: he cries when his team loses online.} — « Sensible » (émotif) = \eng{sensitive} ; \eng{sensible} = raisonnable.
\end{enumerate}
\end{exoresolu}

\begin{exercice}[À vous de jouer]
1. Retrouvez et corrigez les huit erreurs du texte d'introduction (\emph{A weekend of digital recreation}).
2. Traduisez en anglais : a) « J'assiste actuellement à un tournoi de jeux vidéo à Douala. » b) « Mon petit frère est raisonnable : il ne joue qu'une heure par jour sur sa console. » c) « Elle a pris une photo et l'a publiée sur Internet. »
3. Complétez avec \eng{attend}, \eng{assist}, \eng{actually} ou \eng{currently} : a) I \ldots\ my grandmother with her smartphone every Sunday. b) \ldots , I don't like online games at all. c) Did you \ldots\ the film screening at the cultural centre?
\end{exercice}

\begin{corrige}[Exercice : À vous de jouer]
1. \eng{assisted} $\rightarrow$ \eng{attended} ; \eng{Actually} $\rightarrow$ \eng{In fact} (ou suppression) ; \eng{sensible} $\rightarrow$ \eng{not sensible} (le contexte exige l'inverse : jouer cinq heures d'affilée n'est pas raisonnable) ; \eng{library} $\rightarrow$ \eng{bookshop} ; \eng{consoled} $\rightarrow$ \eng{gave} (ou \eng{offered}) ; \eng{made a photo} $\rightarrow$ \eng{took a photo} ; \eng{in Internet} $\rightarrow$ \eng{on the Internet} ; \eng{looked at TV} $\rightarrow$ \eng{watched TV}.
2. a) \eng{I am currently attending a video game tournament in Douala.} b) \eng{My little brother is sensible: he only plays on his console for one hour a day.} c) \eng{She took a photo and posted it on the Internet.}
3. a) \eng{assist} (j'aide ma grand-mère) ; b) \eng{Actually} (en fait) ; c) \eng{attend} (as-tu assisté à la projection ?).
\end{corrige}

\begin{savaistu}[Les faux amis au Bac]
Au baccalauréat, l'emploi d'un faux ami grossier comme \eng{I assisted the match} ou \eng{actually I live in Yaoundé} coûte des points en compétence lexicale : les correcteurs y sont particulièrement attentifs car ces erreurs révèlent une pensée « traduite du français ». À l'inverse, l'emploi correct de \eng{attend}, \eng{currently} ou \eng{sensible} est immédiatement remarqué et valorisé. Les faux amis existent d'ailleurs dans les deux sens : un Anglophone qui apprend le français dira \emph{« je suis sensible »} en voulant dire \emph{I am sensible} (raisonnable) — la proximité entre les deux langues est un piège permanent !
\end{savaistu}

\begin{aretenir}[L'essentiel de la section]
\begin{itemize}
\item \eng{actually} = en fait ; « actuellement » = \eng{currently}, \eng{at present}.
\item \eng{to attend} = assister à ; \eng{to assist} = aider.
\item \eng{library} = bibliothèque ; « librairie » = \eng{bookshop}.
\item \eng{sensible} = raisonnable ; « sensible » = \eng{sensitive}.
\item \eng{game console} = console de jeux ; \eng{to console} = consoler.
\item Calques à bannir : \eng{take a photo}, \eng{on the Internet}, \eng{watch TV}, \eng{listen to music}, \eng{addicted to}.
\item Stratégie : vérifier au dictionnaire et apprendre chaque mot dans une phrase complète.
\end{itemize}
\end{aretenir}

\coursec{Word families, pronunciation and spelling}

Dans la section précédente, nous avons démasqué les faux amis et les calques du français (\eng{actually}, \eng{to assist}, \eng{a library}…). Restons dans la même logique de précision : pour bien employer le vocabulaire des loisirs numériques, il ne suffit pas de connaître un mot isolé. Il faut connaître toute sa \cle{famille} — verbe, nom, adjectifs —, savoir le \cle{prononcer} correctement et l'\cle{orthographier} sans faute. C'est l'objet de cette section.

\courssub{Observing word families in context}

\begin{textebox}[A weekend of digital leisure in Yaoundé]
Last Saturday, the Ngo Bassa family stayed at home and enjoyed a very \textbf{relaxing} weekend. Mr Ngo Bassa, who is usually stressed by his job, felt completely \textbf{relaxed} for the first time in weeks. His wife watched a comedy show on a \textbf{streaming} platform; she found it extremely \textbf{entertaining} and said she had not been so well \textbf{entertained} for a long time. Their son Junior spent the afternoon playing \textbf{video games} \textbf{online}. His mother worries that these games are \textbf{addictive}, but Junior insists he is not \textbf{addicted}: “Gaming is just my favourite form of \textbf{entertainment}, Mum. It helps me to \textbf{relax}.” Their daughter Sandrine prefers to \textbf{download} \textbf{podcasts} and listen to them offline. The whole family agrees on one thing: a good \textbf{Wi-Fi} \textbf{connection} has become essential for their \textbf{leisure} time. When the \textbf{connectivity} is poor, everyone complains! On Sunday evening, they all watched an \textbf{amusing} film together, and the children were so \textbf{amused} that they laughed until midnight.
\end{textebox}

\textbf{Glossaire :} \emph{stressed} (adj., stressé) ; \emph{to worry} (v., s'inquiéter) ; \emph{to insist} (v., insister) ; \emph{poor} (adj., faible, mauvais) ; \emph{to complain} (v., se plaindre) ; \emph{until} (prép., jusqu'à).

\textbf{Questions d'observation :}
\begin{enumerate}
\item Trouvez dans le texte trois adjectifs en \eng{-ing} et trois adjectifs en \eng{-ed}. Qui ou quoi décrivent-ils ?
\item Relevez tous les mots construits sur la racine \eng{entertain}. Quelle est la nature grammaticale de chacun ?
\item Comment sont orthographiés \eng{Wi-Fi}, \eng{online} et \eng{video games} ?
\end{enumerate}

\courssub{The main word families of digital leisure}

Une \cle{famille de mots} (\eng{word family}) regroupe tous les mots construits sur une même racine. En anglais, on passe du verbe au nom ou à l'adjectif grâce à des \cle{suffixes} : \eng{-ment}, \eng{-tion}, \eng{-ing}, \eng{-ed}, \eng{-ive}, \eng{-ity}… Maîtriser ces familles multiplie votre vocabulaire sans effort et enrichit vos copies.

\begin{definition}[Les cinq familles à connaître]
\begin{itemize}
\item \eng{to entertain} (v., divertir) → \eng{entertainment} (n., divertissement) → \eng{entertaining} (adj., divertissant) / \eng{entertained} (adj., diverti).
\item \eng{to relax} (v., se détendre) → \eng{relaxation} (n., détente) → \eng{relaxing} (adj., relaxant) / \eng{relaxed} (adj., détendu).
\item \eng{to amuse} (v., amuser) → \eng{amusement} (n., amusement) → \eng{amusing} (adj., amusant) / \eng{amused} (adj., amusé).
\item \eng{addict} (n., accro) → \eng{addiction} (n., dépendance) → \eng{addictive} (adj., qui rend accro) / \eng{addicted} (adj., accro, dépendant).
\item \eng{to connect} (v., connecter) → \eng{connection} (n., connexion) → \eng{connected} (adj., connecté) / \eng{connectivity} (n., connectivité).
\end{itemize}
\end{definition}

\begin{exemplebox}[Exemples traduits]
\eng{Video games are a popular form of entertainment.} « Les jeux vidéo sont une forme de divertissement populaire. »

\eng{This game is very addictive.} « Ce jeu rend très accro. »

\eng{Many teenagers are addicted to social media.} « Beaucoup d'adolescents sont accros aux réseaux sociaux. »

\eng{Listening to music is my favourite relaxation.} « Écouter de la musique est ma détente préférée. »

\eng{The connection in our neighbourhood is very slow.} « La connexion dans notre quartier est très lente. »

\eng{The comedian entertained the whole audience.} « L'humoriste a diverti tout le public. »
\end{exemplebox}

\begin{propriete}[La distinction -ing / -ed : qualité de la chose ou sentiment de la personne]
C'est LA règle capitale de ces familles :
\begin{itemize}
\item L'adjectif en \cle{-ing} décrit la \textbf{chose} ou la \textbf{situation} qui provoque le sentiment : \eng{The film is entertaining.} « Le film est divertissant. » (c'est le film qui divertit).
\item L'adjectif en \cle{-ed} décrit le \textbf{sentiment ressenti par la personne} : \eng{I am entertained.} « Je suis diverti. »
\end{itemize}
Même logique pour toutes les paires : \eng{relaxing}/\eng{relaxed}, \eng{amusing}/\eng{amused}, \eng{addictive}/\eng{addicted}, \eng{boring}/\eng{bored}, \eng{interesting}/\eng{interested}.

\eng{This podcast is boring, so I am bored.} « Ce podcast est ennuyeux, donc je m'ennuie. »
\end{propriete}

\begin{attention}[La préposition qui accompagne \eng{addicted}]
On dit toujours \eng{addicted \textbf{to} + nom ou -ing} : \eng{He is addicted to video games.} « Il est accro aux jeux vidéo. » ; \eng{She is addicted to scrolling.} « Elle est accro au défilement (des réseaux). » Ne dites jamais « addicted of » ni « addicted with ».
\end{attention}

\begin{center}
\begin{tabularx}{\linewidth}{|l|X|X|X|}\hline
\rowcolor{popblueL} \textbf{Racine} & \textbf{Verbe} & \textbf{Nom} & \textbf{Adjectifs} \\ \hline
entertain & to entertain (divertir) & entertainment (divertissement) & entertaining / entertained \\ \hline
relax & to relax (se détendre) & relaxation (détente) & relaxing / relaxed \\ \hline
amuse & to amuse (amuser) & amusement (amusement) & amusing / amused \\ \hline
addict & to addict (rendre accro, rare) & addiction (dépendance) & addictive / addicted \\ \hline
connect & to connect (connecter) & connection (connexion) & connected ; connectivity (n.) \\ \hline
\end{tabularx}
\end{center}

\begin{popfigure}
\begin{tikzpicture}[
  grow=down,
  level distance=1.4cm,
  level 1/.style={sibling distance=4.2cm},
  level 2/.style={sibling distance=2.6cm},
  every node/.style={draw, rounded corners, align=center, font=\small, fill=popcream},
  edge from parent/.style={draw=popblue, thick}
]
\node[fill=popblue, text=white, font=\bfseries] {entertain}
  child { node[fill=popblueL, align=center]{verb\\ \emph{to entertain}}
    child { node[align=center]{The singer entertains\\ the crowd.} } }
  child { node[fill=poporangeL, align=center]{noun\\ \emph{entertainment}}
    child { node[align=center]{Video games are\\ cheap entertainment.} } }
  child { node[fill=popgreenL, align=center]{adjectives\\ \emph{entertaining / entertained}}
    child { node[align=center]{an entertaining show\\ I feel entertained} } };
\end{tikzpicture}
\legende{Arbre lexical de la famille \eng{entertain} : une racine, trois natures, un exemple par branche.}
\end{popfigure}

\begin{methode}[Enrichir une copie : transformer un verbe en nom ou en adjectif]
Pour varier votre style à l'écrit, remplacez une structure verbale simple par un nom ou un adjectif de la même famille.
\begin{enumerate}
\item Repérez le verbe de votre phrase : \eng{I relax by gaming.}
\item Identifiez le nom correspondant : \eng{relaxation}.
\item Reconstruisez la phrase autour du nom : \eng{Gaming is my favourite relaxation.}
\end{enumerate}
Autres transformations : \eng{The show entertained us.} → \eng{The show was very entertaining.} / \eng{He is addicted to his phone.} → \eng{He has a phone addiction.} Cette technique, appelée \cle{nominalisation}, donne un style plus riche et plus scolaire — très appréciée dans les épreuves de composition du Baccalauréat.
\end{methode}

\courssub{Pronunciation of key words}

L'anglais ne se prononce pas comme il s'écrit. Voici la prononciation des mots essentiels du thème, notée en lettres françaises.

\begin{center}
\begin{tabularx}{\linewidth}{|l|X|X|}\hline
\rowcolor{popblueL} \textbf{Mot} & \textbf{Prononciation} & \textbf{Piège à éviter} \\ \hline
leisure & « lé-djeur » & ne pas dire « lé-sure » ni « lé-zeur » \\ \hline
download & « daoun-loud » & accent sur la première syllabe \\ \hline
upload & « ap-loud » & ne pas dire « oup-loud » \\ \hline
website & « ouèb-saït » & \eng{web} = « ouèb », pas « wèb » \\ \hline
podcast & « pod-caste » & \eng{-cast} comme dans \eng{broadcast} \\ \hline
streaming & « stri-ming » & \eng{ea} se prononce « i » long \\ \hline
Wi-Fi & « ouaï-faï » & les deux syllabes riment \\ \hline
\end{tabularx}
\end{center}

\begin{attention}[Prononciation : les erreurs typiques des francophones]
\begin{itemize}
\item \eng{leisure} se prononce « lé-djeur » et non « lé-sure » : le \eng{s} y sonne « dj ».
\item Dans \eng{download} et \eng{upload}, \eng{-load} se prononce « loud » (comme \eng{loud}, « fort »), pas « loude ».
\item \eng{website} commence par « ouèb » : le \eng{w} anglais est un « ou » bref, jamais un « v ».
\item Ne prononcez pas le \eng{e} final muet : \eng{game} = « guèim » en une seule syllabe.
\end{itemize}
\end{attention}

\courssub{Spelling rules for digital vocabulary}

\begin{propriete}[Orthographe : un mot ou deux ? majuscules ? trait d'union ?]
\begin{itemize}
\item \eng{Wi-Fi} : deux majuscules et un trait d'union (nom déposé, souvent expliqué comme \emph{Wireless Fidelity}).
\item \eng{website} : un seul mot, tout en minuscules.
\item \eng{online} et \eng{offline} : un seul mot, sans trait d'union.
\item \eng{e-mail} ou \eng{email} : les deux graphies sont correctes ; la forme sans trait d'union est la plus moderne.
\item \eng{video game} : deux mots séparés.
\end{itemize}
\end{propriete}

\begin{attention}[Orthographe : les fautes qui coûtent des points]
\begin{itemize}
\item On écrit \eng{online} en un seul mot — jamais « on line » ni « on-line ».
\item \eng{Wi-Fi} prend deux majuscules — jamais « wifi » ni « Wifi » dans un devoir soigné.
\item \eng{website} est un seul mot — jamais « web site ».
\item \eng{video game} s'écrit en deux mots — jamais « videogame ».
\item Attention au nom : on écrit \eng{entertainment} avec \eng{-tain-} au milieu (comme \eng{entertain}), pas « entertaiment ».
\end{itemize}
\end{attention}

\begin{savaistu}[D'où vient le mot \eng{entertainment} ?]
Le suffixe \eng{-tainment} de \eng{entertainment} remonte au latin \emph{tenere}, qui signifie « tenir ». Le divertissement, c'est littéralement ce qui « retient » votre attention entre (\emph{inter-}) les mains ! La même racine se retrouve dans \eng{to maintain} (maintenir), \eng{to contain} (contenir)… et dans le français \emph{entretenir}. La prochaine fois qu'une série « retient » toute votre famille devant l'écran, vous saurez que l'étymologie avait tout prévu.
\end{savaistu}

\begin{exoresolu}[Word families in action]
Énoncé — Complétez chaque phrase avec la forme correcte du mot entre parenthèses.

1. Watching comedies is very \ldots (relax). \quad 2. The children were \ldots by the magician's tricks (amuse). \quad 3. Gaming is his favourite form of \ldots (entertain). \quad 4. This app is so \ldots that I can't stop using it (addict). \quad 5. We need a better internet \ldots in our neighbourhood (connect).

\tcblower
Solution détaillée :

1. \eng{relaxing} : on décrit la \emph{chose} (watching comedies) qui provoque la détente → adjectif en \eng{-ing}. « Regarder des comédies est très relaxant. »

2. \eng{amused} : on décrit le \emph{sentiment des enfants} → adjectif en \eng{-ed}. « Les enfants étaient amusés par les tours du magicien. »

3. \eng{entertainment} : après la préposition \eng{of}, il faut un \emph{nom}. « Le jeu vidéo est sa forme de divertissement préférée. »

4. \eng{addictive} : l'application \emph{provoque} la dépendance → \eng{-ive}. « Cette application rend tellement accro que je ne peux pas arrêter de l'utiliser. »

5. \eng{connection} : après l'adjectif \eng{better} et l'article, il faut un \emph{nom}. « Nous avons besoin d'une meilleure connexion internet dans notre quartier. »
\end{exoresolu}

\begin{exercice}[À vous de jouer]
\begin{enumerate}
\item Transformez pour varier le style : \eng{She relaxes by listening to podcasts.} → utilisez le nom \eng{relaxation}.
\item Choisissez la bonne forme : \eng{The film was (entertaining / entertained), so everybody felt (entertaining / entertained).}
\item Corrigez l'orthographe : « I play videogames on line when the wifi is good. »
\item Prononcez à voix haute : \eng{leisure, download, website, streaming}, puis faites-vous corriger par un camarade.
\end{enumerate}
\end{exercice}

\begin{corrige}[Exercice « À vous de jouer »]
\begin{enumerate}
\item \eng{Listening to podcasts is her favourite relaxation.} (nominalisation du verbe \eng{relax}).
\item \eng{The film was entertaining, so everybody felt entertained.} (le film = qualité → \eng{-ing} ; les spectateurs = sentiment → \eng{-ed}).
\item \eng{I play video games online when the Wi-Fi is good.} (\eng{video game} en deux mots, \eng{online} en un seul mot, \eng{Wi-Fi} avec deux majuscules et un trait d'union).
\item « lé-djeur », « daoun-loud », « ouèb-saït », « stri-ming ».
\end{enumerate}
\end{corrige}

\begin{aretenir}[L'essentiel de la section]
\begin{itemize}
\item Chaque verbe du thème forme une famille : \eng{entertain → entertainment → entertaining / entertained}.
\item \eng{-ing} = qualité de la chose ; \eng{-ed} = sentiment de la personne.
\item \eng{addicted} se construit toujours avec la préposition \eng{to}.
\item Prononciation : \eng{leisure} « lé-djeur », \eng{website} « ouèb-saït ».
\item Orthographe : \eng{Wi-Fi}, \eng{website}, \eng{online}/\eng{offline} en un mot, \eng{video game} en deux mots.
\item Pour enrichir une copie, transformez le verbe en nom : \eng{I relax} → \eng{my favourite relaxation}.
\end{itemize}
\end{aretenir}

\coursec{Using the vocabulary in context}

Dans les sections précédentes, nous avons travaillé sur les \cle{familles de mots} (\eng{entertain, entertainment, entertaining}), la \cle{prononciation} et l'\cle{orthographe} du lexique des loisirs numériques. Mais un mot n'est vraiment « su » que lorsqu'on sait l'employer dans une phrase, un dialogue ou un paragraphe personnel. Cette dernière section est donc consacrée au \cle{réemploi en contexte} : lire, réutiliser, produire. C'est exactement le parcours exigé au baccalauréat, où le vocabulaire est évalué d'abord en compréhension (paper 1), puis en production écrite (paper 2).

\begin{popfigure}
\begin{center}
\begin{tikzpicture}[node distance=1.1cm]
\node[draw=popblue, fill=popblueL, rounded corners, thick, minimum width=3.4cm, minimum height=1.4cm, align=center] (u) {\textbf{1. Understand}\\ \eng{read and listen}};
\node[draw=poporange, fill=poporangeL, rounded corners, thick, minimum width=3.4cm, minimum height=1.4cm, align=center, right=of u] (r) {\textbf{2. Reuse}\\ \eng{sentences, dialogue}};
\node[draw=popgreen, fill=popgreenL, rounded corners, thick, minimum width=3.4cm, minimum height=1.4cm, align=center, right=of r] (p) {\textbf{3. Produce}\\ \eng{your own paragraph}};
\draw[-{Stealth}, very thick, popdark] (u) -- (r);
\draw[-{Stealth}, very thick, popdark] (r) -- (p);
\end{tikzpicture}
\end{center}
\legende{Les trois étapes du réemploi du lexique : comprendre, réutiliser, produire.}
\end{popfigure}

\courssub{Guided reading: spotting the vocabulary in an article}

\begin{textebox}[Digital leisure among young Cameroonians]
In many Cameroonian towns, young people relax with ICTs in different ways. At home, those who can afford Wi-Fi stream films and chat on social media. Others walk to the nearest cybercafé or gaming centre, where they pay a small fee to browse the Internet or play football video games with friends. For most teenagers, a weekly data bundle is a treasure: it means music, jokes, and video calls with cousins abroad.

In Douala and Yaoundé, Saturday afternoons have a new rhythm. Groups of boys gather around a single screen to follow the English Premier League, while their sisters download the latest Nigerian and Cameroonian hits onto their phones. A funny video can go viral in a few hours: by Monday morning, the whole school is laughing about it. Some students even binge-watch entire series during the holidays, finishing twelve episodes in one weekend.

Parents and teachers, however, are not always pleased. They worry about screen time: a teenager who spends six hours a day online has little time left for homework, sport or sleep. The cost of data is another problem. When the bundle is finished, some young people skip lunch to buy a new one. A few become truly addicted: they check their phones every two minutes and feel nervous without them.

Yet ICTs also open doors. Through video calls, families stay close to relatives in Europe or America. Creative teenagers record songs, edit short films and post them online, sometimes earning their first money. The challenge, therefore, is not to reject digital leisure, but to use it wisely: enjoy the screen, but do not let the screen control you.
\end{textebox}

\textbf{Glossaire.}

\begin{center}
\begin{tabularx}{\linewidth}{|l|l|X|}
\hline
\rowcolor{popblueL} \textbf{Mot anglais} & \textbf{Nature} & \textbf{Traduction} \\
\hline
to afford & verbe & avoir les moyens de \\
\hline
a fee & nom & un droit, un tarif (d'entrée ou d'utilisation) \\
\hline
a data bundle & nom & un forfait internet \\
\hline
to go viral & expression & devenir viral, se propager très vite \\
\hline
to binge-watch & verbe & regarder (une série) en boucle, d'une traite \\
\hline
screen time & nom & le temps passé devant un écran \\
\hline
to skip lunch & expression & sauter le déjeuner \\
\hline
addicted & adjectif & accro, dépendant \\
\hline
wisely & adverbe & sagement, avec discernement \\
\hline
\end{tabularx}
\end{center}

\textbf{Comprehension questions.}

\begin{enumerate}
\item Where do young people without Wi-Fi at home go to use the Internet?
\item What does a weekly data bundle mean for most teenagers?
\item Give two reasons why parents worry about digital leisure.
\item How can ICTs “open doors” for creative teenagers?
\item Find in the text one collocation with \eng{screen} and one idiom about popularity online.
\end{enumerate}

\begin{corrige}[Questions de compréhension]
\begin{enumerate}
\item They go to the nearest cybercafé or gaming centre, where they pay a small fee.
\item It means music, jokes, and video calls with cousins abroad.
\item Parents worry about excessive screen time (less time for homework, sport and sleep) and about the cost of data (some teenagers skip lunch to buy a bundle).
\item They can record songs, edit short films, post them online and sometimes earn their first money.
\item Collocation: \eng{screen time}. Idiom: \eng{to go viral}.
\end{enumerate}
\end{corrige}

\courssub{Gap-fill practice: reusing collocations and verbs}

\begin{exoresolu}[Complete the sentences]
Complétez avec : \eng{stream} — \eng{binge-watch} — \eng{data bundle} — \eng{go viral} — \eng{screen time} — \eng{hooked on}.

Énoncé :
\begin{enumerate}
\item My little brother can \ldots\ a whole series in two days.
\item I bought a weekly \ldots\ of 2 GB to listen to music.
\item Her dance video will certainly \ldots\ ; everyone is sharing it.
\item We cannot \ldots\ films at home because the Wi-Fi is too slow.
\item Doctors advise teenagers to reduce their \ldots\ before sleeping.
\item Many students are \ldots\ social media and forget their homework.
\end{enumerate}
\tcblower
Solution détaillée :
\begin{enumerate}
\item \eng{binge-watch} : le complément \eng{a whole series in two days} appelle l'idée de « tout regarder d'une traite ».
\item \eng{data bundle} : \eng{2 GB} indique un forfait internet ; attention à l'ordre des mots (\eng{data bundle}, jamais \eng{bundle data}).
\item \eng{go viral} : \eng{everyone is sharing it} = propagation rapide ; verbe irrégulier \eng{go, went, gone}.
\item \eng{stream} : on « streame » un film ; impossible avec un Wi-Fi lent.
\item \eng{screen time} : collocation nominale invariable (pas de pluriel).
\item \eng{hooked on} : expression idiomatique « être accro à », suivie de la préposition \eng{on} et d'un nom ou d'un gérondif (\eng{hooked on playing}).
\end{enumerate}
\end{exoresolu}

\begin{attention}[Pièges des francophones]
\begin{itemize}
\item \eng{hooked on} exige la préposition \eng{on} : \eng{I am hooked on music apps.} Jamais \eng{hooked to} ni \eng{hooked with}. En revanche, \eng{addicted} prend \eng{to} : \eng{He is addicted to online games.}
\item \eng{data} est ici invariable dans la collocation \eng{a data bundle} ; ne dites pas \eng{a datas bundle}.
\item \eng{to go viral} se conjugue comme \eng{go} : \eng{The video went viral yesterday.} (« La vidéo est devenue virale hier. »)
\item Ne traduisez pas « forfait » par \eng{package} tout seul : dites \eng{a data bundle} ou \eng{a data plan}.
\item « Passer du temps devant l'écran » se dit \eng{to spend time in front of a screen} ou \eng{to spend time online}, et non \eng{to pass time} (calque du français).
\end{itemize}
\end{attention}

\courssub{Model dialogue: at home vs at the cybercafé}

\begin{textebox}[A dialogue between two friends, Amina and Roland]
Amina: Hi Roland! What did you do last weekend?

Roland: I spent Saturday at the gaming centre near the market. We played football video games against other teams online. And you?

Amina: I stayed at home. We have Wi-Fi now, so I streamed a film and chatted with my cousin in Bamenda on a video call.

Roland: Lucky you! At the cybercafé, I pay a fee for one hour, and the connection is sometimes slow.

Amina: I know. But be careful: last month you bought three data bundles in one week!

Roland: You are right. My screen time is too high. I am a bit hooked on online games.

Amina: Then come and watch the match at my place on Saturday. It is cheaper, and we can share the Wi-Fi.

Roland: Great idea! I will bring the controller, and we will play together instead of paying at the gaming centre.
\end{textebox}

\begin{exemplebox}[Exemples traduits]
\eng{I am hooked on music apps, but I control my screen time.} « Je suis accro aux applications musicales, mais je contrôle mon temps d'écran. »

\eng{At the cybercafé, we play against other teams online.} « Au cybercafé, nous affrontons d'autres équipes en ligne. »

\eng{Her video went viral in one night.} « Sa vidéo est devenue virale en une nuit. »

\eng{I stream music because CDs are expensive here.} « J'écoute de la musique en streaming parce que les CD sont chers ici. »
\end{exemplebox}

\begin{attention}[Variez les verbes dans un dialogue]
Dans un dialogue ou une rédaction, évitez de répéter \eng{play} à chaque phrase. Montrez au correcteur l'étendue de votre lexique en variant : \eng{stream} (un film), \eng{chat} (avec un ami), \eng{browse} (le Web), \eng{watch} (un match), \eng{download} (une chanson), \eng{post} (une photo), \eng{share} (une vidéo). Une bonne copie d'examen contient au moins quatre verbes différents du champ lexical.
\end{attention}

\courssub{Guided paragraph: “My favourite digital leisure activity”}

\begin{methode}[Rédiger un paragraphe de 5 à 6 lignes]
\begin{enumerate}
\item \textbf{Name the activity} : nommez l'activité (\eng{My favourite digital leisure activity is streaming music.}).
\item \textbf{Say where and with whom} : dites où et avec qui (\eng{I usually listen at home, with my sister, after school.}).
\item \textbf{Use two collocations and one idiom} : employez deux collocations et une expression idiomatique (\eng{data bundle}, \eng{screen time}, \eng{to be hooked on}, \eng{to go viral}).
\item \textbf{Conclude with an advantage or a danger} : concluez par un avantage ou un danger (\eng{It relaxes me, but I limit my screen time before exams.}).
\end{enumerate}
\end{methode}

\begin{exemplebox}[Modèle de paragraphe]
\eng{My favourite digital leisure activity is watching comedy videos on my phone. Every evening, at home, my brother and I browse the Internet with our weekly data bundle. We are hooked on a Cameroonian comedian whose sketches often go viral. This activity makes us laugh and forget our problems. However, I control my screen time, because too much phone before sleeping is bad for my eyes.}

Traduction : « Mon loisir numérique préféré est de regarder des vidéos comiques sur mon téléphone. Chaque soir, à la maison, mon frère et moi naviguons sur Internet avec notre forfait hebdomadaire. Nous sommes accros à un humoriste camerounais dont les sketchs deviennent souvent viraux. Cette activité nous fait rire et oublier nos problèmes. Cependant, je contrôle mon temps d'écran, car trop de téléphone avant de dormir est mauvais pour les yeux. »
\end{exemplebox}

\courssub{Pair discussion: advantages and dangers of digital leisure}

\begin{experience}[Discussion en binômes]
Par groupes de deux, discutez pendant cinq minutes des avantages et des dangers des loisirs numériques. Chaque élève doit employer au moins cinq mots du thème. Banque d'idées :

\begin{itemize}
\item \textbf{Advantages} : \eng{relaxing, cheap entertainment, video calls with family abroad, learning English through films, going viral can bring fame and money}.
\item \textbf{Dangers} : \eng{too much screen time, addiction, the cost of data bundles, less time for homework and sport, skipping meals to buy data}.
\end{itemize}

Phrases utiles : \eng{In my opinion\ldots} / \eng{I agree, but\ldots} / \eng{On the one hand\ldots on the other hand\ldots} / \eng{The main danger is\ldots}
\end{experience}

\begin{savaistu}[Ce lexique au baccalauréat]
Au Bac (paper 2, essay), un sujet fréquent porte sur \eng{the advantages and disadvantages of the Internet or social media for young people}. Tout le lexique de cette leçon — \eng{screen time, data bundle, go viral, addicted, stream, cybercafé} — y est directement réemployable. Un paragraphe qui contient cinq mots précis du champ lexical vaut toujours mieux qu'un paragraphe vague avec \eng{good} et \eng{bad}.
\end{savaistu}

\courssub{Self-assessment: my vocabulary checklist}

\begin{exercice}[Checklist d'auto-évaluation]
Pour chaque mot, cochez ce que vous savez faire : le définir (D), le prononcer (P), l'employer dans une phrase (E).
\end{exercice}

\begin{center}
\begin{tabularx}{\linewidth}{|X|c|c|c|}
\hline
\rowcolor{popblueL} \textbf{Word / expression} & \textbf{D} & \textbf{P} & \textbf{E} \\
\hline
to stream & & & \\
\hline
to binge-watch & & & \\
\hline
a data bundle & & & \\
\hline
to go viral & & & \\
\hline
screen time & & & \\
\hline
to be hooked on & & & \\
\hline
a cybercafé / a gaming centre & & & \\
\hline
to browse the Internet & & & \\
\hline
addicted / addiction & & & \\
\hline
a video call & & & \\
\hline
\end{tabularx}
\end{center}

\begin{corrige}[Auto-évaluation]
Il n'y a pas de correction unique : c'est un outil personnel. Règle du professeur : si une case « E » (employer en phrase) n'est pas cochée, écrivez immédiatement une phrase avec ce mot, par exemple \eng{My uncle makes a video call to Douala every Sunday.} (« Mon oncle fait un appel vidéo vers Douala chaque dimanche. ») Un mot que l'on sait définir mais pas employer n'est qu'à moitié appris.
\end{corrige}

\begin{aretenir}[L'essentiel de la section]
\begin{itemize}
\item Le lexique se maîtrise en trois étapes : \textbf{comprendre} (lire un texte), \textbf{réutiliser} (phrases à trous, dialogue), \textbf{produire} (paragraphe personnel).
\item En production écrite ou orale : variez les verbes (\eng{stream, chat, browse, download, post, share}), employez des collocations exactes (\eng{screen time, data bundle}) et au moins une expression idiomatique (\eng{go viral, be hooked on}).
\item Attention aux prépositions : \eng{hooked on}, \eng{addicted to}, \eng{play against}, \eng{chat with}.
\item Ce vocabulaire est une réserve directe pour l'essai du baccalauréat sur Internet et les réseaux sociaux.
\end{itemize}
\end{aretenir}

\newpage

\coursec{Activité d'intégration}

\courssub{Objectifs de l'activité}

À l'issue de cette activité d'intégration, l'élève sera capable de :

\begin{itemize}
\item \cle{réemployer} à l'oral et à l'écrit le vocabulaire de la leçon relatif aux TIC comme ressources de loisir (\eng{ICTs as recreational resources}) : équipements, lieux, activités en ligne et hors ligne ;
\item utiliser correctement les \cle{collocations} étudiées (\eng{to browse the Internet}, \eng{to stream music}, \eng{to top up data}...) et au moins une expression idiomatique ;
\item éviter les \cle{faux amis} et les calques du français signalés dans la leçon (\eng{actually} $\neq$ « actuellement », \eng{a library} $\neq$ « une librairie »...) ;
\item produire un \cle{dialogue oral} structuré (émission de radio) avec un registre de langue adapté : accroche, interviews, conseils, conclusion ;
\item développer l'écoute active : repérer dans la production des camarades les mots du thème grâce à une grille d'écoute.
\end{itemize}

\courssub{Situation}

La radio communautaire de ta ville, \eng{Mount Cameroon Community Radio} (Buea), prépare une émission jeunesse hebdomadaire. Le thème du prochain numéro : les loisirs numériques des jeunes de la communauté. Le producteur recherche des groupes d'élèves de Terminale capables de tenir une émission de cinq minutes, en anglais, riche en vocabulaire précis. Ta classe a été sélectionnée : à vous de convaincre !

\begin{textebox}[Announcement — Mount Cameroon Community Radio, Buea]
\textbf{CALL FOR YOUNG RADIO TALENTS!}

Are you between 16 and 19? Do you love radio? \eng{MCCR Youth Hour} is looking for student teams to produce a five-minute show on:

“\textbf{Digital Leisure in Our Community}”

How do young people in Buea relax with technology? At home on a smartphone, or at the gaming centre down the street? Streaming music, chatting on social media, playing video games online — or just downloading films for the weekend? And what do parents and teachers think about screen time and the cost of data?

Your show must include: one presenter, two young guests with different habits, and one adult adviser (parent or teacher). Use rich, accurate English. The best team will be broadcast live on Saturday at 5 p.m.!

Send your script to the station or perform it live in your school hall.
\end{textebox}

\begin{definition}[Glossaire utile de l'annonce]
\begin{itemize}
\item \eng{to look for} (v.) : rechercher ;
\item \eng{a presenter} (n.) : un animateur, une animatrice ;
\item \eng{screen time} (n.) : le temps passé devant les écrans ;
\item \eng{data} (n.) : les données mobiles, le « forfait Internet » ;
\item \eng{to broadcast} (v.) : diffuser (verbe irrégulier : \eng{broadcast -- broadcast -- broadcast}) ;
\item \eng{a script} (n.) : un script, le texte écrit de l'émission.
\end{itemize}
\end{definition}

\courssub{Consignes}

Travail en groupes de quatre. L'activité se déroule en cinq étapes guidées.

\begin{enumerate}
\item \textbf{Compréhension et repérage (10 min).} Lisez l'annonce de la radio. Soulignez les mots du champ lexical des loisirs numériques qu'elle contient. Classez-les en trois colonnes : \eng{equipment and places}, \eng{activities}, \eng{people}. Ajoutez à chaque colonne au moins trois autres mots appris dans la leçon.

\item \textbf{Répartition des rôles et plan de l'émission (10 min).} Attribuez les rôles : le \cle{présentateur} (\eng{the presenter/host}), deux \cle{invités} (\eng{guests}) aux habitudes différentes — l'un préfère les loisirs numériques \emph{à la maison}, l'autre au \eng{cybercafé} ou au \eng{gaming centre} — et un \cle{conseiller adulte} (parent ou enseignant). Rédigez un plan en quatre temps : accroche et présentation des invités ; interview du premier invité ; interview du second ; conseils de l'adulte et conclusion.

\item \textbf{Rédaction du script (25 min).} Écrivez le script complet (environ 350 mots). Contraintes obligatoires :
\begin{itemize}
\item au moins \cle{10 mots du vocabulaire} de la leçon (soulignez-les dans votre script) ;
\item au moins \cle{3 collocations} correctes, par exemple \eng{to browse the Internet}, \eng{to stream music}, \eng{to download apps}, \eng{to top up my data}, \eng{to spend time online} ;
\item au moins \cle{1 expression idiomatique} : \eng{to be glued to one's screen} (« être rivé à son écran »), \eng{to kill time} (« tuer le temps »), \eng{in the blink of an eye} (« en un clin d'œil ») ;
\item \cle{aucun faux ami} : attention à \eng{actually} (= « en fait », pas « actuellement »), \eng{to assist} (= « aider quelqu'un », pas « assister à »), \eng{a library} (= « une bibliothèque », pas « une librairie »).
\end{itemize}

\item \textbf{Répétition et prononciation (10 min).} Répétez l'émission. Vérifiez la prononciation des mots pièges : \eng{technology} « tèk-no-lo-dji », \eng{leisure} « lè-jeur », \eng{data} « dé-i-te », \eng{wireless} « ouaïeur-lèss ». Le présentateur doit parler clairement et poser des questions fermées avec l'intonation montante (\eng{Do you play online games?}) et les questions ouvertes avec l'intonation descendante (\eng{What do you use your phone for?}).

\item \textbf{Performance et écoute active (15 min).} Chaque groupe joue son émission devant la classe (ou la fait enregistrer sur un téléphone). Pendant l'écoute, les autres élèves cochent sur leur grille les mots du thème entendus et relèvent les collocations et l'expression idiomatique. Vote final : le meilleur groupe est celui qui réemploie le lexique le plus riche et le plus correct.
\end{enumerate}

\begin{exemplebox}[Grille d'écoute — Listening grid]
\begin{center}
\begin{tabularx}{\linewidth}{|X|X|X|}\hline
\rowcolor{popblueL} \textbf{Category} & \textbf{Words I heard} & \textbf{Tick} \\ \hline
Equipment / places & smartphone, laptop, headset, cybercafé, gaming centre, Wi-Fi hotspot & \\ \hline
Activities & to stream, to download, to browse, to chat, to play online games, to upload & \\ \hline
Collocations & to top up data, to spend time online, to browse the Internet & \\ \hline
Idiom & to be glued to one's screen, to kill time & \\ \hline
\end{tabularx}
\end{center}
\end{exemplebox}

\courssub{Ressources à mobiliser}

\begin{aretenir}[Boîte à outils de l'émission]
\textbf{Lexique du thème :}
\begin{itemize}
\item Équipements et lieux : \eng{smartphone, laptop, tablet, games console, headset, keyboard, Wi-Fi connection, cybercafé, gaming centre, Internet café}.
\item Activités : \eng{to browse/surf the Internet, to stream music or films, to download/upload files, to chat on social media, to play online/video games, to watch videos, to post photos, to top up one's data}.
\item Personnes : \eng{a user, a gamer, a blogger, a subscriber, a viewer}.
\end{itemize}
\textbf{Structures utiles pour l'interview :}
\begin{itemize}
\item \eng{How often do you...?} / \eng{What do you use your phone for?}
\item \eng{I spend about two hours a day online.} (présent simple + expression de fréquence)
\item \eng{I prefer streaming music to downloading it.} (\eng{to prefer ... to ...} + gérondif)
\item Conseils : \eng{You should limit your screen time.} / \eng{Why don't you switch off your phone at night?}
\end{itemize}
\textbf{Prononciation :} accent tonique sur la première syllabe de \eng{CYbercafé} « saï-beur-ka-fé » et de \eng{LEIsure} « lè-jeur » ; on écrit \eng{the Internet} avec une majuscule et l'article défini.
\end{aretenir}

\begin{attention}[Pièges à éviter dans votre script]
\begin{itemize}
\item \eng{actually} signifie « en fait » ; pour « actuellement », dites \eng{currently} ou \eng{at the moment}.
\item On dit \eng{to play games} (jouer \emph{à} des jeux), jamais \eng{to play at games}.
\item \eng{to assist} = aider quelqu'un ; « assister à un match » = \eng{to attend a match} ou \eng{to watch a match}.
\item \eng{data} est ici invariable et singulier : \eng{data is expensive}, jamais \eng{datas}.
\item Pas de calque : « passer du temps » = \eng{to spend time} (et non \eng{to pass time}) ; « surfer sur Internet » = \eng{to surf the Internet} (avec l'article).
\end{itemize}
\end{attention}

\courssub{Proposition de production et corrigé commenté}

\begin{exoresolu}[Script modèle — “Digital Leisure in Our Community”]
\textbf{Énoncé.} Rédiger et jouer une émission de radio de cinq minutes respectant toutes les contraintes (10 mots du thème, 3 collocations, 1 expression idiomatique, aucun faux ami).

\tcblower

\textbf{Production modèle (script) :}

\eng{\textbf{Presenter:} Good afternoon, dear listeners, and welcome to MCCR Youth Hour, live from Buea! Today's topic is “Digital Leisure in Our Community”. How do young people here relax with technology? With me in the studio are two students, Miriam and Elvis, and their teacher, Mr Tabi. Miriam, let's start with you. What do you use your smartphone for?}

\eng{\textbf{Miriam:} Well, I do almost everything at home because we have a Wi-Fi connection. I browse the Internet for my homework, I stream music while I cook, and I chat with my friends on social media in the evening. At the weekend, I download films and watch them with my little brother.}

\eng{\textbf{Presenter:} How many hours a day do you spend online?}

\eng{\textbf{Miriam:} Honestly, about three hours. My mother says I am glued to my screen!}

\eng{\textbf{Presenter:} Elvis, your habits are quite different, aren't they?}

\eng{\textbf{Elvis:} Yes, completely! We have no Wi-Fi at home, so I go to the gaming centre near the market. I play online games there with my friends — football games mostly. I also top up my data twice a week, but data is expensive, so I use it carefully. Actually, the gaming centre is cheaper for me than buying a lot of data.}

\eng{\textbf{Presenter:} Mr Tabi, as a teacher, what is your advice?}

\eng{\textbf{Mr Tabi:} Technology is a wonderful recreational resource, but young people should limit their screen time. Why don't you switch off your phone one hour before sleeping? And remember: the Internet is also a library — use it to learn, not only to kill time.}

\eng{\textbf{Presenter:} Thank you all! That was MCCR Youth Hour. Stay connected... but not too connected!}

\medskip
\textbf{Commentaire des choix de langue :}
\begin{itemize}
\item \textbf{Lexique du thème (plus de 10 mots) :} \eng{smartphone, Wi-Fi connection, browse the Internet, stream music, social media, download films, glued to my screen, gaming centre, online games, top up my data, screen time, library}.
\item \textbf{Collocations respectées :} \eng{to browse the Internet}, \eng{to stream music}, \eng{to top up my data}, \eng{to spend hours online}, \eng{to limit screen time}.
\item \textbf{Expression idiomatique :} \eng{to be glued to one's screen} (« être rivé à son écran ») et, en bonus, \eng{to kill time} (« tuer le temps »).
\item \textbf{Faux amis évités :} \eng{actually} est bien employé au sens de « en fait » ; \eng{library} désigne ici une bibliothèque virtuelle, pas une librairie ; \eng{data} reste invariable.
\item \textbf{Grammaire :} présent simple pour les habitudes (\eng{I go, I play, I use}), fréquence avec \eng{twice a week}, question fermée à intonation montante (\eng{aren't they?}), conseils avec \eng{should} et \eng{Why don't you...?}.
\item \textbf{Registre radiophonique :} accroche (\eng{Good afternoon, dear listeners}), transitions (\eng{let's start with you}), chute humoristique finale (\eng{Stay connected... but not too connected!}).
\end{itemize}
\end{exoresolu}

\courssub{Bilan de l'activité}

Cette activité a permis de mobiliser l'ensemble des compétences de la leçon dans une tâche de communication authentique :

\begin{itemize}
\item \textbf{Lexique :} les élèves ont réemployé activement le vocabulaire des équipements, des lieux et des activités numériques de loisir, dans un contexte camerounais réaliste (cybercafé, coût des données, Wi-Fi à la maison).
\item \textbf{Collocations et idiomes :} la contrainte des trois collocations et de l'expression idiomatique a obligé chaque groupe à dépasser le simple mot isolé pour construire des combinaisons naturelles de l'anglais.
\item \textbf{Vigilance linguistique :} l'interdiction des faux amis a renforcé la conscience des écarts entre français et anglais (\eng{actually}, \eng{library}, \eng{data} invariable).
\item \textbf{Oral et écrit :} le script écrit a préparé la performance orale ; la grille d'écoute a transformé le reste de la classe en auditeurs actifs, développant la compréhension orale.
\item \textbf{Compétence de vie :} le débat sur le temps d'écran et le coût des données invite à un usage raisonné et responsable des TIC, à la maison comme dans la communauté.
\end{itemize}

\begin{experience}[Pour aller plus loin]
En devoir facultatif, chaque élève peut rédiger un court article de blog (120 mots) intitulé \eng{“My Digital Leisure Habits”} en réutilisant au moins huit mots du thème et une expression idiomatique différente de celle de son groupe. Les meilleurs articles seront affichés au tableau d'anglais de la classe.
\end{experience}

\newpage

\coursec{Méthodes et exercices résolus}

\coursec{Lesson 57 — Vocabulary: Words and expressions related to using ICTs as recreational resources at home or in the community}

\begin{textebox}[Blog post — Digital leisure in Yaoundé]
\eng{In Yaoundé, as in many African cities, the smartphone has become the primary gateway to entertainment. After school or work, young people gather around Wi-Fi hotspots in neighbourhoods like Mvog-Ada or Essos to stream music videos, watch football highlights on YouTube, or challenge each other on mobile games like PUBG Mobile and Free Fire. Those who can afford data bundles stream series on Netflix or Showmax at home, while others prefer the social atmosphere of an Internet café near the market, where a gaming session costs only a few hundred francs. WhatsApp and TikTok are the undisputed kings of social media: users spend hours scrolling through feeds, sharing comedy skits, and sending voice notes. However, parents and teachers often worry about excessive screen time and the risk of addiction. The challenge for the youth is to balance online fun with offline responsibilities — sports, reading, and family time.}
\end{textebox}

\courssub{ICT equipment and places for recreation}

\begin{methode}[Bien choisir le mot du champ lexical des TIC récréatives]
Pour employer correctement le vocabulaire des loisirs numériques, suis ces étapes :
\begin{enumerate}
\item \cle{Classe le mot} dans la bonne sous-catégorie : \emph{equipment} (\eng{smartphone, laptop, console, headset, router}), \emph{places} (\eng{Internet café, community telecentre, gaming lounge, Wi-Fi hotspot}), \emph{activities} (\eng{streaming, gaming, browsing, chatting, downloading}).
\item \cle{Vérifie la nature du mot} : nom, verbe ou adjectif. \eng{to stream} (verbe), \eng{a stream} (nom), \eng{streaming} (nom d'activité).
\item \cle{Choisis le bon verbe d'action} : on \eng{plays} a game, on \eng{watches} a video, on \eng{listens to} music, on \eng{surfs} the Internet, on \eng{downloads} a film, on \eng{posts} a comment.
\item \cle{Contrôle la préposition} : \eng{to listen \textbf{to}}, \eng{to chat \textbf{with} friends}, \eng{to be addicted \textbf{to}}, \eng{to log \textbf{on to}} a website.
\item \cle{Relis ta phrase} et demande-toi si un anglophone dirait exactement cela.
\end{enumerate}
Réflexe : apprends chaque mot avec son verbe et sa préposition habituels, jamais isolément.
\end{methode}

\begin{exoresolu}[Compléter avec le bon mot du champ lexical]
Énoncé — Complete the sentences with the correct word from the box: \emph{stream, headset, Internet café, download, surf}.\par
1. In the evening, many students in Yaoundé \ldots\ the Internet on their phones to relax.\par
2. My brother bought a new \ldots\ so that he can hear his teammates clearly while gaming.\par
3. It is cheaper to \ldots\ the film at night when data bundles are cheaper.\par
4. Since we have no Wi-Fi at home, I sometimes go to the \ldots\ near the market to play online games.\par
5. We \ldots\ music on a free platform instead of buying CDs.
\tcblower
Solution détaillée :\par
1. \eng{surf}. Le verbe idiomatique pour « naviguer sur Internet » est \eng{to surf the Internet}. Sujet pluriel, présent simple : \eng{surf}. « Many students in Yaoundé surf the Internet on their phones to relax. »\par
2. \eng{headset}. L'indice est \eng{hear his teammates clearly} : il s'agit d'un casque-écouteur. Attention à l'article : \eng{a new headset} (nom comptable au singulier).\par
3. \eng{download}. L'indice \eng{data bundles are cheaper} évoque le téléchargement de fichiers. Après la structure \eng{it is cheaper to + base verbale}, on met l'infinitif sans \emph{to} : \eng{download}.\par
4. \eng{Internet café}. L'indice \eng{no Wi-Fi at home} et \eng{near the market} désigne un lieu public d'accès à Internet. En anglais, le terme standard est \eng{Internet café} (ou \eng{cybercafe} sans accent) ; article défini \eng{the} ici car le lieu est précisé par \eng{near the market}.\par
5. \eng{stream}. Écouter de la musique en ligne sans la télécharger = \eng{to stream music}. Présent simple, sujet \eng{we} : \eng{stream}.\par
Conclusion : chaque réponse est justifiée par un indice de sens dans la phrase et par la grammaire (forme verbale, article).
\end{exoresolu}

\courssub{Recreational activities online and offline}

\begin{methode}[Réemployer le vocabulaire dans un dialogue ou un paragraphe]
Pour produire un texte oral ou écrit riche et naturel :
\begin{enumerate}
\item \cle{Prépare une banque de 8 à 10 mots} du thème avant de parler ou d'écrire (ex. \eng{gaming console, to stream, social media, screen time, offline, Wi-Fi hotspot, data bundle, voice note}).
\item \cle{Structure ta production} : introduction (situation), développement (activités, avantages, inconvénients), conclusion (opinion personnelle).
\item \cle{Utilise des connecteurs} : \eng{first of all, moreover, however, as a result, in my opinion}.
\item \cle{Varie les formulations} : ne répète pas \eng{play} partout ; alterne \eng{enjoy, spend time on, be fond of, take part in}.
\item \cle{Donne des exemples concrets} du contexte camerounais : \eng{watching football matches online, sharing videos on WhatsApp, playing games at a gaming lounge in Douala}.
\item \cle{À l'oral}, parle lentement, articule les finales (\eng{downloads, chats}) et utilise des phrases courtes.
\end{enumerate}
\end{methode}

\begin{exoresolu}[Rédiger un court paragraphe avec le lexique imposé]
Énoncé — Write a paragraph of about 80 words on: “How young people in your community use ICTs for recreation.” Use at least five of these words: \emph{smartphone, stream, social media, gaming, online, screen time, relax}.
\tcblower
Solution détaillée :\par
Démarche : on identifie d'abord les mots obligatoires, puis on construit un plan simple : (a) les outils utilisés, (b) les activités, (c) un point de vigilance, (d) une opinion.\par
Réponse rédigée :\par
\eng{In my community, most young people own a smartphone, which they use for many recreational activities. After school, they go online to chat with friends on social media such as WhatsApp and Facebook. Many of them stream music and watch comedy videos to relax, while others prefer gaming with their friends at the weekend. However, parents often complain that their children's screen time is too long. In my opinion, ICTs are wonderful tools for entertainment, but young people should balance online fun with outdoor activities.}\par
Justification : les six mots imposés sont employés avec leurs collocations correctes (\eng{own a smartphone, go online, stream music, screen time}) ; les connecteurs \eng{while, however, in my opinion} structurent le texte ; le présent simple est utilisé pour décrire des habitudes. Conclusion : un paragraphe clair, naturel, d'environ 85 mots, conforme à la consigne.
\end{exoresolu}

\courssub{Collocations and idiomatic expressions}

\begin{methode}[Apprendre et utiliser les collocations des loisirs numériques]
\begin{enumerate}
\item \cle{Mémorise les blocs de mots}, pas les mots seuls : \eng{to spend time online, to kill time, to binge-watch a series, to scroll through a feed, to go viral, to log in / log out, to recharge data}.
\item \cle{Distingue les paires proches} : \eng{to make friends online} (pas \eng{to do friends}), \eng{to take a selfie} (pas \eng{to make a selfie}), \eng{to have fun} (pas \eng{to do fun}).
\item \cle{Repère le verbe support} : en anglais, c'est souvent le verbe qui impose le nom : \eng{play a game, watch a video, send a message, post a photo}.
\item \cle{Teste la collocation} en la remplaçant par un synonyme : si la phrase devient bizarre, la combinaison est fautive.
\item \cle{Tiens un carnet de collocations} classé par verbe (\eng{play, watch, spend, log}).
\end{enumerate}
\end{methode}

\begin{exoresolu}[Corriger les collocations fautives]
Énoncé — Each sentence contains a wrong collocation. Correct it.\par
1. My sister makes selfies every weekend and posts them online.\par
2. We did a lot of fun at the gaming lounge in Buea.\par
3. He spends three hours to watch videos every day.\par
4. I want to do new friends on social media.\par
5. Turn on your account before you leave the computer room.
\tcblower
Solution détaillée :\par
1. \eng{My sister \textbf{takes} selfies every weekend and posts them online.} La collocation correcte est \eng{to take a selfie} (comme \eng{to take a photo}) ; \eng{make} est un calque du français « faire ».\par
2. \eng{We \textbf{had} a lot of fun at the gaming lounge in Buea.} On dit \eng{to have fun} ; \eng{fun} est ici un nom indénombrable, d'où \eng{a lot of fun}.\par
3. \eng{He spends three hours \textbf{watching} videos every day.} Structure obligatoire : \eng{to spend + time + V-ing} ; jamais d'infinitif après \eng{spend time}.\par
4. \eng{I want to \textbf{make} new friends on social media.} Paradoxalement, c'est ici \eng{make} qui convient : \eng{to make friends}. Le français « se faire des amis » piège souvent dans l'autre sens.\par
5. \eng{\textbf{Log out of} your account before you leave the computer room.} Pour un compte, on dit \eng{to log in / log out} ; \eng{turn on/off} s'emploie pour l'appareil (\eng{turn off the computer}), pas pour la session.\par
Conclusion : chaque correction repose sur une collocation figée de l'anglais ; il faut les apprendre comme des unités.
\end{exoresolu}

\courssub{False friends and French traps}

\begin{methode}[Déjouer les faux amis du domaine des TIC]
\begin{enumerate}
\item \cle{Repère les mots « trop faciles »} qui ressemblent au français : \eng{actually} ne veut pas dire « actuellement » mais « en fait » ; « actuellement » se dit \eng{currently} ou \eng{nowadays}.
\item \cle{Méfie-toi des calques de structure} : « je suis connecté » = \eng{I am connected / online}, pas une forme bizarre ; « passer du temps » = \eng{to spend time} (pas \eng{to pass time}, sauf \eng{to pass the time} = tromper l'ennui).
\item \cle{Vérifie les noms} : \eng{a library} = bibliothèque, mais une « librairie » = \eng{a bookshop} ; \eng{data} est traité comme indénombrable au sens courant (\eng{mobile data is expensive}).
\item \cle{Attention aux verbes} : « enregistrer une émission » = \eng{to record a programme} ; \eng{to register} = s'inscrire.
\item \cle{En cas de doute}, reformule avec un mot simple et sûr plutôt que de risquer un faux ami.
\end{enumerate}
\end{methode}

\begin{exoresolu}[Traduire sans tomber dans les pièges]
Énoncé — Translate into English:\par
1. Actuellement, beaucoup de jeunes passent leur temps libre sur les réseaux sociaux.\par
2. J'ai enregistré le match de football pour le regarder ce soir.\par
3. Les données mobiles sont chères au Cameroun, alors je vais au cybercafé.
\tcblower
Solution détaillée :\par
1. \eng{Nowadays, many young people spend their free time on social media.} Piège n°1 : « actuellement » ne se traduit jamais par \eng{actually} ; on choisit \eng{nowadays} ou \eng{currently}. Piège n°2 : « passer du temps » = \eng{to spend time} ; « les réseaux sociaux » = \eng{social media} (pluriel d'usage).\par
2. \eng{I recorded the football match to watch it this evening.} « Enregistrer » au sens de capturer une émission = \eng{to record} ; \eng{to register} signifie « s'inscrire ». Le prétérit simple \eng{recorded} convient car l'action est terminée ; \eng{to watch it} exprime le but.\par
3. \eng{Mobile data is expensive in Cameroon, so I go to the Internet café.} \eng{Data} est indénombrable au sens courant : verbe au singulier \eng{is}, sans article. « Alors » de conséquence = \eng{so} ; \eng{Internet café} est le terme standard (on évite \eng{cybercafé} avec accent).\par
Conclusion : la traduction correcte exige de neutraliser trois pièges classiques : \eng{actually}, \eng{pass time}, \eng{register}.
\end{exoresolu}

\courssub{Word families, pronunciation and spelling}

\begin{methode}[Former et prononcer correctement les mots de la même famille]
\begin{enumerate}
\item \cle{Identifie la racine} puis les suffixes : \eng{entertain} (verbe) → \eng{entertainment} (nom) → \eng{entertaining} (adjectif) ; \eng{connect} → \eng{connection} → \eng{connected}.
\item \cle{Retiens les suffixes typiques} : \emph{-ment}, \emph{-tion/-sion}, \emph{-ing} (noms) ; \emph{-ive}, \emph{-ing}, \emph{-ed} (adjectifs) ; \emph{-ise/-ize} (verbes : \eng{computerise}).
\item \cle{Respecte les règles d'orthographe} : consonne doublée après voyelle courte accentuée (\eng{chat → chatting, chatted}) ; \emph{e} muet qui saute devant \emph{-ing} (\eng{game → gaming, share → sharing}).
\item \cle{Soigne la prononciation} : \eng{technology} se dit « tek-no-lo-dji » ; \eng{download} (nom) s'accentue sur la première syllabe (« downe-lode ») ; le \emph{th} de \eng{the} se prononce avec la langue entre les dents.
\item \cle{Vérifie la nature grammaticale} exigée par la phrase avant de choisir la forme du mot.
\end{enumerate}
\end{methode}

\begin{exoresolu}[Trouver la bonne forme du mot]
Énoncé — Complete each sentence with the correct form of the word in brackets.\par
1. Online games are a popular form of \ldots\ (entertain) for teenagers.\par
2. The \ldots\ (connect) in our village is too slow for streaming.\par
3. She spends hours \ldots\ (chat) with her cousins in Bamenda.\par
4. This documentary about digital life is very \ldots\ (interest).\par
5. \ldots\ (game) can be fun, but it should not replace sport.
\tcblower
Solution détaillée :\par
1. \eng{entertainment}. Après la préposition \eng{of}, il faut un nom : suffixe \emph{-ment} → \eng{entertainment} (« a popular form of entertainment »).\par
2. \eng{connection}. L'article \eng{the} appelle un nom ; \eng{connect} → \eng{connection} (suffixe \emph{-tion}). « The connection in our village is too slow for streaming. »\par
3. \eng{chatting}. Structure \eng{to spend time + V-ing}. Orthographe : voyelle courte \emph{a} + consonne finale \emph{t} → on double : \eng{chatting} (comme \eng{sitting, swimming}).\par
4. \eng{interesting}. Le documentaire \emph{provoque} l'intérêt : adjectif en \emph{-ing} pour la chose, \emph{-ed} pour la personne qui ressent (\eng{I am interested}). Ici le sujet est la chose : \eng{interesting}.\par
5. \eng{Gaming}. En tête de phrase comme sujet, on utilise le gérondif : \eng{Gaming can be fun}. Orthographe : \emph{e} muet de \eng{game} qui saute devant \emph{-ing} → \eng{gaming} (et non « gameing »).\par
Conclusion : la nature grammaticale (nom, adjectif, gérondif) et les règles d'orthographe (consonne doublée, \emph{e} muet) guident chaque réponse.
\end{exoresolu}

\courssub{Using the vocabulary in context}

\begin{experience}[Jeu de rôle — « Weekend plans at the gaming lounge »]
En binômes. L'élève A est un habitué du \eng{gaming lounge} de Douala ; l'élève B est un nouveau venu qui veut savoir comment s'inscrire, quels jeux sont disponibles, le prix à l'heure, et s'il y a une connexion fibre. Utilisez au moins 8 mots du thème (\eng{headset, console, stream, lag, data bundle, Wi-Fi hotspot, screen time, recharge}). Changez de rôle. L'enseignant note la fluidité, la prononciation des finales (\eng{-s, -ed}) et l'usage correct des prépositions (\eng{addicted to, log on to, spend time on}).
\end{experience}

\begin{savaistu}[Le numérique au Cameroun : enjeux et réalités]
Le Cameroun compte plus de 10 millions d'abonnés Internet (données ARTEC 2023), majoritairement via mobile (3G/4G). Les \eng{data bundles} (forfaits data) sont le mode d'accès principal ; les opérateurs (MTN, Orange, Camtel, Nexttel) proposent des offres « nuit » moins chères pour le streaming et le téléchargement. Les \eng{Internet cafés} (parfois appelés \eng{cybercafés}) restent populaires dans les quartiers périphériques et les zones rurales où la fibre n'est pas déployée. Le \eng{mobile money} (Mobile Money, MoMo) sert souvent à recharger son crédit data. Le football (Ligue 1, Premier League, CAN) est le contenu le plus streamé. Attention : \eng{screen time} (temps d'écran) devient un sujet de santé publique ; le MINESEC encourage l'éducation aux médias.
\end{savaistu}

\begin{exercice}[Entraînement personnel — Vocabulary in use]
\begin{enumerate}
\item Complete with the correct word: \emph{stream, scroll, recharge, lag, hotspot, addicted, bundle, console}.
\begin{enumerate}
\item My phone creates a Wi-Fi \ldots\ so my laptop can connect.
\item During the match, the video started to \ldots\ because of a bad connection.
\item She is \ldots\ to TikTok; she spends five hours a day on it.
\item I need to \ldots\ my data \ldots\ before midnight.
\item They bought a new PlayStation \ldots\ for the living room.
\item We just \ldots\ through the feed without really watching.
\end{enumerate}
\item Choose the correct collocation.
\begin{enumerate}
\item She \textbf{makes / takes / does} a selfie with her friends.
\item We \textbf{had / made / did} a lot of fun at the LAN party.
\item He wants to \textbf{make / do / take} new friends on Discord.
\item Don't forget to \textbf{log in / log on / turn on} your account.
\end{enumerate}
\item Write the correct form of the word in brackets.
\begin{enumerate}
\item The \ldots\ (connect) was lost during the storm.
\item This new series is really \ldots\ (entertain).
\item \ldots\ (download) large files takes time on 3G.
\item He is \ldots\ (interest) in coding.
\item They spent the afternoon \ldots\ (game).
\end{enumerate}
\item Translate into English.
\begin{enumerate}
\item Actuellement, je télécharge un film pour ce soir.
\item Les jeunes passent trop de temps sur les réseaux sociaux.
\item J'ai enregistré l'émission pour la regarder plus tard.
\end{enumerate}
\end{enumerate}
\end{exercice}

\begin{corrige}[Exercice n°1 -- Vocabulary in use]
\begin{enumerate}
\item \begin{enumerate}
\item \eng{hotspot} (Wi-Fi hotspot = point d'accès).
\item \eng{lag} (verbe : la vidéo saccade).
\item \eng{addicted to} (piège : \eng{addicted} + \eng{to}).
\item \eng{recharge ... bundle} (collocation : \eng{recharge a data bundle}).
\item \eng{console} (gaming console).
\item \eng{scroll} (to scroll through a feed).
\end{enumerate}
\item \begin{enumerate}
\item \eng{takes} (\eng{take a selfie}).
\item \eng{had} (\eng{have fun}).
\item \eng{make} (\eng{make friends}).
\item \eng{log in} (ou \eng{log on to}) — pour une session ; \eng{turn on} = allumer l'appareil.
\end{enumerate}
\item \begin{enumerate}
\item \eng{connection} (nom, suffixe \emph{-tion}).
\item \eng{entertaining} (adjectif en \emph{-ing} pour la chose qui divertit).
\item \eng{Downloading} (gérondif sujet, \emph{e} muet sauté).
\item \eng{interested} (adjectif en \emph{-ed} pour la personne qui ressent).
\item \eng{gaming} (gérondif, \emph{e} muet sauté).
\end{enumerate}
\item \begin{enumerate}
\item \eng{Currently, I am downloading a film for this evening.} (Pas \eng{actually} ; \eng{downloading} = en train de télécharger).
\item \eng{Young people spend too much time on social media.} (\eng{spend time}, pas \eng{pass time} ; \eng{social media}).
\item \eng{I recorded the programme to watch it later.} (\eng{recorded} = enregistré ; \eng{later} = plus tard).
\end{enumerate}
\end{enumerate}
\end{corrige}

\begin{attention}[Les 5 erreurs qui coûtent des points au Bac]
\begin{enumerate}
\item \textbf{Faux ami \eng{actually}.} \eng{Actually} signifie « en fait », jamais « actuellement ». Écris \eng{Nowadays / Currently, teenagers spend hours online.}\par
\item \textbf{Calque « passer du temps ».} On dit \eng{to \textbf{spend} time online} et jamais « to pass time » ; et après \eng{spend}, le verbe est en \emph{-ing} : \eng{She spends two hours \textbf{watching} videos.}\par
\item \textbf{Collocations avec \eng{make/do/take}.} \eng{to \textbf{take} a selfie}, \eng{to \textbf{make} friends}, \eng{to \textbf{have} fun}, \eng{to \textbf{play} a game}. Intervertir ces verbes est systématiquement sanctionné.\par
\item \textbf{Orthographe des formes en \emph{-ing}.} Consonne doublée : \eng{chatting, running, stopping} ; \emph{e} muet supprimé : \eng{gaming, sharing, downloading} (mais \eng{download} garde sa forme car pas d'\emph{e} final).\par
\item \textbf{Prépositions et mots indénombrables.} \eng{to listen \textbf{to} music}, \eng{to be addicted \textbf{to} gaming}, \eng{to log \textbf{on to} a site} ; \eng{data} et \eng{information} sont indénombrables : \eng{Mobile data \textbf{is} expensive}, sans \emph{s} ni article indéfini.
\end{enumerate}
\end{attention}

\newpage

\coursec{Exercices}

\courssub{Je vérifie mes connaissances}

\begin{exercice}[QCM : le bon mot]
Entoure la lettre de la réponse correcte.
\begin{enumerate}
\item At weekends, I usually \ldots\ videos on my phone.
\begin{enumerate}
\item[a)] stream \item[b)] steam \item[c)] scream
\end{enumerate}
\item My little brother is \ldots\ video games; he plays for hours.
\begin{enumerate}
\item[a)] glued on \item[b)] hooked on \item[c)] hanged on
\end{enumerate}
\item This funny video went \ldots\ on social media in one night.
\begin{enumerate}
\item[a)] virtual \item[b)] virus \item[c)] viral
\end{enumerate}
\item I need to buy a new \ldots\ to watch the match online.
\begin{enumerate}
\item[a)] data bundle \item[b)] date bundle \item[c)] data candle
\end{enumerate}
\item Many young people in Douala play games at a \ldots\ near the market.
\begin{enumerate}
\item[a)] gaming centre \item[b)] gambling centre \item[c)] gaming sentry
\end{enumerate}
\end{enumerate}
\end{exercice}

\begin{exercice}[Vrai ou faux ?]
Réponds par \eng{True} ou \eng{False} et justifie ta réponse par une courte phrase en anglais.
\begin{enumerate}
\item \eng{A cybercafé is a place where people only drink coffee.}
\item \eng{“Screen time” means the time spent in front of a screen.}
\item \eng{To “binge-watch” means to watch only one episode of a series.}
\item \eng{“To kill time” means to do something while waiting.}
\item \eng{“Actually” means “actuellement” in English.}
\end{enumerate}
\end{exercice}

\begin{exercice}[Vocabulaire : définitions]
Trouve le mot ou l'expression de la leçon qui correspond à chaque définition.
\begin{enumerate}
\item A small room or shop with computers and internet access that people pay to use: \ldots
\item To transfer a file from the internet onto your device: \ldots
\item Unable to stop doing something, like a game or a series: \ldots
\item Websites and applications like Facebook, TikTok or WhatsApp: \ldots
\item To move quickly through posts on your phone with your finger: \ldots
\end{enumerate}
\end{exercice}

\begin{exercice}[Familles de mots : à compléter]
Recopie et complète le tableau avec la forme qui manque.
\begin{center}
\begin{tabular}{|l|l|l|}\hline
\rowcolor{popblueL} Verb & Noun & Adjective \\ \hline
to entertain & \ldots & entertaining \\ \hline
to relax & relaxation & \ldots \\ \hline
to addict & addiction & \ldots \\ \hline
to connect & \ldots & connected \\ \hline
--- & recreation & \ldots \\ \hline
\end{tabular}
\end{center}
\end{exercice}

\courssub{Je m'entraîne}

\begin{exercice}[Traduction]
Traduis les phrases suivantes en anglais en utilisant le vocabulaire de la leçon.
\begin{enumerate}
\item Les jeunes de mon quartier passent trop de temps devant l'écran.
\item Hier soir, j'ai regardé toute la saison d'une série d'un coup.
\item Cette vidéo de danse est devenue virale au Cameroun.
\item Mon forfait internet est épuisé, je ne peux plus télécharger de musique.
\item Pour tuer le temps, elle fait défiler son fil d'actualité.
\end{enumerate}
\end{exercice}

\begin{exercice}[Appariement : collocations]
Relie chaque verbe ou expression de la colonne A à son complément de la colonne B.
\begin{center}
\begin{tabular}{|l|l|}\hline
\rowcolor{popblueL} A & B \\ \hline
1. to play & a) viral \\ \hline
2. to binge- & b) time online \\ \hline
3. to go & c) through one's feed \\ \hline
4. to scroll & d) watch a series \\ \hline
5. to kill & e) online games \\ \hline
\end{tabular}
\end{center}
\end{exercice}

\begin{exercice}[Reformulation avec un mot de la même famille]
Réécris chaque phrase en utilisant la forme du mot entre parenthèses, sans changer le sens.
\begin{enumerate}
\item \eng{Gaming relaxes me.} (relaxation)
\item \eng{The show entertained the children.} (entertainment)
\item \eng{Many teenagers are addicted to social media.} (addiction)
\item \eng{We connect to the internet at the cybercafé.} (connection)
\item \eng{Playing football is a recreational activity.} (recreation)
\end{enumerate}
\end{exercice}

\begin{exercice}[Dialogue à compléter : au cybercafé]
Complète le dialogue entre deux amis avec les expressions suivantes : \eng{glued to the screen}, \eng{kill time}, \eng{hooked on}, \eng{data bundle}, \eng{screen time}.
\begin{textebox}[At the cybercafé in Bamenda]
A: Hi Sandra! What are you doing here?\\
B: Oh, just trying to \ldots\ while I wait for my brother.\\
A: Look at that boy over there. He has been \ldots\ for three hours!\\
B: I know. He is completely \ldots\ online games.\\
A: And you? Is your \ldots\ very high these days?\\
B: Not really. My \ldots\ is finished, so I can only use the computers here.
\end{textebox}
\end{exercice}

\courssub{Je me prépare au Bac}

\begin{exercice}[Texte à trous type Bac]
Complète l'article suivant avec les dix mots et expressions de la liste : \eng{stream}, \eng{data bundle}, \eng{binge-watch}, \eng{go viral}, \eng{gaming centre}, \eng{screen time}, \eng{hooked on}, \eng{download}, \eng{cybercafé}, \eng{social media}.
\begin{textebox}[Digital leisure among young Cameroonians]
In many towns of Cameroon, young people now spend their free time differently. Instead of playing outside all evening, many of them (1) \ldots\ music and films on their phones. Those who do not have a smartphone go to the nearest (2) \ldots\ , where they pay a small fee to use a computer. Others prefer the (3) \ldots\ in the neighbourhood, where they play online games with friends.

However, parents worry about their children's (4) \ldots\ , which can reach six or seven hours a day. Some teenagers are so (5) \ldots\ their favourite series that they (6) \ldots\ a whole season in one weekend. Funny videos can (7) \ldots\ in a few hours and be shared by thousands of users on (8) \ldots\ . To enjoy all this, you need a good (9) \ldots\ , which is not always cheap. When it is finished, you can no longer (10) \ldots\ anything until you buy another one.
\end{textebox}
\end{exercice}

\begin{exercice}[Compréhension et langue : texte type Bac]
Lis le texte suivant, puis réponds aux questions en anglais, avec des phrases complètes.
\begin{textebox}[A new way to relax]
Every Saturday afternoon, the small gaming centre near the main market in Buea is full of young people. Some come to play online football games, others to chat with friends on social media or to watch videos. For many of them, these activities are a cheap and pleasant way to relax after a long week at school.

“I come here to kill time and have fun with my friends,” says Junior, a sixteen-year-old student. “At home, my data bundle is never enough, so I prefer the gaming centre.” His friend Divine, however, spends most of her free time on her phone. “I am hooked on a new series,” she admits. “Last weekend I binge-watched ten episodes in two days.”

Not everyone is happy with this new form of recreation. Some parents complain that their children are glued to their screens and no longer help with housework. Teachers also notice that too much screen time can affect students' concentration in class. “ICTs are wonderful tools,” says Mr Ndi, an English teacher, “but young people must learn to use them wisely and not let them take over their lives.”
\end{textebox}
\textbf{Comprehension questions}
\begin{enumerate}
\item Where is the gaming centre described in the text?
\item Why does Junior prefer the gaming centre to his home?
\item What did Divine do last weekend?
\item Give two complaints adults make about young people's use of ICTs.
\item What advice does Mr Ndi give?
\end{enumerate}
\textbf{Language work}
\begin{enumerate}
\item Find in the text two expressions that mean “to spend time doing something unimportant”.
\item Find in the text the adjective formed from the verb \eng{to addict}.
\item Put into the plural: \eng{a gaming centre}, \eng{a data bundle}, \eng{an episode}.
\item Rewrite in the passive voice: \eng{Young people share funny videos on social media.}
\end{enumerate}
\end{exercice}

\begin{exercice}[Composition type Bac]
Rédige une composition en anglais sur le sujet suivant :
\begin{center}
\emph{“Discuss the advantages and disadvantages of using ICTs for recreation among young people in your community.”}
\end{center}
Consignes :
\begin{itemize}
\item Longueur attendue : entre 200 et 250 mots.
\item Structure imposée : introduction, deux paragraphes (avantages / inconvénients), conclusion avec ton opinion personnelle.
\item Tu dois employer au moins huit mots ou expressions du lexique de la leçon, par exemple : \eng{screen time}, \eng{hooked on}, \eng{binge-watch}, \eng{go viral}, \eng{data bundle}, \eng{gaming centre}, \eng{cybercafé}, \eng{social media}, \eng{kill time}, \eng{glued to the screen}.
\item Soigne la grammaire, l'orthographe et la ponctuation.
\end{itemize}
\end{exercice}

\begin{methode}[Réussir la composition type Bac]
\begin{enumerate}
\item \textbf{Analyse le sujet} : il s'agit d'une discussion équilibrée (avantages \emph{et} inconvénients), située dans ta communauté. Note les mots-clés : \eng{advantages}, \eng{disadvantages}, \eng{recreation}, \eng{young people}, \eng{your community}.
\item \textbf{Fais un plan rapide} au brouillon : deux ou trois idées par paragraphe, chacune illustrée par un exemple local (cybercafé de quartier, salle de jeux, forfait internet).
\item \textbf{Rédige l'introduction} : présente le sujet en deux ou trois phrases et annonce le plan, par exemple : \eng{This essay will discuss the advantages and disadvantages of using ICTs for recreation among young people in my community.}
\item \textbf{Développe} avec des connecteurs : \eng{first of all}, \eng{moreover}, \eng{on the one hand / on the other hand}, \eng{however}, \eng{in addition}, \eng{as a result}.
\item \textbf{Conclus} avec ton opinion personnelle introduite par \eng{in my opinion} ou \eng{to my mind}, sans apporter d'idée nouvelle.
\item \textbf{Relis-toi} : vérifie les accords sujet-verbe, les temps (présent de vérité générale), les articles, les faux amis (\eng{actually} $\neq$ « actuellement ») et le nombre de mots.
\end{enumerate}
\end{methode}

\begin{exemplebox}[Modèle d'introduction et de transition]
\eng{Nowadays, information and communication technologies have changed the way young people spend their free time in my community. Instead of playing outside, many teenagers prefer their phones, the cybercafé or the gaming centre. This essay will examine the advantages and disadvantages of this new form of recreation.}

\eng{On the one hand, ICTs offer cheap and pleasant entertainment. On the other hand, excessive screen time can become a serious problem.}
\end{exemplebox}

\begin{attention}[Pièges à éviter dans la composition]
\begin{itemize}
\item Ne traduis pas « passer du temps » par \eng{pass time} : dis \eng{spend time}.
\item \eng{Actually} signifie « en fait » ; pour « actuellement », utilise \eng{nowadays} ou \eng{currently}.
\item \eng{To be addicted to} et \eng{to be hooked on} sont toujours suivis de la préposition \eng{to}/\eng{on}, jamais de \eng{for}.
\item \eng{Social media} se construit sans article dans le sens général : \eng{on social media}.
\item Évite le calque « les TIC » : en anglais, on dit \eng{ICTs} ou \eng{information and communication technologies}.
\end{itemize}
\end{attention}

\newpage

\coursec{Corrigés des exercices}

\begin{corrige}[Exercice 1 — QCM : le bon mot]
\begin{enumerate}
\item \textbf{a)} \eng{stream} — « regarder en streaming ». \eng{stream} (verbe) = diffuser/regarder en continu ; \eng{steam} = « vapeur » ; \eng{scream} = « crier ». Phrase complète : \eng{At weekends, I usually stream videos on my phone.}
\item \textbf{b)} \eng{hooked on} — « accro à ». La collocation correcte est \eng{to be hooked on something} ; \eng{glued} se combine avec \eng{to} (\eng{glued to the screen}) et non \eng{on} ; \eng{hanged on} n'existe pas dans ce sens.
\item \textbf{c)} \eng{viral} — \eng{to go viral} = « devenir viral ». \eng{virtual} = « virtuel » (faux ami partiel avec « viral »), \eng{virus} = « virus ».
\item \textbf{a)} \eng{data bundle} — « forfait internet ». \eng{date} = « date » ou « rendez-vous » ; \eng{candle} = « bougie ».
\item \textbf{a)} \eng{gaming centre} — « salle de jeux vidéo ». \eng{gambling centre} = « salle de jeux d'argent », sens différent ; \eng{sentry} = « sentinelle ».
\end{enumerate}
\end{corrige}

\begin{corrige}[Exercice 2 — Vrai ou faux ?]
\begin{enumerate}
\item \eng{False. A cybercafé is a place where people pay to use computers and the internet, not a place where people only drink coffee.}
\item \eng{True. “Screen time” is the amount of time someone spends looking at a screen.}
\item \eng{False. To binge-watch means to watch many episodes of a series one after another, not just one.}
\item \eng{True. To kill time means to do something unimportant while you are waiting.}
\item \eng{False. “Actually” means “en fait / en réalité”; “actuellement” is translated by “currently” or “at present”.}
\end{enumerate}
Point de méthode : au Bac, une justification courte mais complète (sujet + verbe) est exigée ; un simple \eng{False} sans phrase ne rapporte pas tous les points.
\end{corrige}

\begin{corrige}[Exercice 3 — Vocabulaire : définitions]
\begin{enumerate}
\item \eng{a cybercafé} (on accepte aussi \eng{an internet café}) — « cybercafé ».
\item \eng{to download} — « télécharger » (de l'internet vers l'appareil ; l'inverse est \eng{to upload}).
\item \eng{addicted} (ou l'expression \eng{hooked on}) — « accro, dépendant ».
\item \eng{social media} — « réseaux sociaux » (sans article dans le sens général).
\item \eng{to scroll (through one's feed)} — « faire défiler (son fil d'actualité) ».
\end{enumerate}
\end{corrige}

\begin{corrige}[Exercice 4 — Familles de mots : à compléter]
\begin{center}
\begin{tabular}{|l|l|l|}\hline
\rowcolor{popblueL} Verb & Noun & Adjective \\ \hline
to entertain & \textbf{entertainment} & entertaining \\ \hline
to relax & relaxation & \textbf{relaxing / relaxed} \\ \hline
to addict & addiction & \textbf{addictive / addicted} \\ \hline
to connect & \textbf{connection} & connected \\ \hline
--- & recreation & \textbf{recreational} \\ \hline
\end{tabular}
\end{center}
Remarques : \eng{relaxing} décrit la chose qui détend (\eng{a relaxing game}), \eng{relaxed} décrit la personne détendue (\eng{a relaxed student}) ; même distinction entre \eng{addictive} (\eng{an addictive game}) et \eng{addicted} (\eng{an addicted teenager}). \eng{recreation} n'a pas de verbe courant correspondant ; on utilise \eng{to relax}, \eng{to have fun} ou \eng{to enjoy oneself}. Prononciation : \eng{recreation} (loisir) se dit « ré-kri-é-cheune », à ne pas confondre avec \eng{re-creation} « ri-kri-é-cheune » (nouvelle création).
\end{corrige}

\begin{corrige}[Exercice 5 — Traduction]
\begin{enumerate}
\item \eng{The young people in my neighbourhood spend too much time in front of their screens.} (ou \eng{...spend too much screen time.})
\item \eng{Last night, I binge-watched a whole season of a series.} (ou \eng{...watched a whole season of a series in one go.})
\item \eng{This dance video has gone viral in Cameroon.} (prétérit accepté : \eng{went viral})
\item \eng{My data bundle is finished, so I can no longer download music.}
\item \eng{To kill time, she scrolls through her news feed.}
\end{enumerate}
Points de vigilance : « passer du temps » se traduit par \eng{to spend time} (jamais \eng{to pass time}) ; « d'un coup » se rend par \eng{binge-watch} ou \eng{in one go} ; « fil d'actualité » = \eng{news feed} ; « épuisé » pour un forfait = \eng{finished} ou \eng{exhausted}, pas \eng{tired}.
\end{corrige}

\begin{corrige}[Exercice 6 — Appariement : collocations]
1 -- e) \eng{to play online games} (« jouer à des jeux en ligne ») ; 2 -- d) \eng{to binge-watch a series} (« regarder une série en rafale ») ; 3 -- a) \eng{to go viral} (« devenir viral ») ; 4 -- c) \eng{to scroll through one's feed} (« faire défiler son fil d'actualité ») ; 5 -- b) \eng{to kill time online} (« tuer le temps en ligne »).

Justification : en anglais, les collocations sont fixes ; on ne peut pas dire \eng{play viral} ni \eng{kill a series}. Apprends chaque verbe avec son complément habituel comme un seul bloc.
\end{corrige}

\begin{corrige}[Exercice 7 — Reformulation avec un mot de la même famille]
\begin{enumerate}
\item \eng{Gaming gives me relaxation.} (ou \eng{I find relaxation in gaming.})
\item \eng{The show provided entertainment for the children.} (ou \eng{The show was entertainment for the children.})
\item \eng{Many teenagers have an addiction to social media.} (ou \eng{Many teenagers suffer from an addiction to social media.})
\item \eng{We get an internet connection at the cybercafé.} (ou \eng{Our connection to the internet is at the cybercafé.})
\item \eng{Playing football is a form of recreation.}
\end{enumerate}
Point de grammaire : quand on passe du verbe au nom, il faut souvent introduire un verbe support (\eng{give}, \eng{provide}, \eng{have}, \eng{get}) et vérifier l'article (\eng{an addiction}, \eng{a form of}).
\end{corrige}

\begin{corrige}[Exercice 8 — Dialogue à compléter : au cybercafé]
Dans l'ordre : \eng{kill time} ; \eng{glued to the screen} ; \eng{hooked on} ; \eng{screen time} ; \eng{data bundle}.

Dialogue complet :
\eng{A: Hi Sandra! What are you doing here? — B: Oh, just trying to kill time while I wait for my brother. — A: Look at that boy over there. He has been glued to the screen for three hours! — B: I know. He is completely hooked on online games. — A: And you? Is your screen time very high these days? — B: Not really. My data bundle is finished, so I can only use the computers here.}

Justifications : après \eng{trying to}, on met la base verbale (\eng{kill time}) ; \eng{has been glued to the screen} est un present perfect avec \eng{for three hours} ; \eng{hooked on} est suivi du nom \eng{online games} ; la question \eng{Is your screen time very high?} porte sur une quantité ; la dernière réplique parle d'argent et d'internet, d'où \eng{data bundle}.
\end{corrige}

\begin{corrige}[Exercice 9 — Texte à trous type Bac]
(1) \eng{stream} ; (2) \eng{cybercafé} ; (3) \eng{gaming centre} ; (4) \eng{screen time} ; (5) \eng{hooked on} ; (6) \eng{binge-watch} ; (7) \eng{go viral} ; (8) \eng{social media} ; (9) \eng{data bundle} ; (10) \eng{download}.

Justifications : (1) le sujet \eng{many of them} impose un verbe au présent pluriel ; (2) le lieu où l'on paie pour utiliser un ordinateur est le \eng{cybercafé} ; (3) le lieu des jeux en ligne est le \eng{gaming centre} ; (4) le complément de \eng{their children's} est un nom, \eng{screen time} ; (5) la structure \eng{so + adjective + that} impose \eng{hooked on} ; (6) après \eng{they}, base verbale \eng{binge-watch} ; (7) après \eng{can}, base verbale \eng{go viral} ; (8) \eng{on social media} est la préposition attendue ; (9) \eng{a good} + nom singulier : \eng{data bundle} ; (10) après \eng{can no longer}, base verbale \eng{download}.

Barème indicatif : 1 point par réponse correcte (10 points) ; une faute d'orthographe dans le mot recopié coûte 0,5 point.
\end{corrige}

\begin{corrige}[Exercice 10 — Compréhension et langue : texte type Bac]
\textbf{Comprehension}
\begin{enumerate}
\item \eng{The gaming centre is near the main market in Buea.}
\item \eng{Junior prefers the gaming centre because his data bundle is never enough at home.}
\item \eng{Last weekend, Divine binge-watched ten episodes of a new series in two days.}
\item \eng{Parents complain that their children are glued to their screens and no longer help with housework; teachers notice that too much screen time affects students' concentration in class.}
\item \eng{Mr Ndi advises young people to use ICTs wisely and not to let them take over their lives.}
\end{enumerate}
\textbf{Language work}
\begin{enumerate}
\item \eng{to kill time} (l'expression qui signifie « occuper son temps à quelque chose de futile ») ; on accepte aussi \eng{to have fun} comme activité sans but précis.
\item \eng{addicted} ; le texte utilise son synonyme \eng{hooked on}, mais l'adjectif de la famille de \eng{to addict} est \eng{addicted} (préposition \eng{to}).
\item \eng{gaming centres}, \eng{data bundles}, \eng{episodes} (pluriel régulier en \eng{-s} ; pas de changement d'orthographe).
\item \eng{Funny videos are shared on social media (by young people).} — le complément d'objet \eng{funny videos} devient sujet, le verbe \eng{share} passe à \eng{are shared} (présent passif : \eng{be} + participe passé).
\end{enumerate}
Barème indicatif : compréhension 5 $\times$ 2 points = 10 points ; langue 4 $\times$ 2,5 points = 10 points. Toute réponse de compréhension doit être une phrase complète ; recopier un fragment du texte sans verbe conjugué ne vaut que la moitié des points.
\end{corrige}

\begin{corrige}[Exercice 11 — Composition type Bac]
\textbf{Barème indicatif (sur 20)} : respect du sujet et de la structure (4) ; richesse et exactitude du lexique de la leçon, au moins huit items (4) ; correction grammaticale (4) ; orthographe et ponctuation (3) ; organisation et connecteurs (3) ; longueur respectée, 200 à 250 mots (2).

\textbf{Proposition de rédaction modèle} :

\eng{Nowadays, information and communication technologies have changed the way young people spend their free time in my community. Instead of playing outside, many teenagers prefer their phones, the cybercafé or the gaming centre. This essay will discuss the advantages and disadvantages of using ICTs for recreation.}

\eng{On the one hand, ICTs offer cheap and pleasant entertainment. With a small data bundle, a student can stream music, watch films and chat with friends on social media. The gaming centre near the market is also a good place to kill time and relax after a long week at school. Moreover, funny videos that go viral can make the whole country laugh.}

\eng{On the other hand, excessive screen time can become a serious problem. Some teenagers are so hooked on video games that they forget their homework. Others are glued to the screen all day and no longer help their parents. In addition, when young people binge-watch series all night, they are tired in class the next morning.}

\eng{In my opinion, ICTs are wonderful tools for recreation, but young people must use them wisely and not let them take over their lives.}

(Environ 190 mots ; à allonger avec un exemple personnel pour atteindre 200 mots.)

\textbf{Points de vigilance} : \eng{spend time} et non \eng{pass time} ; \eng{actually} = « en fait », « actuellement » = \eng{nowadays} ; \eng{hooked on} et \eng{addicted to} avec la bonne préposition ; \eng{on social media} sans article ; présent de vérité générale tout au long de la discussion.
\end{corrige}

\newpage

\coursec{Fiche bilan}

\begin{popfigure}
\begin{tikzpicture}[mindmap, grow cyclic, every node/.style={concept}, concept color=popblue!70, root concept/.append style={concept color=popblue, font=\bfseries, minimum size=2.6cm}, level 1 concept/.append style={sibling angle=60, level distance=3.6cm, minimum size=2.2cm, font=\small}, level 2 concept/.append style={sibling angle=45, level distance=2.4cm, minimum size=1.7cm, font=\scriptsize}]
\node[root concept, text=white]{ICTs for recreation}
  child[concept color=poporange!80]{ node {Equipment and places}
    child[concept color=poporange!50]{ node {device, console, cybercafé} }
    child[concept color=poporange!50]{ node {headphones, router} }
  }
  child[concept color=popgreen!80]{ node {Online activities}
    child[concept color=popgreen!50]{ node {stream, download, browse} }
    child[concept color=popgreen!50]{ node {chat, play online} }
  }
  child[concept color=poppurple!80]{ node {Offline activities}
    child[concept color=poppurple!50]{ node {watch DVDs, edit photos} }
    child[concept color=poppurple!50]{ node {burn CDs} }
  }
  child[concept color=poppink!80]{ node {Collocations}
    child[concept color=poppink!50]{ node {surf the Net} }
    child[concept color=poppink!50]{ node {log on / log off} }
  }
  child[concept color=popgold!90]{ node {False friends}
    child[concept color=popgold!60]{ node {actually $\neq$ actuellement} }
    child[concept color=popgold!60]{ node {a library $\neq$ librairie} }
  }
  child[concept color=popblue!50]{ node {Word families}
    child[concept color=popblue!30]{ node {entertain, entertainment} }
    child[concept color=popblue!30]{ node {connect, connection} }
  };
\end{tikzpicture}
\legende{Carte mentale de la leçon : les TIC comme ressources de loisir.}
\end{popfigure}

\begin{bilan}
\textbf{1. Lexique clé à connaître par cœur}

\begin{center}
\begin{tabularx}{\linewidth}{|X|X|}\hline
\rowcolor{popblueL} \textbf{English} & \textbf{Français} \\ \hline
ICT (information and communication technology) & TIC (technologies de l'information et de la communication) \\ \hline
a device (smartphone, tablet, laptop) & un appareil (smartphone, tablette, ordinateur portable) \\ \hline
a games console & une console de jeux \\ \hline
headphones / earphones & un casque / des écouteurs \\ \hline
a cybercafé / an internet café & un cybercafé \\ \hline
free Wi-Fi / a hotspot & le Wi-Fi gratuit / un point d'accès \\ \hline
to surf the Net / to browse the web & surfer sur Internet \\ \hline
to stream music / films & écouter / regarder en streaming \\ \hline
to download / to upload & télécharger (vers son appareil / vers le réseau) \\ \hline
to chat online / instant messaging & bavarder en ligne / messagerie instantanée \\ \hline
social networks / social media & les réseaux sociaux \\ \hline
to play video games / online games & jouer aux jeux vidéo / en ligne \\ \hline
to log on / to log off & se connecter / se déconnecter \\ \hline
a screen / screen time & un écran / le temps passé devant un écran \\ \hline
entertainment / leisure & le divertissement / les loisirs \\ \hline
\end{tabularx}
\end{center}

\textbf{2. Collocations et expressions idiomatiques}
\begin{itemize}
\item \eng{to surf the Net}, \eng{to hang out online} « traîner en ligne », \eng{to be glued to one's screen} « être rivé à son écran ».
\item \eng{to kill time online} « passer le temps sur Internet », \eng{to binge-watch a series} « regarder une série en rafale ».
\item \eng{to go viral} « devenir viral », \eng{to keep in touch with friends} « rester en contact avec ses amis ».
\item Prépositions : \eng{to chat \textbf{with} friends}, \eng{to listen \textbf{to} music}, \eng{to play \textbf{with} a console}, \eng{to search \textbf{for} information}, \eng{to be addicted \textbf{to} video games}.
\end{itemize}

\textbf{3. Faux amis et pièges du francophone}

\begin{center}
\begin{tabularx}{\linewidth}{|X|X|X|}\hline
\rowcolor{popblueL} \textbf{Mot anglais} & \textbf{Sens réel} & \textbf{Piège français} \\ \hline
actually & en fait, réellement & $\neq$ « actuellement » = \eng{currently} \\ \hline
a library & une bibliothèque & $\neq$ « librairie » = \eng{a bookshop} \\ \hline
to assist & aider & $\neq$ « assister à » = \eng{to attend} \\ \hline
sensible & raisonnable & $\neq$ « sensible » = \eng{sensitive} \\ \hline
an opportunity & une occasion, une chance & attention à « opportunité » (anglicisme) \\ \hline
to demand & exiger & $\neq$ « demander » = \eng{to ask (for)} \\ \hline
\end{tabularx}
\end{center}

\textbf{4. Familles de mots, prononciation, orthographe}
\begin{itemize}
\item \eng{entertain} (v.) $\rightarrow$ \eng{entertainment} (n.) $\rightarrow$ \eng{entertaining} (adj.) ; \eng{connect} $\rightarrow$ \eng{connection} ; \eng{relax} $\rightarrow$ \eng{relaxation} ; \eng{communicate} $\rightarrow$ \eng{communication}.
\item Prononciation : \eng{technology} « tek-no-lo-dji », \eng{leisure} « lé-jeur », \eng{device} « di-vaïce », \eng{Wi-Fi} « ouaï-faï ».
\item Orthographe : \eng{leisure} (pas \emph{liesure}), \eng{download} en un seul mot, \eng{online} en un seul mot.
\end{itemize}

\textbf{5. Méthode express}
\begin{itemize}
\item À l'oral : présenter ses loisirs numériques avec \eng{I usually...}, \eng{In my free time, I enjoy...}, \eng{I am fond of...}.
\item À l'écrit : varier les verbes (stream, browse, share, post) et employer les prépositions correctes.
\item Toujours vérifier les faux amis avant de traduire mot à mot.
\end{itemize}
\end{bilan}

\begin{aretenir}[Checklist avant le Bac]
\begin{itemize}
\item[$\square$] Je connais le vocabulaire des équipements (device, console, headphones, router, hotspot).
\item[$\square$] Je sais nommer les lieux de loisir numérique (cybercafé, internet café, free Wi-Fi zone).
\item[$\square$] J'emploie correctement \eng{to surf the Net}, \eng{to stream}, \eng{to download}, \eng{to upload}.
\item[$\square$] Je maîtrise les prépositions : \eng{chat with}, \eng{listen to}, \eng{search for}, \eng{addicted to}.
\item[$\square$] Je distingue \eng{log on} et \eng{log off}, \eng{download} et \eng{upload}.
\item[$\square$] Je ne confonds plus \eng{actually} et « actuellement », \eng{library} et « librairie ».
\item[$\square$] Je connais les familles de mots : entertain / entertainment / entertaining ; connect / connection.
\item[$\square$] Je prononce correctement \eng{technology}, \eng{leisure}, \eng{device}, \eng{Wi-Fi}.
\item[$\square$] Je sais utiliser deux expressions idiomatiques : \eng{to be glued to one's screen}, \eng{to binge-watch}.
\item[$\square$] Je peux parler une minute de mes loisirs numériques à la maison ou au cybercafé.
\item[$\square$] Je peux écrire un court paragraphe sur les avantages et les dangers des TIC récréatives.
\item[$\square$] Je relis toujours mes productions pour traquer les calques du français.
\end{itemize}
\end{aretenir}

\findecours

\end{document}
